\documentclass[11pt,twoside,openany]{book}
\newcommand{\FgaLang}{spanish}
\usepackage[a4paper,inner=31mm,outer=27mm,top=28mm,bottom=31mm,marginparwidth=17mm]{geometry}
\usepackage{fontspec}
\defaultfontfeatures{Ligatures=TeX,Scale=MatchLowercase}
\setmainfont{TeX Gyre Pagella}
\setsansfont{TeX Gyre Heros}
\setmonofont{TeX Gyre Cursor}
\usepackage{polyglossia}
\setdefaultlanguage{\FgaLang}
\usepackage{microtype}
\usepackage{graphicx}
\usepackage{mathtools}
\usepackage{amsthm}
\usepackage{unicode-math}
\setmathfont{Latin Modern Math}
\usepackage{bookmark}
\usepackage{hyperref}
\hypersetup{hidelinks,unicode=true}

\setcounter{secnumdepth}{3}
\setcounter{tocdepth}{2}
\numberwithin{equation}{section}

\newcommand{\src}[1]{%
  \phantomsection\hypertarget{src:#1}{}%
  \marginpar{\raggedleft\footnotesize\sffamily #1}%
}
\newcommand{\sid}[1]{\phantomsection\hypertarget{fga:#1}{}}
\newcommand{\Osh}{\underline{O}}
\newcommand{\Zsh}{\underline{\mathbf Z}}
\newcommand{\Ccat}{\mathcal C}
\newcommand{\Homsh}{\underline{\operatorname{Hom}}}
\newcommand{\Extsh}{\underline{\operatorname{Ext}}}
\newcommand{\hsh}{\underline{h}}
\newcommand{\fquote}[1]{\textquotedbl#1\textquotedbl}
\newcommand{\seminar}[1]{%
  \begin{flushleft}\sffamily SEMINARIO BOURBAKI\\#1\end{flushleft}%
}
\newcommand{\seminaren}[1]{%
  \begin{flushleft}\sffamily BOURBAKI SEMINAR\\#1\end{flushleft}%
}
\newcommand{\paperhead}[2]{%
  \begin{center}
    {\Large\bfseries #1\par}
    \vspace{0.6\baselineskip}
    #2
  \end{center}
  \vspace{0.8\baselineskip}
}

\theoremstyle{plain}
\newtheorem{proposition}{Proposición}[section]
\newtheorem{theoreme}[proposition]{Teorema}
\newtheorem{theorem}[proposition]{Theorem}
\newtheorem{corollaire}[proposition]{Corolario}
\newtheorem{corollary}[proposition]{Corollary}
\newtheorem{lemme}[proposition]{Lema}
\newtheorem{lemma}[proposition]{Lemma}

\theoremstyle{definition}
\newtheorem{definition}[proposition]{Definición}
\newtheorem{definitionen}[proposition]{Definition}

\theoremstyle{remark}
\newtheorem*{remarque}{Observación}
\newtheorem*{remark}{Remark}

\newcommand{\unit}[1]{%
  \chapter{#1}%
  \markboth{#1}{#1}%
}
\newcommand{\fsection}[2]{%
  \section*{#1. #2}%
  \addcontentsline{toc}{section}{#1. #2}%
}
\newcommand{\result}[2]{%
  \par\medskip\noindent\textbf{#1%
  \if\relax\detokenize{#2}\relax\else\space#2\fi. ---}\enspace
}

% Printed page numbers recur across different works and errata.
% Namespace only the invisible PDF destinations; preserve all body bytes.
\newcommand{\FgaUnitKey}{frontmatter}
\renewcommand{\src}[1]{%
  \phantomsection\hypertarget{src:\FgaUnitKey:#1}{}%
  \marginpar{\raggedleft\footnotesize\sffamily #1}%
}
\setlength{\emergencystretch}{3em}
% The only inherited boldsymbol use is the Greek mu of the group scheme.
% Select its Unicode bold-italic glyph instead of a nonexistent bold math font.
\renewcommand{\boldsymbol}[1]{\symbfit{#1}}
\newcommand{\seminares}[1]{\begin{flushleft}\sffamily SEMINARIO BOURBAKI\\#1\end{flushleft}}
\theoremstyle{plain}
\newtheorem{teorema}[proposition]{Teorema}
\newtheorem{corolario}[proposition]{Corolario}
\newtheorem{lema}[proposition]{Lema}
\theoremstyle{definition}
\newtheorem{definicion}[proposition]{Definición}
\theoremstyle{remark}
\newtheorem*{observacion}{Observación}
\hypersetup{pdftitle={Fundamentos de la geometría algebraica -- Edición española actual -- tramo A},pdfauthor={Alexandre Grothendieck},pdfsubject={Exposiciones 149, 182, 190 y 195, con sus erratas}}
\begin{document}
% BEGIN RIGHTS AND AI TRANSLATION NOTICE
\begingroup
\small
\noindent\textbf{Fuentes y autoría}

La obra original, Fondements de la géométrie algébrique, es de Alexandre Grothendieck. La traducción humana al inglés de \href{https://github.com/thosgood/fga/tree/9976dd5f8e2ddcb388f0a07f3aa334f120d1019f}{Tim Hosgood} constituye una fuente de referencia de esta edición. Su proyecto reconoce las correcciones de Steve Hnizdur y el apoyo técnico de Jon Sterling. La versión 9976dd5f8e2ddcb388f0a07f3aa334f120d1019f identifica el texto inglés utilizado para la comparación de procedencia; no establece la versión histórica exacta incorporada inicialmente. La traducción española y las correcciones y composición posteriores fueron producidas por IA. La autoría de la traducción inglesa permanece con su traductor; no se le atribuye la traducción española. Esta atribución no constituye permiso ni aval.

\noindent\textbf{Derechos, disponibilidad y autoría}

Los derechos sobre las obras originales y sobre las contribuciones protegidas a estas traducciones y transcripciones derivadas corresponden a sus respectivos titulares. Este proyecto no dedica esas obras o contribuciones al dominio público ni concede una licencia sobre ellas. Las menciones de derechos de terceros y las autorizaciones ya concedidas siguen siendo aplicables.

Estas ediciones se ofrecen gratuitamente, con fines no comerciales y exclusivamente educativos, bajo la presunción de que los respectivos titulares de derechos no se opondrían a que estas traducciones y transcripciones estuvieran disponibles gratuitamente junto a los originales franceses en NUMDAM (www.numdam.org/). Esta presunción no significa que se haya obtenido autorización; el acceso gratuito a los originales no autoriza por sí solo su traducción, adaptación o republicación.

Si usted es titular de derechos sobre algún material presentado aquí y se opone a su disponibilidad, puede abrir una incidencia en GitHub (github.com/KokunoYumeto/ega-fga-es/issues). El responsable de la publicación retirará el material afectado a petición del titular de los derechos.

La traducción al español está producida por IA. Las correcciones, la edición y la composición también están producidas por IA. El responsable de la publicación no tradujo el texto. Se conservan la autoría y las atribuciones de las obras originales y de las contribuciones heredadas; no se atribuyen a la IA. Los modelos documentados se identifican en cada edición.
\par
\endgroup
\clearpage
% END RIGHTS AND AI TRANSLATION NOTICE

\frontmatter
\begin{titlepage}
\centering
\vspace*{0.13\textheight}
{\Huge\bfseries Fundamentos de la\par geometría algebraica\par}
\vspace{1.5\baselineskip}
{\Large Alexandre Grothendieck\par}
\vfill
{\Large Edición española actual -- tramo A\par}
\vspace{\baselineskip}
{\large Exposiciones 149, 182, 190 y 195, con sus erratas\par}
\vspace{1.5\baselineskip}
{\normalsize Edición española\par}
\end{titlepage}
\chapter*{Nota sobre esta edición}
\addcontentsline{toc}{chapter}{Nota sobre esta edición}
Este volumen contiene la traducción española de las exposiciones 149, 182, 190 y 195 y sus erratas.

La edición actual y la traducción según el impreso se ofrecen como textos distintos.

Las cifras marginales remiten a las páginas de las exposiciones originales.

Traducción, correcciones y composición producidas por IA: OpenAI Codex, GPT-6 Astra, esfuerzo Ultra.
La autoría de la obra original corresponde a Alexandre Grothendieck.

El PDF se acompaña de fuentes LaTeX editables y de los archivos necesarios para su compilación.

\tableofcontents
\mainmatter
\def\FgaUnitKey{149}\unit{Exposición 149 — Teoremas de dualidad para los haces algebraicos coherentes}
\sid{149}
\src{169}
\seminar{(Mayo de 1957)}
\paperhead{TEOREMAS DE DUALIDAD PARA LOS HACES ALGEBRAICOS COHERENTES}%
  {por Alexandre GROTHENDIECK}

Los resultados que siguen, inspirados en el \fquote{teorema de dualidad algebraica}
de Serre, fueron obtenidos en los inviernos de 1955 y 1956. Se demuestran muy
sencillamente con ayuda de resultados bastante elementales sobre la cohomología de los espacios
proyectivos [3] y mediante el uso intensivo del álgebra homológica de
Cartan--Eilenberg, en la forma de [2].

\fsection{1}{Los Ext de haces de módulos ({[2]}, caps. 3 y 4)}

Sea $X$ un espacio topológico provisto de un haz $\Osh$ de anillos con unidad
(no necesariamente conmutativos). Consideramos la categoría abeliana
$\Ccat^{\Osh}$ de los haces de $\Osh$-módulos, llamados también $\Osh$-Módulos.
Se sabe que todo objeto de esta categoría admite una resolución inyectiva, lo que
permite definir en ella los funtores $\operatorname{Ext}$ con las propiedades
formales bien conocidas. Más precisamente, para evitar confusiones,
denotamos por $\operatorname{Hom}_{\Osh}(X;A,B)$, o simplemente por
$\operatorname{Hom}(X;A,B)$, el grupo abeliano de los $\Osh$-homomorfismos de $A$
en $B$, mientras que $\Homsh_{\Osh}(A,B)$ denotará el haz de gérmenes
de homomorfismos de $A$ en $B$ ($A,B\in\Ccat^{\Osh}$). Para
$A\in\Ccat^{\Osh}$ fijo se definen funtores $h_A$ y $\hsh_A$ con valores, respectivamente,
en la categoría $\Ccat$ de los grupos abelianos y en la categoría
$\Ccat^{\Zsh}$ de los haces abelianos sobre $X$, mediante las fórmulas:
\begin{equation}\tag{1.1}
  h_A(B)=\operatorname{Hom}_{\Osh}(X;A,B),
  \qquad
  \hsh_A(B)=\Homsh_{\Osh}(A,B).
\end{equation}

Los funtores $h_A$ y $\hsh_A$ son funtores covariantes exactos por la izquierda,
de los cuales consideramos los funtores derivados derechos, denotados respectivamente por
$\operatorname{Ext}_{\Osh}^{p}(X;A,B)$ y $\Extsh_{\Osh}^{p}(A,B)$. Por tanto,
por definición,
\begin{equation}\tag{1.2}
\left\{
\begin{aligned}
\operatorname{Ext}_{\Osh}^{p}(X;A,B)
  &= (R^{p}h_A)(B)
   = H^{p}\bigl(\operatorname{Hom}_{\Osh}(X;A,C(B))\bigr),\\
\Extsh_{\Osh}^{p}(A,B)
  &= (R^{p}\hsh_A)(B)
   = H^{p}\bigl(\Homsh_{\Osh}(A,C(B))\bigr),
\end{aligned}
\right.
\end{equation}

donde el símbolo $R^{p}$ designa el paso a los funtores derivados derechos, y donde
$C(B)$ designa una resolución inyectiva arbitraria de $B$ en $\Ccat^{\Osh}$.
Denotemos por
\[
  \Gamma:\Ccat^{\Zsh}\longrightarrow\Ccat
\]
el funtor \fquote{secciones}. Recordemos que sus funtores derivados derechos se
denotan por
\[
  B\longmapsto H^{p}(X,B).
\]

\src{170}
\begin{equation}\tag{1.3}
  H^{p}(X,B)=(R^{p}\Gamma)(B)=H^{p}\bigl(\Gamma(C(B))\bigr).
\end{equation}
Evidentemente, se tiene $h_A=\Gamma\hsh_A$; por otra parte, se comprueba que $\hsh_A$
transforma objetos inyectivos en objetos acíclicos para $\Gamma$. Se concluye del
modo habitual:

\result{PROPOSICIÓN}{1} Para todo $\Osh$-Módulo $A$ existe un funtor
espectral cohomológico sobre $\Ccat^{\Osh}$, que converge al funtor graduado
$\bigl(\operatorname{Ext}_{\Osh}^{*}(X;A,B)\bigr)$ y cuyo término inicial es
\begin{equation}\tag{1.4}
  E_{2}^{p,q}(A,B)=H^{p}\bigl(X,\Extsh_{\Osh}^{q}(A,B)\bigr).
\end{equation}
De ella se deducen \fquote{homomorfismos de borde} y una sucesión exacta de cinco términos
que no escribiremos.

\result{COROLARIO}{1} Si $A$ es localmente isomorfo a $\Osh^{n}$, entonces se tienen
isomorfismos canónicos
\begin{equation}\tag{1.5}
  \operatorname{Ext}_{\Osh}^{p}(X;A,B)
  \overset{\sim}{\longleftarrow}
  H^{p}\bigl(X,\Homsh_{\Osh}(A,B)\bigr)
\end{equation}
(dados por homomorfismos de borde de la sucesión espectral). En particular, se tiene
un isomorfismo canónico
\begin{equation}\tag{1.6}
  \operatorname{Ext}_{\Osh}^{p}(X;\Osh,B)=H^{p}(X,B).
\end{equation}

Para utilizar estos resultados, es preciso saber explicitar los
$\Extsh_{\Osh}^{p}(A,B)$. Son funtores que se calculan localmente,
es decir, si $U$ es un abierto de $X$, se tiene
\[
  \Extsh_{\Osh}^{p}(A,B)|U
  =\Extsh_{\Osh|U}^{p}(A|U,B|U),
\]
como resulta del hecho de que la restricción a $U$ de un $\Osh$-Módulo inyectivo
es un $(\Osh|U)$-Módulo inyectivo. Además, para $x\in X$ fijo, se tienen homomorfismos funtoriales
\begin{equation}\tag{1.7}
  \Homsh_{\Osh}(A,B)_x
  \longrightarrow
  \operatorname{Hom}_{\Osh_x}(A_x,B_x),
\end{equation}
que se prolongan de manera única a un homomorfismo de funtores cohomológicos
en $B$:
\begin{equation}\tag{1.8}
  \Extsh_{\Osh}^{p}(A,B)_x
  \longrightarrow
  \operatorname{Ext}_{\Osh_x}^{p}(A_x,B_x).
\end{equation}

\result{PROPOSICIÓN}{2} Si $A$ es, en un entorno de $x$, isomorfo al
conúcleo de un homomorfismo $\Osh^{m}\longrightarrow\Osh^{n}$, entonces (1.8) es un isomorfismo en grado $p$ cuando $A$ admite localmente una resolución libre parcial hasta el grado $p+1$; es un isomorfismo para todo $p$ cuando $A$ es pseudocoherente. Una sola presentación finita da únicamente la comparación en grado cero.\footnote{El texto impreso dice « (1.7) »; la
presencia de $p$ y la frase anterior exigen « (1.8) ».} Esto ocurre, en
particular, si $A$ es un $\Osh$-Módulo coherente [3].

\src{171}
\result{PROPOSICIÓN}{3} Sea $L_{*}=(L_i)$ una resolución por la izquierda del
$\Osh$-Módulo $A$ por $\Osh$-Módulos cada uno de los cuales es localmente isomorfo a
un $\Osh^{n}$. Entonces $\Extsh_{\Osh}^{*}(A,B)$ se identifica con
$H^{*}\bigl(\Homsh_{\Osh}(L_{*},B)\bigr)$, y
$\operatorname{Ext}_{\Osh}^{*}(X;A,B)$ se identifica con la hipercohomología de $X$
con respecto al complejo $\Homsh_{\Osh}(L_{*},B)$.

La demostración es estándar: se considera el bicomplejo
$\Homsh_{\Osh}(L_{*},C(B))$, donde $C(B)$ es una resolución inyectiva de $B$, y los
homomorfismos naturales de $\Homsh_{\Osh}(L_{*},B)$ y de
$\Homsh_{\Osh}(A,C(B))$ en este último.

Para terminar, señalemos que los dos funtores $\operatorname{Ext}$ introducidos
en (1.2) no solo son funtores cohomológicos en $B$, sino también
bifuntores cohomológicos, covariantes en $B$ y contravariantes en $A$.

\fsection{2}{Ley de composición en los Ext}

Los resultados de este número se deben, independientemente, a CARTIER y a YONEDA; véase
una exposición de CARTIER [1] para más detalles. Sea $\Ccat$ una categoría abeliana.
Sean $K$ y $L$ dos objetos graduados de $\Ccat$. Denotamos por
$\operatorname{Hom}(K,L)$ el grupo abeliano graduado cuya componente de grado
$n$ está formada por los homomorfismos homogéneos de grado $n$ de $K$ en $L$ (es decir,
los sistemas $(u_i)$ de homomorfismos $K^{i}\longrightarrow L^{i+n}$).

Si $K$ y $L$ son complejos (con operadores diferenciales de grado $+1$, para
fijar ideas), se introduce en $\operatorname{Hom}(K,L)$ el operador
diferencial dado por
\begin{equation}\tag{2.1}
  \delta(u)=du+(-1)^{n+1}ud,\qquad\text{donde }n=\deg(u),
\end{equation}
que lo convierte en un complejo con operador diferencial de grado $+1$. Los ciclos de
grado $n$ son los operadores de grado $n$ que anticonmutan con $d$ (como
operadores homogéneos).\footnote{El texto impreso dice « con $u$ »; la ecuación
(2.1) exige « con $d$ ».} Se puede entonces considerar
$H^{*}\bigl(\operatorname{Hom}(K,L)\bigr)$; este es un invariante de los tipos
de homotopía de $K$ y de $L$, que podremos denotar por $H^{*}(K,L)$.

Si se tiene un tercer complejo $M$, la composición de los homomorfismos
define un acoplamiento
\[
  \operatorname{Hom}(K,L)\times\operatorname{Hom}(L,M)
  \longrightarrow\operatorname{Hom}(K,M)
\]
compatible con los operadores diferenciales, de donde, al pasar a la cohomología,
se obtienen acoplamientos
\begin{equation}\tag{2.2}
  H^{*}(K,L)\times H^{*}(L,M)\longrightarrow H^{*}(K,M),
\end{equation}
denotados por $(u,v)\longmapsto vu$. Estos acoplamientos satisfacen la evidente propiedad
de asociatividad. En particular, $H^{*}(K,K)$ es un anillo graduado asociativo
con unidad, y $H^{*}(K,L)$ (resp. $H^{*}(L,K)$) es un módulo graduado por la derecha
(resp. por la izquierda) sobre este anillo, etc. En grado $0$, (2.2) se reduce a la
composición de los morfismos de complejos. Finalmente, una sucesión exacta de
complejos
\[
  0\longrightarrow K'\longrightarrow K\longrightarrow K''\longrightarrow 0,
\]
\src{172}
tal que $K'^{i}$ se identifica con un sumando directo de $K^{i}$ para todo $i$,
da lugar a una sucesión exacta de complejos de grupos
$\operatorname{Hom}(K'',L)$, etc., y de ahí a un operador de conexión
\[
  H^{i}(K',L)\longrightarrow H^{i+1}(K'',L).
\]
Se definen de manera análoga los operadores de conexión relativos a una sucesión exacta en $L$. Los
acoplamientos (2.2) son compatibles, en el sentido usual, con estos operadores de conexión.

Supongamos ahora que $\Ccat$ es una categoría tal que todo objeto $A$
de $\Ccat$ admite una resolución inyectiva $C(A)$. Se comprueba entonces, utilizando
una de las numerosas variantes del teorema del bicomplejo, que
\[
  H^{*}\bigl(C(A),C(B)\bigr)
  =H^{*}\bigl(\operatorname{Hom}(C(A),C(B))\bigr)
\]
es canónicamente isomorfo a
\[
  H^{*}\bigl(\operatorname{Hom}(A,C(B))\bigr)
  =\operatorname{Ext}^{*}(A,B).
\]
Los operadores de conexión considerados anteriormente dan los homomorfismos de conexión de los
$\operatorname{Ext}$. Además, los acoplamientos (2.2) dan aquí
acoplamientos asociativos, compatibles con los homomorfismos de conexión:
\begin{equation}\tag{2.3}
  \operatorname{Ext}^{*}(A,B)\times\operatorname{Ext}^{*}(B,C)
  \longrightarrow\operatorname{Ext}^{*}(A,C).
\end{equation}
En particular, $\operatorname{Ext}^{*}(A,A)$ es un anillo graduado asociativo
con unidad, etc. (Se demuestra de manera análoga que los funtores
$\operatorname{Ext}$ operan en los funtores derivados de un funtor
cualquiera; no utilizaremos aquí este hecho.)

En el caso en que la categoría considerada sea la categoría $\Ccat^{\Osh}$ de los
$\Osh$-Módulos sobre $X$, se obtienen, por tanto, acoplamientos
\begin{equation}\tag{2.4}
  \operatorname{Ext}_{\Osh}^{p}(X;A,B)
  \times\operatorname{Ext}_{\Osh}^{q}(X;B,C)
  \longrightarrow\operatorname{Ext}_{\Osh}^{p+q}(X;A,C),
\end{equation}
que pueden explicitarse como se ha indicado. Por lo demás, los mismos
desarrollos, sustituyendo la categoría de los grupos abelianos por la
categoría de los haces abelianos sobre $X$, y los funtores $\operatorname{Hom}$
por los funtores $\Homsh$, definen también acoplamientos con las mismas
propiedades formales y, esta vez, de \fquote{naturaleza local}:
\begin{equation}\tag{2.5}
  \Extsh_{\Osh}^{p}(A,B)\times\Extsh_{\Osh}^{q}(B,C)
  \longrightarrow\Extsh_{\Osh}^{p+q}(A,C).
\end{equation}
Estos últimos se precisan observando que los homomorfismos (1.8) son
compatibles con los acoplamientos entre los $\operatorname{Ext}$.

\src{173}
Recordemos, por último, que existe también una estructura multiplicativa entre los funtores
$H^{p}(X,A)$ (producto cup). Se comprueba entonces que las sucesiones espectrales de la
proposición 1 son compatibles con las estructuras multiplicativas; más
precisamente, se tiene un acoplamiento de la sucesión espectral $E(A,B)$ y de la sucesión
espectral $E(B,C)$ hacia la sucesión espectral $E(A,C)$ que, en el límite,
se reduce al acoplamiento entre los $\operatorname{Ext}$ globales y, en el
término inicial, al acoplamiento deducido, en el segundo miembro de (1.4), del
producto cup y de los acoplamientos de los $\operatorname{Ext}$ locales. De ello resulta,
en particular, que los \fquote{homomorfismos de borde}
\begin{equation}\tag{2.6}
  \operatorname{Ext}_{\Osh}^{n}(X;A,B)
  \longrightarrow H^{0}\bigl(X,\Extsh_{\Osh}^{n}(A,B)\bigr)
\end{equation}
y
\begin{equation}\tag{2.7}
  H^{n}\bigl(X,\Homsh_{\Osh}(A,B)\bigr)
  \longrightarrow\operatorname{Ext}_{\Osh}^{n}(X;A,B)
\end{equation}
son compatibles con las estructuras multiplicativas. Por tanto, si nos restringimos a los
haces localmente isomorfos a un $\Osh^{m}$, esto explicita por completo la
composición entre los $\operatorname{Ext}$ globales mediante el producto cup, teniendo
en cuenta el isomorfismo (1.5).

\fsection{3}{Resultados de cohomología local}

Sea $A$ un anillo conmutativo con unidad provisto de un ideal $J$. Vamos a
definir, para un $A$-módulo variable $M$, homomorfismos funtoriales
\begin{equation}\tag{3.1}
\left\{
\begin{aligned}
\operatorname{Ext}_{A}^{p}(A/J,M)
  &\longrightarrow
  \operatorname{Hom}_{A}\!\left(\bigwedge^{p}J/J^{2},M\otimes A/J\right),\\
\operatorname{Tor}_{p}^{A}(A/J,M)
  &\longleftarrow
  \left(\bigwedge^{p}J/J^{2}\right)
  \otimes\operatorname{Hom}_{A}(A/J,M),
\end{aligned}
\right.
\end{equation}
donde los productos tensoriales y exteriores se toman sobre el anillo $A$; nótese,
por otra parte, que $J/J^{2}$ es en realidad un $A/J$-módulo, y que sus potencias
exteriores como $A$-módulo o como $A/J$-módulo son las mismas.
Definir los homomorfismos (3.1) equivale a definir, para todo
sistema $\underline{x}=(x_1,\ldots,x_p)$ de elementos de $J$, homomorfismos
$\varphi_{\underline{x}}$:
\begin{equation}\tag{3.2}
\left\{
\begin{aligned}
\varphi_{\underline{x}}:
  \operatorname{Ext}_{A}^{p}(A/J,M)&\longrightarrow M\otimes A/J,\\
\varphi_{\underline{x}}:
  \operatorname{Hom}_{A}(A/J,M)&\longrightarrow
  \operatorname{Tor}_{p}^{A}(A/J,M),
\end{aligned}
\right.
\end{equation}
de tal manera que se satisfagan las condiciones siguientes:

\src{174}
\begin{enumerate}
\item[\textup{i.}] $\varphi_{x_1,\ldots,x_p}$ depende de manera
$A$-multilineal alternada del sistema de los $x_i\in J$.
\item[\textup{ii.}] $\varphi_{x_1,\ldots,x_p}$ es nulo cuando alguno de los $x_i$ está
en $J^{2}$.
\end{enumerate}
En realidad, ii. resulta de i., pues $a\varphi_{\underline{x}}=0$ para $a\in J$,
como se ve observando que todos los módulos de (3.2) son anulados por
$J$.

Para definir $\varphi_{\underline{x}}$, consideremos el complejo
$K_{\underline{x}}$ cuyo $A$-módulo subyacente es $\bigwedge A^{p}$, y
cuyo operador diferencial es el producto interior $i_{\underline{x}}$ por
$\underline{x}$ considerado como forma lineal sobre $A^{p}$ de componentes
$x_1,\ldots,x_p$. El operador diferencial es de grado $-1$, los grados son
positivos, y la homología de este complejo en grado $0$ es
$A/(x_1A+\cdots+x_pA)$. Como los $x_i$ están en $J$, se deduce una
aumentación $K_{\underline{x},0}\longrightarrow A/J$. Así,
$K_{\underline{x}}$ aparece como un complejo libre aumentado, con módulo
de aumentación $A/J$. De ello se deducen, del modo habitual, homomorfismos
\[
\operatorname{Ext}_{A}^{*}\bigl(H_0(K_{\underline{x}}),M\bigr)
\longrightarrow
H^{*}\bigl(\operatorname{Hom}_{A}(K_{\underline{x}},M)\bigr)
\]
y
\[
\operatorname{Tor}_{*}^{A}\bigl(H_0(K_{\underline{x}}),M\bigr)
\longleftarrow
H_{*}\bigl(K_{\underline{x}}\otimes_A M\bigr),
\]
de donde, componiendo con los homomorfismos sobre los $\operatorname{Ext}$ y los
$\operatorname{Tor}$ deducidos del homomorfismo de aumentación
$H_0(K_{\underline{x}})\longrightarrow A/J$, se obtienen homomorfismos
\begin{equation}\tag{3.3}
\left\{
\begin{aligned}
\psi_{\underline{x}}:
\operatorname{Ext}_{A}^{*}(A/J,M)
&\longrightarrow H^{*}\bigl(\operatorname{Hom}_{A}(K_{\underline{x}},M)\bigr),\\
\psi_{\underline{x}}:
\operatorname{Tor}_{*}^{A}(A/J,M)
&\longleftarrow H_{*}\bigl(K_{\underline{x}}\otimes_A M\bigr).
\end{aligned}
\right.
\end{equation}
Ahora bien, se comprueba inmediatamente que, en grado máximo $p$, la cohomología de los segundos
miembros es $M/(x_1M+\cdots+x_pM)$ (resp. es el conjunto de los elementos de $M$
anulados por cada uno de los $x_i$). Como los $x_i$ están en $J$, se deducen
homomorfismos
\begin{equation}\tag{3.4}
\left\{
\begin{aligned}
H^{p}\bigl(\operatorname{Hom}_{A}(K_{\underline{x}},M)\bigr)
&\longrightarrow M\otimes A/J,\\
H_{p}\bigl(K_{\underline{x}}\otimes_A M\bigr)
&\longleftarrow\operatorname{Hom}_{A}(A/J,M).
\end{aligned}
\right.
\end{equation}
Componiendo los homomorfismos (3.3) y (3.4), se obtienen los homomorfismos (3.2)
que queríamos definir. La comprobación de i. es laboriosa, pero no
presenta dificultades.

\src{175}
\result{PROPOSICIÓN}{4} Sea $A$ un anillo conmutativo con unidad, y sea
$(x_1,\ldots,x_p)$ una sucesión de elementos de $A$ tal que, para $1\leq i\leq p$,
la imagen de $x_i$ en el cociente de $A$ por el ideal generado por
$(x_1,\ldots,x_{i-1})$ no sea divisor de $0$. Sea $J$ el ideal generado
por los $x_i$. Entonces $J/J^{2}$ es un $(A/J)$-módulo libre que tiene por base las
imágenes canónicas de los $x_i$, el complejo $K_{\underline{x}}$ es una resolución
libre de $A/J$, y para todo $A$-módulo $M$, los homomorfismos (3.1) en
grado $p$ son biyectivos. Lo mismo ocurre con los homomorfismos análogos
definidos para grados $i$ cualesquiera, siempre que $J.M=0$.

(El punto esencial, del cual resultan todos los demás, es la aciclicidad de
$K_{\underline{x}}$, que es un hecho bien conocido bajo las condiciones indicadas.)

\result{COROLARIO}{1} Sean $A$ y $J$ como arriba, y supongamos además que
$A$ es un álgebra afín regular de $\dim n$ sobre un cuerpo perfecto $k$, y
que $A/J$ es un álgebra afín regular. Denotemos por $\Omega^{1}(A)$,
$\Omega^{1}(A/J)$ los módulos de diferenciales de Kähler. Entonces se tiene un
isomorfismo canónico compatible con la localización,
\begin{equation}\tag{3.5}
  \operatorname{Ext}_{A}^{p}\bigl(\Omega^{n-p}(A/J),\Omega^{n}(A)\bigr)=A/J.
\end{equation}
En efecto, $\Omega^{n-p}(A/J)$ es un $(A/J)$-módulo proyectivo de rango $1$, y asimismo
$\Omega^{n}(A)$ es un $A$-módulo proyectivo de rango $1$,\footnote{El texto impreso
dice « rango $n$ »; los dos módulos de formas de grado máximo tienen rango
$1$, como exige también la identificación con $\bigwedge^{p}(J/J^{2})$ más abajo.}
por lo que el primer miembro es igual a
\[
  \operatorname{Ext}_{A}^{p}(A/J,A)
  \otimes\Omega^{n-p}(A/J)'\otimes\Omega^{n}(A)
\]
(donde el símbolo « $'$ » designa el $(A/J)$-módulo dual). El producto tensorial de los
dos últimos factores se identifica con $\bigwedge^{p}(J/J^{2})$, por lo que el conjunto
se identifica con
\[
  \operatorname{Ext}_{A}^{p}\!\left(A/J,\bigwedge^{p}(J/J^{2})\right),
\]
y, por la proposición, con
\[
  \operatorname{Hom}_{A}\!\left(\bigwedge^{p}J/J^{2},
  \bigwedge^{p}J/J^{2}\right),
\]
es decir, con $A/J$. En particular, en
$\operatorname{Ext}_{A}^{p}\bigl(\Omega^{n-p}(A/J),\Omega^{n}(A)\bigr)$ hay un
elemento distinguido, correspondiente al elemento unidad de $A/J$, llamado
\emph{clase fundamental del ideal $J$ en $A$}. (En realidad, puede
definirse bajo condiciones sensiblemente más generales.) El
corolario 1 puede darse en una forma más geométrica y más global:

\result{COROLARIO}{2} Sean $X$ una variedad no singular sobre un cuerpo
algebraicamente cerrado $k$, $Y$ una subvariedad cerrada no singular de $X$,
$\Osh_X$ el haz de anillos locales de $X$, y $\Osh_Y$ el haz de anillos
locales de $Y$ considerado como haz cociente de $\Osh_X$. Sean $n$ la
dimensión de $X$ y $n-p$ la de $Y$.

\src{176}
Sea $\Omega_X$ (resp. $\Omega_Y$) el haz de gérmenes de formas
diferenciales regulares sobre $X$ (resp. $Y$). Entonces se tienen isomorfismos
canónicos:
\begin{equation}\tag{3.6}
  \Extsh_{\Osh_X}^{p}\bigl(\Omega_Y^{n-p},\Omega_X^{n}\bigr)=\Osh_Y,
\end{equation}
o también
\begin{equation}\tag{3.6 bis}
  \Extsh_{\Osh_X}^{p}\bigl(\Osh_Y,\Omega_X^{n}\bigr)=\Omega_Y^{n-p}.
\end{equation}
La fórmula (3.6 bis) puede servir como definición de $\Omega_Y^{n-p}$ cuando $Y$
es una variedad singular. Más precisamente:

\result{PROPOSICIÓN}{5} Sea $Y$ un subconjunto algebraico de dimensión
$q=n-p$ en una variedad algebraica no singular $X$ de dimensión $n$. Sean
$F$ un haz algebraico coherente sobre $X$ cuyo soporte está contenido en $Y$, y $L$
un haz algebraico localmente libre sobre $X$. Entonces los haces
$\Extsh_{\Osh_X}^{i}(F,L)$ son nulos para $i<p$, mientras que para $i=p$ se tiene un
isomorfismo canónico
\begin{equation}\tag{3.7}
  \Extsh_{\Osh_X}^{p}(F,L)
  =\Homsh_{\Osh_X}\!\left(F,
  \Extsh_{\Osh_X}^{p}(\Osh_X/J,L)\right),
\end{equation}
donde $J$ designa un haz de ideales arbitrario sobre $X$ que anula $F$ y cuyo
conjunto de ceros es $Y$. En particular, si $F$ es un haz algebraico
coherente sobre $Y$, se tiene
\begin{equation}\tag{3.7 bis}
  \Extsh_{\Osh_X}^{p}(F,L)
  =\Homsh_{\Osh_Y}\!\left(F,\Extsh_{\Osh_X}^{p}(\Osh_Y,L)\right).
\end{equation}
Por último, siendo todavía $F$ un haz algebraico coherente sobre $Y$, los haces
\[
  E^{i}(F)=\Extsh_{\Osh_X}^{p+i}(F,\Omega_X^{n})
\]
no dependen de la manera de sumergir el espacio algebraico $Y$ en una variedad
no singular $X$.

Como la cuestión es local, se puede suponer que $X$ es afín y $L=\Osh_X$. Esto reduce
el problema a una cuestión de álgebra conmutativa, e incluso de álgebra local: si $A$ es un
anillo local regular y $M$ un $A$-módulo cuyo soporte tiene dimensión
$\leq q=n-p$, hay que demostrar que $\operatorname{Ext}_{A}^{i}(M,A)=0$ para $i<p$, y
\[
  \operatorname{Ext}_{A}^{p}(M,A)
  =\operatorname{Hom}_{A}\!\left(M,\operatorname{Ext}_{A}^{p}(A/J,A)\right),
\]
donde $J$ es un ideal cualquiera de « dimensión » $\leq q$ que anula $M$. Para el
primer punto se procede por inducción sobre $q$: un devissage inmediato reduce
al caso en que $M$ es de la forma $A/J$ y luego, sustituyendo $J$ por un ideal más
pequeño y utilizando la hipótesis de inducción y la sucesión exacta de los
$\operatorname{Ext}$, al caso en que $J$ está generado por un « sistema de parámetros »
como en la proposición 4, donde el resultado es inmediato. El resultado
anterior implica que, si $J$ es un ideal

\src{177}
fijo de « dimensión » $\leq q$, entonces el funtor contravariante
$E(M)=\operatorname{Ext}_{A}^{p}(M,A)$ sobre la categoría de los $A/J$-módulos es
exacto por la izquierda; además, transforma sumas directas en productos directos,
de donde resulta fácilmente que
\[
  E(M)=\operatorname{Hom}_{A}\bigl(M,E(A/J)\bigr).
\]
\footnote{El texto impreso dice $E(A)$; en la categoría considerada, el objeto
representante es $A/J$, y la igualdad requerida para (3.7) es $E(A/J)$.} Por último, la
última afirmación de la proposición 5 es más sutil y resulta de una
caracterización intrínseca de los $E^{i}(F)$ mediante un teorema de dualidad
local que no puede exponerse aquí.

\result{COROLARIO}{} Denotemos por $\omega_Y^{q}$ el haz
$\Extsh_{\Osh_X}^{p}(\Osh_Y,\Omega_X^{n})$. Entonces se tiene un isomorfismo
funtorial para los haces algebraicos coherentes $F$ sobre $Y$:
\begin{equation}\tag{3.8}
  \Extsh_{\Osh_X}^{p}(F,\Omega_X^{n})
  =\Homsh_{\Osh_X}(F,\omega_Y^{q}).
\end{equation}

\fsection{4}{Clase de cohomología asociada a una subvariedad}

En todo lo que sigue, $X$ designa un conjunto algebraico de dimensión $n$, definido
sobre un cuerpo $k$, que supondremos algebraicamente cerrado para simplificar. Salvo
en el n\textsuperscript{o} 6, se supone que $X$ es no singular. Denotamos por
$\Osh_X$ el haz de anillos locales sobre $X$, y por
\[
  \Omega_X^{*}=\bigoplus_p\Omega_X^{p}
\]
el haz de gérmenes de formas diferenciales sobre $X$.\footnote{El texto
impreso emplea un signo de unión; la estructura graduada es la suma directa
de los haces $\Omega_X^p$.} Si $Y$ es una parte cerrada de $X$, identificamos
los haces algebraicos coherentes sobre $Y$ con haces coherentes sobre $X$
nulos fuera de $Y$; en particular, así ocurre con $\Osh_Y$ y $\Omega_Y$.

\result{LEMA}{1} Sean $F$ un haz algebraico coherente sobre $X$ cuyo
soporte tiene dimensión $\leq n-p$, y $L$ un haz algebraico coherente
localmente libre sobre $X$. Entonces $\operatorname{Ext}_{\Osh_X}^{i}(X;F,L)$ es nulo
para $i<p$, y se tiene un isomorfismo canónico
\begin{equation}\tag{4.1}
  \operatorname{Ext}_{\Osh_X}^{p}(X;F,L)
  =H^{0}\bigl(X,\Extsh_{\Osh_X}^{p}(F,L)\bigr).
\end{equation}
Si $F$ es un haz algebraico coherente sobre una parte cerrada $Y$ de $X$ de
dimensión $\leq n-p$, se tiene un isomorfismo canónico
\begin{equation}\tag{4.1 bis}
  \Extsh_{\Osh_X}^{p}(F,L)
  =\Homsh_{\Osh_X}\!\left(
    F\otimes L'\otimes\Omega_X^{n},\omega_Y^{n-p}
  \right),
\end{equation}
donde $\omega_Y^{n-p}$ es el haz sobre $Y$ definido en el corolario de la
proposición 5 (que se identifica con $\Omega_Y^{n-p}$ si $Y$ es no singular).

La fórmula (4.1) es una consecuencia inmediata de la sucesión espectral de la
proposición 1 y de la proposición 5; en virtud de la fórmula (3.8), se puede
escribir

\src{178}
\begin{align*}
\Extsh_{\Osh_X}^{p}(F,L)
&=L\otimes(\Omega_X^{n})'\otimes
  \Extsh_{\Osh_X}^{p}(F,\Omega_X^{n})\\
&=L\otimes(\Omega_X^{n})'\otimes
  \Homsh_{\Osh_X}(F,\omega_Y^{q})\\
&=\Homsh_{\Osh_X}\!\left(
  F\otimes L'\otimes\Omega_X^{n},\omega_Y^{q}\right),
\end{align*}
donde $q=n-p$, de donde resulta la fórmula (4.1 bis).

Tomando, en particular, $F=\Osh_Y$, $L=\Omega_X^{p}$, se obtiene, teniendo en cuenta
que $\Omega_X^{n}\otimes(\Omega_X^{p})'=\Omega_X^{n-p}$, un isomorfismo
canónico
\begin{equation}\tag{4.2}
  \operatorname{Ext}_{\Osh_X}^{p}(X;\Osh_Y,\Omega_X^{p})
  =\operatorname{Hom}_{\Osh_X}
  (X;\Omega_X^{n-p},\omega_Y^{n-p}).
\end{equation}
Supongamos ahora que $Y$ es no singular, para simplificar, de modo que
$\omega_Y^{n-p}=\Omega_Y^{n-p}$. Se tiene un homomorfismo natural de
$\Omega_X^{n-p}$ en $\Omega_Y^{n-p}$, de donde resulta una sección canónica $s_Y$ del
haz $\Extsh_{\Osh_X}^{p}(\Osh_Y,\Omega_X^{p})$, que llamaremos, si
todas las componentes de $Y$ tienen dimensión $n-p$, \emph{sección
fundamental} del haz $\Extsh_{\Osh_X}^{p}(\Osh_Y,\Omega_X^{p})$. Esta
sección define, en virtud de (4.1), un elemento de
$\operatorname{Ext}_{\Osh_X}^{p}(X;\Osh_Y,\Omega_X^{p})$. Ahora bien, el homomorfismo
natural $\Osh_X\longrightarrow\Osh_Y$ define un homomorfismo
\[
  \operatorname{Ext}_{\Osh_X}^{p}(X;\Osh_Y,\Omega_X^{p})
  \longrightarrow
  \operatorname{Ext}_{\Osh_X}^{p}(X;\Osh_X,\Omega_X^{p})
  =H^{p}(X,\Omega_X^{p}).
\]
Se obtiene así un elemento de $H^{p}(X,\Omega_X^{p})$, denotado por $P_X(Y)$ y llamado
\emph{clase de cohomología de $Y$ en $X$}. Se deduce, por tanto, de la sección
$s_Y$ de $\Extsh_{\Osh_X}^{p}(\Osh_Y,\Omega_X^{p})$ mediante el diagrama
de homomorfismos siguiente:
\begin{equation}\tag{4.3}
\begin{aligned}
H^{p}(X,\Omega_X^{p})
&=\operatorname{Ext}_{\Osh_X}^{p}(X;\Osh_X,\Omega_X^{p})\\
&\longleftarrow\operatorname{Ext}_{\Osh_X}^{p}(X;\Osh_Y,\Omega_X^{p})
\xrightarrow{\ \sim\ }
H^{0}\!\left(X,\Extsh_{\Osh_X}^{p}(\Osh_Y,\Omega_X^{p})\right).
\end{aligned}
\end{equation}

Llamemos \emph{ciclo no singular de dimensión $n-p$} a un elemento del grupo
abeliano libre generado por las subvariedades irreducibles no singulares de
dimensión $n-p$ en $X$. Entonces la función $Y\longmapsto P(Y)$ se prolonga a
un homomorfismo del grupo de ciclos no singulares de dimensión $n-p$ sobre $X$
en el grupo $H^{p}(X,\Omega_X^{p})$.

Sean $Z^{n-p}$ y $Z'^{n-p'}$ ciclos no singulares de dimensiones $n-p$ y
$n-p'$; se dice que se \emph{intersecan transversalmente} si toda componente
de $Z$ interseca transversalmente toda componente de $Z'$. Entonces el ciclo $Z.Z'$ está
definido y es un ciclo no singular de dimensión $n-p-p'$. Dicho esto, se tiene el

\src{179}
\result{TEOREMA}{1} Si $Z^{n-p}$ y $Z'^{n-p'}$ son ciclos no singulares
que se intersecan transversalmente, entonces se tiene
\begin{equation}\tag{4.4}
  P_X(Z.Z')=P_X(Z).P_X(Z'),
\end{equation}
donde el producto del segundo miembro es el producto cup
\[
  H^{p}(X,\Omega_X^{p})\times H^{p'}(X,\Omega_X^{p'})
  \longrightarrow H^{p+p'}(X,\Omega_X^{p+p'}).
\]
(Se supone que $X$ es isomorfo a una parte localmente cerrada de un espacio
proyectivo.)

Esta última hipótesis sirve únicamente para poder concluir que todo
haz algebraico coherente sobre $X$ es cociente de un haz algebraico
coherente localmente libre (SERRE) y, por tanto, admite una resolución por la izquierda mediante
haces localmente libres. Para demostrar el teorema 1, se puede suponer que
$Z$ y $Z'$ son subvariedades irreducibles no singulares $Y$ y $Y'$ que se
intersecan transversalmente. Sea $L_{*}$ una resolución por la izquierda de $\Osh_Y$ mediante
haces localmente libres; entonces, en virtud de la proposición 3, el diagrama
de homomorfismos (4.3) se identifica con el diagrama
\[
R^{p}\Gamma(\Omega_X^{p})
\xleftarrow{\alpha}
\mathbb R^{p}\Gamma\!\left(\Homsh_{\Osh_X}(L_{*},\Omega_X^{p})\right)
\xrightarrow[\sim]{\beta}
\Gamma\!\left(H^{p}\bigl(\Homsh_{\Osh_X}(L_{*},\Omega_X^{p})\bigr)\right),
\]
donde $\Gamma$ es el funtor « grupo de secciones » sobre la categoría de los
haces abelianos sobre $X$, y $\mathbb R^{p}\Gamma$ designa
la hipercohomología de grado $p$ de dicho funtor, y $R^{p}\Gamma$ su
$p$-ésimo funtor derivado. Supondremos, para simplificar, que $L_0=\Osh_X$ y que
la aumentación $L_0\longrightarrow\Osh_Y$ es el homomorfismo natural (lo cual
es lícito); entonces $\alpha$ se deduce del homomorfismo de complejos
$\Osh_X\longrightarrow L_{*}$ (considerando $\Osh_X$ como un complejo concentrado
en grado $0$), teniendo en cuenta que $\mathbb R^{p}\Gamma(K)=R^{p}\Gamma(K_0)$ si $K$
es un complejo de haces concentrado en grado $0$. El homomorfismo $\beta$ es un
« homomorfismo de borde » bien conocido. Consideremos un diagrama análogo, relativo a
una resolución localmente libre $L'_{*}$ de $\Osh_{Y'}$, y consideremos el
diagrama conmutativo de acoplamientos:
\begin{equation}\tag{4.5}
\resizebox{\textwidth}{!}{$
\begin{array}{ccc}
R^{p}\Gamma(\Omega_X^{p})
&\longleftarrow
\mathbb R^{p}\Gamma\!\left(\Homsh_{\Osh_X}(L_{*},\Omega_X^{p})\right)
&\xrightarrow{\sim}
\Gamma H^{p}\!\left(\Homsh_{\Osh_X}(L_{*},\Omega_X^{p})\right)
\\[3pt]
\times&\times&\times\\[3pt]
R^{p'}\Gamma(\Omega_X^{p'})
&\longleftarrow
\mathbb R^{p'}\Gamma\!\left(\Homsh_{\Osh_X}(L'_{*},\Omega_X^{p'})\right)
&\xrightarrow{\sim}
\Gamma H^{p'}\!\left(\Homsh_{\Osh_X}(L'_{*},\Omega_X^{p'})\right)
\\[3pt]
\downarrow&\downarrow&\downarrow\\[3pt]
R^{p+p'}\Gamma(\Omega_X^{p+p'})
&\longleftarrow
\mathbb R^{p+p'}\Gamma\!\left(
\Homsh_{\Osh_X}(L_{*}\otimes L'_{*},\Omega_X^{p+p'})\right)
&\xrightarrow{\sim}
\Gamma H^{p+p'}\!\left(
\Homsh_{\Osh_X}(L_{*}\otimes L'_{*},\Omega_X^{p+p'})\right).
\end{array}
$}
\end{equation}

\src{180}
Los acoplamientos de las dos columnas de la derecha se deducen del acoplamiento de
complejos de haces
\[
\Homsh_{\Osh_X}(L_{*},\Omega_X^{p})
\times\Homsh_{\Osh_X}(L'_{*},\Omega_X^{p'})
\longrightarrow
\Homsh_{\Osh_X}(L_{*}\otimes L'_{*},\Omega_X^{p+p'}),
\]
que se define mediante el producto exterior
$\Omega_X^{p}\times\Omega_X^{p'}\longrightarrow\Omega_X^{p+p'}$.
El acoplamiento de la primera columna es el producto cup (relativo al producto
exterior). Afirmo que la última fila de (4.5) se identifica con el diagrama
de isomorfismos análogo a (4.3), donde $Y$ se sustituye por $Y\cap Y'$ y $p$ por
$p+p'$. Para ello basta demostrar que $L\otimes L'$ es una resolución
(evidentemente localmente libre) de $\Osh_{Y\cap Y'}$. En efecto, se tiene
\[
  H_0(L\otimes L')=\Osh_Y\otimes\Osh_{Y'}=\Osh_{Y\cap Y'}
\]
y
\[
  H_i(L\otimes L')=\operatorname{Tor}_{i}^{\Osh_X}(\Osh_Y,\Osh_{Y'})=0
  \qquad(i>0),
\]
por el hecho de que $Y$ y $Y'$ se intersecan transversalmente. El teorema 1 resulta
ahora de la fórmula:
\begin{equation}\tag{4.6}
  s_Y.s_{Y'}=s_{Y.Y'}
\end{equation}
(donde el producto del primer miembro es el de la última columna de (4.5)).
Esta fórmula (4.6), de naturaleza puramente local, se comprueba sin dificultad
tomando para $L_{*}$ y $L'_{*}$ las resoluciones consideradas en la proposición
4. Se demuestra asimismo (mediante una demostración más sencilla) que
$Z\longmapsto P_X(Z)$ es compatible con el producto cartesiano:
\begin{equation}\tag{4.7}
  P_{X\times X'}(Z\times Z')=P_X(Z)\otimes P_{X'}(Z')
\end{equation}
(fórmula válida si $Z$, resp. $Z'$, es un ciclo no singular sobre la variedad
no singular $X$, resp. $X'$, considerando $Z\times Z'$ como un ciclo no
singular sobre $X\times X'$). De (4.4) y (4.7) se deduce que $P_X(Z)$ es también
compatible con la operación « imagen inversa » por un morfismo de variedades no
singulares $f:X\longrightarrow X'$:
\begin{equation}\tag{4.8}
  P_X\bigl(f^{-1}(Z')\bigr)=f^{*}\bigl(P_{X'}(Z')\bigr),
\end{equation}
fórmula válida si $Z'$ es un ciclo no singular sobre $X'$ tal que $f$ sea
« transversal » a $Z'$, es decir, tal que la gráfica de $f$ sea transversal al ciclo
$X\times Z'$ en $X\times X'$.

\src{181}
\result{COROLARIO}{1} Sean $X$, $X'$ dos variedades no singulares
localmente cerradas en un espacio proyectivo, y supongamos que $X'$ es completa. Sean $U$
un ciclo no singular sobre $X\times X'$, $a$ y $b$ dos puntos de $X'$ tales que
$U$ interseca transversalmente los ciclos $X\times(a)$ y $X\times(b)$, y sean $Z$
y $Z'$ los ciclos no singulares sobre $X$ tales que
\[
  Z\times(a)=(X\times(a)).U,
  \qquad
  Z'\times(b)=(X\times(b)).U.
\]
Entonces se tiene
\[
  P_X(Z)=P_X(Z').
\]

En efecto, sea $f_a:X\longrightarrow X\times X'$ definido por
$f_a(x)=(x,a)$; entonces, en virtud de (4.8), se tiene la fórmula
$P(Z)=f_a^{*}(P(U))$, y análogamente $P(Z')=f_b^{*}(P(U))$. Ahora bien, utilizando la fórmula de
Künneth
\[
  H^{*}(X\times X',\Omega_{X\times X'}^{*})
  =H^{*}(X,\Omega_X^{*})\otimes H^{*}(X',\Omega_{X'}^{*})
\]
y el hecho de que $H^{0}(X',\Osh_{X'})$ se reduce a los escalares, se obtiene
fácilmente $f_a^{*}=f_b^{*}$, de donde se sigue la conclusión.

Para todo $x\in X$, $(x)$ es una subvariedad no singular de $X$ de
codimensión $n$, y define por tanto un elemento $\varepsilon_x$ de
$H^{n}(X,\Omega_X^{n})$. Si $X$ es una variedad proyectiva no singular, del
corolario 1 resulta que $\varepsilon_x$ no depende del punto $x$ elegido; se
denota por $\varepsilon_X$ y se llama \emph{clase fundamental} de
$H^{n}(X,\Omega_X^{n})$.

\result{OBSERVACIÓN}{} Para disponer de una teoría satisfactoria, habría que definir
$P_X(Z)$ para ciclos $Z$ cualesquiera y demostrar el teorema 1 para una
intersección propia de ciclos. (En el momento de escribir esta exposición, esto aún no se
ha hecho con toda generalidad.) Admitiendo que se haya hecho, el corolario 1
pasa a ser: si $Z$ y $Z'$ son dos ciclos algebraicamente equivalentes, entonces
$P_X(Z)=P_X(Z')$ (afirmación que no parece resultar de lo anterior, incluso si
$Z$ y $Z'$ son no singulares).

\fsection{5}{El teorema de dualidad}

En este número, $X$ designa una variedad proyectiva no singular de dimensión
$n$.

\result{TEOREMA}{2} La clase fundamental $\varepsilon_X$ de
$H^{n}(X,\Omega_X^{n})$ es una base de este espacio vectorial.

(La demostración se dará más adelante.) En virtud del teorema anterior, se puede
identificar $H^{n}(X,\Omega_X^{n})$ con el cuerpo $k$. Consideremos ahora los
acoplamientos contemplados en el n\textsuperscript{o} 2, que dan en particular un
acoplamiento

\src{182}
\[
\operatorname{Ext}_{\Osh_X}^{p}(X;\Osh_X,F)
\times\operatorname{Ext}_{\Osh_X}^{n-p}(X;F,\Omega_X^{n})
\longrightarrow
\operatorname{Ext}_{\Osh_X}^{n}(X;\Osh_X,\Omega_X^{n}),
\]
es decir,
\begin{equation}\tag{5.1}
  H^{p}(X,F)
  \times\operatorname{Ext}_{\Osh_X}^{n-p}(X;F,\Omega_X^{n})
  \longrightarrow H^{n}(X,\Omega_X^{n}).
\end{equation}
Teniendo en cuenta el teorema 2, este acoplamiento define un homomorfismo
\begin{equation}\tag{5.2}
  \operatorname{Ext}_{\Osh_X}^{n-p}(X;F,\Omega_X^{n})
  \longrightarrow\bigl(H^{p}(X,F)\bigr)'.
\end{equation}
Este homomorfismo es funtorial en $F$ y conmuta con los homomorfismos de conexión
relativos a las sucesiones exactas $0\longrightarrow F'\longrightarrow F
\longrightarrow F''\longrightarrow0$.

\result{TEOREMA}{3} El homomorfismo (5.2) es un isomorfismo.

Se recupera, en particular, el resultado de Serre:

\result{COROLARIO}{} Sean $E$ un fibrado vectorial algebraico sobre $X$ y
$\Osh_X(E)$ el haz de gérmenes de secciones regulares de $E$; entonces se tienen
isomorfismos canónicos:
\begin{equation}\tag{5.3}
  \bigl(H^{p}(X,\Osh_X(E))\bigr)'
  =H^{n-p}\bigl(X,\Omega_X^{n}\otimes\Osh_X(E')\bigr).
\end{equation}
Basta aplicar el teorema 3 y el corolario 1 a la proposición 1.

Los teoremas 2 y 3 resultarán del enunciado siguiente:

\noindent\textbf{(D) El homomorfismo siguiente}
\begin{equation}\tag{5.2 bis}
  \operatorname{Ext}_{\Osh_X}^{n-p}(X;F,\Omega_X^{n})
  \longrightarrow\bigl(H^{p}(X,F)\bigr)'\otimes L,
  \qquad L=H^{n}(X,\Omega_X^{n}),
\end{equation}
deducido del acoplamiento (5.1) es un isomorfismo.

Demostremos, en efecto, que (D) implica el teorema 2. Sea
$k_x=\Osh_{(x)}$ el haz de anillos locales de la variedad reducida al punto
$x\in X$; consideremos el homomorfismo canónico
$\Osh_X\longrightarrow k_x$ y el homomorfismo asociado
\begin{equation}\tag{5.3 bis}
  H^{0}(X,\Osh_X)\longrightarrow H^{0}(X,k_x),
\end{equation}
su transpuesto se identifica con el homomorfismo
\begin{equation}\tag{5.4}
  \operatorname{Ext}_{\Osh_X}^{n}(X;k_x,\Omega_X^{n})\otimes L'
  \longrightarrow
  \operatorname{Ext}_{\Osh_X}^{n}(X;\Osh_X,\Omega_X^{n})\otimes L'
\end{equation}
deducido del homomorfismo sobre los $\operatorname{Ext}^{n}$ asociado a
$\Osh_X\longrightarrow k_x$:
\begin{equation}\tag{5.5}
  \operatorname{Ext}_{\Osh_X}^{n}(X;k_x,\Omega_X^{n})
  \longrightarrow
  \operatorname{Ext}_{\Osh_X}^{n}(X;\Osh_X,\Omega_X^{n})
  =H^{n}(X,\Omega_X^{n}).
\end{equation}

\src{183}
Como (5.3 bis) es un isomorfismo,\footnote{El texto impreso da una
segunda fórmula « (5.3) » en la página 182 y remite aquí a ella; el identificador
« (5.3 bis) » distingue esta fórmula de la primera (5.3).} lo mismo ocurre
con (5.4), y por tanto también con (5.5). Como $s_{(x)}$ es una base de
$\operatorname{Ext}_{\Osh_X}^{n}(X;k_x,\Omega_X^{n})$ en virtud de (4.2), se
sigue que su imagen $\varepsilon_X$ es una base de
$H^{n}(X,\Omega_X^{n})$.

Queda por demostrar el enunciado (D), que resultará de manera puramente formal de los hechos
elementales resumidos en los lemas siguientes. En ellos se supone que $X$ es una
parte cerrada (singular o no) del espacio proyectivo $P$ de dimensión $r$.
Se utiliza la notación $\Osh_P(m)$ para el haz sobre $P$ denotado por $\Osh(m)$ en
[3], y la notación $\Osh_X(m)$ para el haz análogo sobre $X$.

\result{LEMA}{2} El enunciado (D) es verdadero si $X=P$ y $F=\Osh_P(m)$.

Este lema se comprueba mediante un cálculo directo. El cálculo explícito de los
$H^{i}(P,\Osh_P(m))$ se encuentra en [3], pero puede hacerse de manera más
elemental. El cálculo del producto cup
\[
  H^{i}(P,\Osh_P(m))\times H^{j}(P,\Osh_P(m'))
  \longrightarrow H^{i+j}(P,\Osh_P(m+m'))
\]
necesario para explicitar el acoplamiento (5.1) no presenta
dificultades.\footnote{El texto impreso dice $\Osh_P(r+r')$; los dos factores
$\Osh_P(m)$ y $\Osh_P(m')$ exigen $\Osh_P(m+m')$.}

\result{LEMA}{3} Todo haz algebraico coherente $F$ sobre $X$ es isomorfo a
un haz cociente de un haz $\Osh_X(-m)^{k}$, donde se puede suponer que $m$ es
tan grande como se quiera.

Resulta del hecho de que $F\otimes\Osh_X(m)$ está « generado por sus secciones » para
$m$ grande, cf. [3].

\result{LEMA}{4} Sea $i>0$. Entonces
$H^{r-i}(P,\Osh_P(-m))=0$ para $m$ grande; y para todo haz algebraico
coherente $B$ sobre $X$, se tiene
\[
  \operatorname{Ext}_{\Osh_X}^{i}(X;\Osh_X(-m),B)=0
\]
para $m$ grande.

El primer enunciado resulta de los cálculos explícitos mencionados anteriormente; para el
segundo, se observa que se tiene un isomorfismo
\[
  \operatorname{Ext}_{\Osh_X}^{i}(X;\Osh_X(-m),B)
  =H^{i}(X,B\otimes\Osh_X(m))
\]
(proposición 1, corolario 1), de donde se sigue la conclusión en virtud de un resultado bien
conocido de [3]. Combinando los dos lemas anteriores, se obtiene el

\result{COROLARIO}{} Sea $i>0$. Entonces el funtor
$F\longmapsto H^{r-i}(P,F)$ sobre la categoría de los haces algebraicos coherentes
sobre $P$ es borrable, y lo mismo ocurre con el funtor
$\operatorname{Ext}_{\Osh_X}^{i}(X;F,B)$ sobre la categoría de los haces
algebraicos coherentes sobre $X$.

\src{184}
\result{LEMA}{5} Sean $A$ y $B$ dos haces algebraicos coherentes sobre
$X$, y pongamos $A(m)=A\otimes\Osh_X(m)$. Entonces, para $m$ suficientemente grande,
el homomorfismo canónico
\[
  \operatorname{Ext}_{\Osh_X}^{1}(X;A(-m),B)
  \longrightarrow H^{0}\bigl(X,\Extsh_{\Osh_X}^{1}(A(-m),B)\bigr)
  =H^{0}\bigl(X,\Extsh_{\Osh_X}^{1}(A,B)(m)\bigr)
\]
es un isomorfismo.

Esto resulta inmediatamente de la sucesión espectral de la proposición 1 aplicada a
$A(-m)$ y $B$, pues se tendrá
\[
  E_{2}^{p,q}(A(-m),B)
  =H^{p}\bigl(X,\Extsh_{\Osh_X}^{q}(A(-m),B)\bigr)
  =H^{p}\bigl(X,\Extsh_{\Osh_X}^{q}(A,B)(m)\bigr),
\]
que es nulo para $p>0$ y $m$ grande.

Demostremos ahora (D) en el caso en que $X=P$. Probamos primero que (5.2 bis)
es un isomorfismo para $p=n$; como los dos miembros son entonces
funtores exactos por la izquierda (pues $H^{r+1}(P,F)=0$), del lema 3 resulta
que basta demostrar la afirmación para $F=\Osh_P(-m)$, pero esta se encuentra
contenida en el lema 2. Como los homomorfismos (5.2 bis) son funtoriales y
compatibles con los operadores de conexión, y como para $p<n$ los dos miembros de
(5.2 bis) son funtores borrables en $F$ (corolario del lema 4), se
sigue entonces, mediante un razonamiento estándar, que (5.2 bis) es un isomorfismo
para todo $p$. Esto demuestra el teorema de dualidad para el espacio proyectivo.
Supongamos ahora que $X$ es cualquiera, pero no singular. Por el teorema de
dualidad para $P$, se tiene un isomorfismo
\[
  H^{n}(X,F)'=H^{n}(P,F)'
  =\operatorname{Ext}_{\Osh_P}^{r-n}(P;F,\Omega_P^{r}).
\]
En virtud del lema 1 (n\textsuperscript{o} 4), el último miembro se identifica con
\[
  \operatorname{Hom}_{\Osh_P}(P;F,\omega_X^{n})
  =\operatorname{Hom}_{\Osh_X}(X;F,\Omega_X^{n})
  =\operatorname{Ext}_{\Osh_X}^{0}(X;F,\Omega_X^{n}),
\]
de donde se obtiene un isomorfismo
\begin{equation}\tag{5.6}
  H^{n}(X,F)'=\operatorname{Hom}_{\Osh_X}(X;F,\Omega_X^{n})
  =\operatorname{Ext}_{\Osh_X}^{0}(X;F,\Omega_X^{n}).
\end{equation}
Tomando $F=\Omega_X^{n}$, se obtiene un isomorfismo
\begin{equation}\tag{5.7}
  \gamma:H^{n}(X,\Omega_X^{n})\xrightarrow{\sim}k.
\end{equation}

\src{185}
Se comprueba que el isomorfismo (5.6) no es otro que (5.2 bis) para $p=n$, cuando
se toma $L=k$ gracias a (5.7). Por consiguiente, (5.2 bis) es un isomorfismo para
$p=n$. Para demostrar que es un isomorfismo para todo $p$, basta nuevamente
demostrar que, para $p<n$, los dos miembros de (5.2 bis) son funtores
borrables en $F$, y a fortiori (teniendo en cuenta el lema 3) que ambos miembros
son nulos cuando se toma $F=\Osh_X(-m)$ con $m$ grande. Ahora bien, para el primer
miembro, esto es cierto en virtud del lema 4, y para el segundo miembro se escribe,
utilizando el teorema de dualidad sobre $P$:
\[
  H^{p}(X,\Osh_X(-m))'
  =\operatorname{Ext}_{\Osh_P}^{r-p}
    (P;\Osh_X(-m),\Omega_P^{r}).
\]
El segundo miembro es nulo para $p<n$ y $m$ grande, como resulta del lema 5
(donde se toma $X=P$) y del hecho de que $\Osh_X$, como haz algebraico
coherente sobre $P$, tiene dimensión cohomológica $\leq r-n$ (pues $X$ es no
singular), de donde
\[
  \Extsh_{\Osh_P}^{r-p}(\Osh_X,\Omega_P^{r})=0
  \qquad\text{para }p<n.
\]

\fsection{6}{El teorema de dualidad para las variedades singulares}

Sea $X$ una parte cerrada de dimensión $n$ del espacio proyectivo $P$ de
dimensión $r$. La fórmula (5.6) debe escribirse aquí
\begin{equation}\tag{6.1}
  H^{n}(X,F)'
  \simeq\operatorname{Hom}_{\Osh_X}(X;F,\omega_X^{n})
  =\operatorname{Ext}_{\Osh_X}^{0}(X;F,\omega_X^{n}),
\end{equation}
donde se ha puesto
\begin{equation}\tag{6.2}
  \omega_X^{n}=E^{0}(\Osh_X)
  =\Extsh_{\Osh_P}^{r-n}(\Osh_X,\Omega_P^{r}).
\end{equation}
Como se dijo en la proposición 5, el haz así introducido no depende
en realidad de la inmersión elegida de $X$ en una variedad no singular $P$.
Tomando, en (6.1), $F=\omega_X^{n}$, se obtiene
\begin{equation}\tag{6.2 bis}
  H^{n}(X,\omega_X^{n})'
  \simeq\operatorname{Hom}_{\Osh_X}(X;\omega_X^{n},\omega_X^{n}),
\end{equation}
\footnote{El texto impreso numera también esta fórmula « (6.2) »; el sufijo
« bis » distingue las dos apariciones.} de donde resulta la existencia de un elemento
distinguido en $H^{n}(X,\omega_X^{n})'$, correspondiente al homomorfismo
identidad de $\omega_X^{n}$ en sí mismo:
\begin{equation}\tag{6.3}
  \gamma:H^{n}(X,\omega_X^{n})\longrightarrow k.
\end{equation}

\src{186}
Consideremos entonces los acoplamientos definidos por la composición de los $\operatorname{Ext}$:
\begin{equation}\tag{6.4}
  H^{p}(X,F)\times
  \operatorname{Ext}_{\Osh_X}^{n-p}(X;F,\omega_X^{n})
  \longrightarrow H^{n}(X,\omega_X^{n}),
\end{equation}
y compongámoslos con el homomorfismo $\gamma$ de (6.3); se obtienen
homomorfismos funtoriales, compatibles con los operadores de borde:
\begin{equation}\tag{6.5}
  \operatorname{Ext}_{\Osh_X}^{n-p}(X;F,\omega_X^{n})
  \longrightarrow H^{p}(X,F)'
\end{equation}
(que generalizan (5.2)). Se comprueba que para $p=n$ se obtiene así
el isomorfismo (6.1). Sentado esto, se tiene el

\result{TEOREMA}{3 bis} Las cuatro condiciones siguientes sobre $X$ son
equivalentes para un entero dado $k\geq0$:
\begin{enumerate}
  \item[\textit{i.}] El homomorfismo funtorial (6.5) es un isomorfismo para
    $n-k\leq p\leq n$.
  \item[\textit{ii.}] Se tiene $H^{p}(X,\Osh_X(-m))=0$ para $m$ grande y para
    $n-k\leq p<n$.
  \item[\textit{iii.}] El funtor $H^{p}(X,F)$ sobre la categoría de los haces
    algebraicos coherentes sobre $X$ es borrable para $n-k\leq p<n$.
  \item[\textit{iv.}] Se tiene
    $E^{i}(\Osh_X)=\Extsh_{\Osh_P}^{r-n+i}(\Osh_X,\Omega_P^{r})=0$
    para $0<i\leq k$.
\end{enumerate}

\noindent\textbf{DEMOSTRACIÓN.} --- \textit{i.}$\Rightarrow$\textit{ii.} en
virtud del lema 4, \textit{ii.}$\Rightarrow$\textit{iii.} en virtud del lema 3,
\textit{iii.}$\Rightarrow$\textit{i.} mediante un razonamiento estándar bien conocido,
teniendo en cuenta que los dos miembros de (6.5) son entonces funtores borrables
para $n-k\leq p<n$ (el primero lo es en virtud del lema 4). Por último, demostremos
\textit{ii.}$\Leftrightarrow$\textit{iv.} Esto resulta del corolario de la
proposición siguiente:

\result{PROPOSICIÓN}{6} Sea $F$ un haz algebraico coherente sobre $X$, y
sea $i$ un entero. Entonces, para $m$ suficientemente grande, se tiene un isomorfismo:
\begin{equation}\tag{6.6}
  H^{i}(X,F(-m))'
  \simeq H^{0}\bigl(X,E^{n-i}(F)(m)\bigr),
\end{equation}
donde se pone (compárese n\textsuperscript{o} 3, proposición 5):
\begin{equation}\tag{6.7}
  E^{j}(F)=\Extsh_{\Osh_P}^{r-n+j}(F,\Omega_P^{r}).
\end{equation}
En efecto, en virtud del teorema de dualidad sobre $P$, el primer miembro de (6.6)
es isomorfo a
$\operatorname{Ext}_{\Osh_P}^{r-i}(P;F(-m),\Omega_P^{r})$, por lo que (6.6) resulta
del lema 5.

\src{187}
\result{COROLARIO}{} Para que se tenga $H^{i}(X,F(-m))=0$ para $m$ grande, es
necesario y suficiente que $E^{n-i}(F)=0$.

Recordemos que los $E^{j}(F)$ no dependen en realidad de la inmersión proyectiva
considerada. La condición del corolario es de naturaleza puramente local; por consiguiente,
si se verifica para $F$, se verifica para todo haz localmente
isomorfo a un haz $F^{n}$. En particular, si esta condición se verifica
para $\Osh_X$, se verifica para todo haz algebraico coherente localmente
libre. Este es, por ejemplo, el caso para todo $i<n$ si $X$ es no singular, para
$i=0$ si ninguna componente de $X$ se reduce a un solo punto, y para $i=0,1$
si $X$ es normal y todas sus componentes tienen dimensión $>1$ (cf.
[3]).\footnote{El texto impreso dice aquí « $S$ es normal »; el párrafo
solo se refiere a la variedad $X$.} Para que esto ocurra para todo
$i<k$, es necesario y suficiente, por definición, que los anillos locales
$\Osh_x$ ($x\in X$) tengan « codimensión homológica » $\geq k$ (cf. [4]
para esta noción). Si $k=n$, esto significa, en virtud del teorema 3 bis, que el
teorema de dualidad es verdadero para $X$, es decir, que (6.5) es un isomorfismo para
todo $p$ y todo $F$. Se pueden dar numerosas condiciones equivalentes sobre
los anillos locales $\Osh_x$ para que esto ocurra (Nagata), por ejemplo que
satisfagan el teorema de equidimensionalidad de Cohen--Macaulay. Así ocurre,
por ejemplo, si localmente $X$ es una « intersección completa » en una
variedad ambiente no singular.

\fsection{7}{La dualidad de Poincaré}

Sea $X$ una variedad proyectiva no singular de dimensión $n$. Entonces
\[
  H^{**}(X)=H^{*}(X,\Omega_X^{*})
\]
es un álgebra anticonmutativa bigraduada de dimensión finita, que graduamos
por el grado total, de modo que $H^{p,q}(X)=H^{p}(X,\Omega_X^{q})$ tiene grado
$p+q$. Los grados de $H^{*}(X)$ están comprendidos entre $0$ y $2n$. En virtud del
teorema 2 y del corolario del teorema 3, $H^{*}(X)$ es un « álgebra de
Poincaré » de dimensión $2n$, es decir, $H^{2n}(X)$ está provisto de un isomorfismo con
el cuerpo base $k$, y el producto
\[
  H^{m}(X)\times H^{2n-m}(X)\longrightarrow H^{2n}(X)=k
\]
es una dualidad entre $H^{m}(X)$ y $H^{2n-m}(X)$. Además, si $Y$ es otra
variedad proyectiva no singular, la fórmula de Künneth para los
haces algebraicos coherentes da
\begin{equation}\tag{7.1}
  H^{*}(X\times Y)=H^{*}(X)\otimes H^{*}(Y),
\end{equation}
isomorfismo compatible con las estructuras de álgebras de Poincaré. Además,
$H^{*}(X)$ es, como álgebra conmutativa, un funtor contravariante en
$X$, pues un morfismo $f:Y\longrightarrow X$ define de manera evidente un
homomorfismo de álgebras graduadas

\src{188}
\begin{equation}\tag{7.2}
  f^{*}:H^{*}(X)\longrightarrow H^{*}(Y).
\end{equation}
Como se trata de álgebras de Poincaré, se obtiene por transposición un
homomorfismo de espacios vectoriales:
\begin{equation}\tag{7.3}
  f_{*}:H^{*}(Y)\longrightarrow H^{*}(X).
\end{equation}
Se vio en el n\textsuperscript{o} 4 que el efecto de $f^{*}$ sobre las clases de
cohomología correspondientes a ciclos no singulares se interpreta
geométricamente tomando las clases de cohomología que corresponden a sus
imágenes inversas. En el caso actual importa demostrar que (7.3) corresponde
análogamente a la operación « imagen directa » de ciclos. Esto resulta, al menos bajo
condiciones adecuadas de no singularidad, del caso particular siguiente:

\result{TEOREMA}{4} Si $f$ es la inclusión de una subvariedad no
singular $Y^{m}$ de $X^{n}$ en $X^{n}$, entonces, denotando por $1_Y$
el elemento unidad de $H(Y)$, se tiene
\begin{equation}\tag{7.4}
  f_{*}(1_Y)=P_X(Y),
\end{equation}
donde el segundo miembro es la clase de cohomología en $X$ asociada a $Y$.

Esta fórmula equivale a
\begin{equation}\tag{7.4 bis}
  \bigl\langle\xi^{m,m},P_X(Y)\bigr\rangle\varepsilon_Y
  =f^{*}(\xi^{m,m})
  \qquad
  \bigl(\xi^{m,m}\in H^{m}(X,\Omega_X^{m})\bigr),
\end{equation}
\footnote{El texto impreso dice $f_{*}(\xi^{m,m})$. Puesto que
$\xi^{m,m}\in H^{m}(X,\Omega_X^{m})$ y que ambos miembros pertenecen a
$H^{m}(Y,\Omega_Y^{m})$, la aplicación requerida es $f^{*}$.} donde
$\varepsilon_Y$ es el elemento fundamental en $H^{m}(Y,\Omega_Y^{m})$, y
constituye, en el caso de las variedades proyectivas no singulares, una nueva
definición de la clase de cohomología asociada a $Y$. Para demostrar el
teorema 4, se interpreta, gracias al teorema 3, el transpuesto del homomorfismo
\[
  H^{m}(X,\Omega_X^{m})
  \longrightarrow H^{m}(Y,\Omega_Y^{m})
  =H^{m}(X,\Omega_Y^{m})
\]
como el homomorfismo
\begin{equation}\tag{7.5}
\begin{split}
  H^{n-m}(X,\Omega_X^{n-m})
  &\simeq
  \operatorname{Ext}_{\Osh_X}^{n-m}(X;\Omega_X^{m},\Omega_X^{n})\\
  &\longleftarrow
  \operatorname{Ext}_{\Osh_X}^{n-m}(X;\Omega_Y^{m},\Omega_X^{n})
  \simeq
  \operatorname{Hom}_{\Osh_X}(X;\Omega_Y^{m},\Omega_Y^{m}).
\end{split}
\end{equation}
Se comprueba que el elemento $1_Y$ del dual de $H^{m}(Y,\Omega_Y^{m})$ se identifica
con el elemento del miembro de la derecha correspondiente al endomorfismo identidad de
$\Omega_Y^{m}$, y por otra parte que la imagen de este elemento en
$H^{n-m}(X,\Omega_X^{n-m})$ es precisamente $P_X(Y)$. Estas compatibilidades, que podrían
haberse dado ya en el n\textsuperscript{o} 4, pueden enunciarse, y son
en realidad válidas, para variedades no singulares cualesquiera; la segunda

\src{189}
resulta, por ejemplo, de la conmutatividad del diagrama siguiente de endomorfismos
canónicos:
\begin{equation}\tag{7.6}
\resizebox{\textwidth}{!}{$
\left\{
\begin{array}{ccccc}
\operatorname{Ext}_{\Osh_X}^{n-m}(X;\Omega_X^{m},\Omega_X^{n})
&\longleftarrow&
\operatorname{Ext}_{\Osh_X}^{n-m}(X;\Omega_Y^{m},\Omega_X^{n})
&\simeq&
\operatorname{Hom}_{\Osh_X}(X;\Omega_Y^{m},\Omega_Y^{m})\\[3pt]
\Big\downarrow&&&&\Big\downarrow\\[3pt]
\operatorname{Ext}_{\Osh_X}^{n-m}(X;\Osh_X,\Omega_X^{n-m})
&\longleftarrow&
\operatorname{Ext}_{\Osh_X}^{n-m}(X;\Osh_Y,\Omega_X^{n-m})
&\simeq&
\operatorname{Hom}_{\Osh_X}(X;\Omega_X^{m},\Omega_Y^{m})
\end{array}
\right.
$}
\end{equation}

Se obtiene así un equivalente exacto del formalismo de la dualidad de Poincaré
sobre las variedades orientadas compactas. En particular, el teorema 4 permite
determinar la clase de cohomología asociada a la diagonal de $X\times X$. Mediante
un razonamiento bien conocido, se deduce, por ejemplo, una fórmula de Lefschetz:

\result{TEOREMA}{5} Sea $f$ un endomorfismo de una variedad proyectiva no
singular $X$, tal que los puntos fijos de $f$ tengan multiplicidad 1.
Entonces el número de estos puntos fijos es congruente, módulo la característica de
$k$, con la suma alternada de las trazas de los endomorfismos de los $H^{i}(X)$ definidos
por $f$.

La restricción sobre $f$ que ha sido necesario imponer se debe a las dificultades mencionadas en
la observación del n\textsuperscript{o} 4. No obstante, se observará que la fórmula de
Lefschetz sigue siendo verdadera si $f$ es « homotópico » a un endomorfismo cuyos
puntos fijos tienen todos multiplicidad 1.

\fsection{8}{Generalización del teorema de dualidad}

Sea $X$ una variedad algebraica no singular tal que todo haz
algebraico coherente $F$ sobre $X$ sea isomorfo a un cociente de un haz
algebraico coherente localmente libre (lo cual ocurre si $X$ es localmente
cerrado en un espacio proyectivo). Entonces todo haz algebraico coherente $F$ sobre
$X$ admite una resolución finita $L$ mediante haces localmente libres, y para
dos resoluciones de esta clase siempre se puede encontrar una tercera y un
homomorfismo de esta en las dos primeras compatible con las
aumentaciones. Análogamente, si $L$ es una resolución finita localmente libre de
$F$, y si se da un homomorfismo $F'\longrightarrow F$, entonces existe
una resolución finita localmente libre $L'$ de $F'$ y un homomorfismo
$L'\longrightarrow L$ compatible con $F'\longrightarrow F$, que se puede
suponer sobreyectivo si $F'\longrightarrow F$ lo es. Esto permite definir, dados
dos enteros $r,s\geq0$, dos multifuntores cohomológicos en los
argumentos $A_1,\ldots,A_r$; $B_1,\ldots,B_s$,

\src{190}
que varían en la categoría de los haces algebraicos coherentes sobre $X$, y cuyos
valores están en la categoría de los haces algebraicos coherentes sobre $X$,
resp. en la categoría de los módulos sobre $H^{0}(X,\Osh_X)$, mediante las fórmulas:
\begin{equation}\tag{8.1}
\resizebox{\textwidth}{!}{$
\left\{
\begin{aligned}
\underline{T}_{r}^{s*}(A_1,\ldots,A_r;B_1,\ldots,B_s)
  &={}H^{*}\!\left(
    \Homsh_{\Osh_X}\!\left(
      L(A_1)\otimes\cdots\otimes L(A_r),
      L(B_1)\otimes\cdots\otimes L(B_s)
    \right)\right),\\[3pt]
T_{r}^{s*}(A_1,\ldots,A_r;B_1,\ldots,B_s)
  &={}\mathbf R^{*}\Gamma\!\left(
    \Homsh_{\Osh_X}\!\left(
      L(A_1)\otimes\cdots\otimes L(A_r),
      L(B_1)\otimes\cdots\otimes L(B_s)
    \right)\right).
\end{aligned}
\right.
$}
\end{equation}
En esta fórmula, $L(F)$ designa una resolución finita localmente libre del
haz algebraico coherente $F$, y $\mathbf R^{*}\Gamma(K)$ designa
la hipercohomología del espacio $X$ con respecto al complejo de haces $K$. Si
$r$ o $s$ es nulo, se sustituye el producto tensorial de los $L(A_i)$, resp. de los
$L(B_j)$, por $\Osh_X$. En particular, $\underline{T}_{0}^{0}$ y $T_{0}^{0}$
son funtores graduados de 0 argumentos, $\underline{T}_{0}^{0}$ se reduce al
grado 0, donde es el haz $\Osh_X$, mientras que $T_{0}^{0}$ es idéntico a
$H^{*}(X,\Osh_X)$. El hecho de que los segundos miembros de (8.1) no dependan
de las resoluciones elegidas es, en cualquier caso, evidente para la primera línea
(pues la cuestión es entonces local), y para la segunda resulta de las observaciones
generales precedentes, teniendo en cuenta la sucesión espectral de
hipercohomología del complejo de haces
\[
  K=\Homsh_{\Osh_X}\bigl(L(A_1)\otimes\cdots,
      L(B_1)\otimes\cdots\bigr),
\]
que converge a la hipercohomología de $X$ con respecto a $K$ y cuyo término
inicial es $H^{p}(X,H^{q}(K))$, es decir,
\begin{equation}\tag{8.2}
  E_{2}^{p,q}
  =H^{p}\!\left(X,
    (\underline{T}_{r}^{s})^{(q)}
      (A_1,\ldots,A_r;B_1,\ldots,B_s)\right).
\end{equation}
Se ve entonces que esta sucesión espectral tampoco depende de las resoluciones
elegidas; su límite es el primer miembro de (8.1). Se definen
fácilmente los operadores de conexión relativos a los diversos argumentos $A_i$, $B_j$,
observando que toda sucesión exacta
$0\longrightarrow F'\longrightarrow F\longrightarrow F''\longrightarrow0$
puede resolverse mediante una sucesión exacta de complejos finitos localmente libres.

En el sistema de los funtores $\underline{T}_{r}^{s*}$, resp.
$T_{r}^{s*}$, se definen operaciones análogas a las del cálculo tensorial, cuya
definición es inmediata a partir de las fórmulas definitorias (8.1). Así, se tiene
una composición (que generaliza la considerada en el n\textsuperscript{o} 2):
\begin{equation}\tag{8.3}
\resizebox{\textwidth}{!}{$
\begin{aligned}
&T_{r}^{s*}(A_1,\ldots,A_r;B_1,\ldots,B_s)
\times T_{r'}^{s'*}(A'_1,\ldots,A'_{r'};B'_1,\ldots,B'_{s'})\\
&\qquad\longrightarrow
T_{r+r'}^{s+s'*}
(A_1,\ldots,A_r,A'_1,\ldots,A'_{r'};
 B_1,\ldots,B_s,B'_1,\ldots,B'_{s'}),
\end{aligned}
$}
\end{equation}
que satisface las propiedades evidentes de asociatividad y de compatibilidad con
los homomorfismos funtoriales, los homomorfismos de conexión y las sucesiones
espectrales. Asimismo, se tienen operaciones de simetría cuya explicitación se deja al
lector. Se tiene además una operación de contracción cada vez que
uno de los argumentos $A_i$

\src{191}
es igual a uno de los argumentos $B_j$:
\begin{equation}\tag{8.4}
\resizebox{\textwidth}{!}{$
\begin{aligned}
&T_{r}^{s*}(A_1,\ldots,A_{i-1},C,A_{i+1},\ldots,A_r;
 B_1,\ldots,B_{j-1},C,B_{j+1},\ldots,B_s)\\
&\qquad\longrightarrow
T_{r-1}^{s-1*}(A_1,\ldots,\widehat{A_i},\ldots,A_r;
 B_1,\ldots,\widehat{B_j},\ldots,B_s).
\end{aligned}
$}
\end{equation}
Además, si un argumento $A_i$ es un haz localmente libre, puede
suprimirse a condición de sustituir uno de los $A_j$ ($j\neq i$) por
$A_j\otimes A_i$, o uno de los $B_k$ por $B_k\otimes A_i'$ (donde
$A_i'=\Homsh_{\Osh_X}(A_i,\Osh_X)$), y se tiene una regla de cálculo análoga en
el caso en que uno de los argumentos $B_i$ es localmente libre. En particular, siempre
se puede suprimir un argumento que sea idéntico a $\Osh_X$. Si todos los
argumentos son localmente libres, salvo a lo sumo uno de los argumentos $B_i$, la regla
que acaba de enunciarse da un isomorfismo funtorial
\begin{equation}\tag{8.5}
  T_{r}^{s*}(A_1,\ldots,A_r;B_1,\ldots,B_s)
  =H^{*}\bigl(X,A_1'\otimes\cdots\otimes A_r'
    \otimes B_1\otimes\cdots\otimes B_s\bigr)
\end{equation}
(pues se reduce al caso en que $r=0$, $s=1$, donde esto es inmediato; también se
puede utilizar directamente la sucesión espectral de término inicial (8.2)). Para los
$\underline{T}_{r}^{s}$ se definen operaciones correspondientes a las anteriores.
Las relaciones entre las diversas operaciones así
introducidas son las mismas que para las operaciones análogas del cálculo
tensorial.

Sea $n$ la dimensión de $X$. Aplicando sucesivamente una composición
tensorial (8.3) y contracciones (8.4) sobre los argumentos repetidos, se
obtiene un acoplamiento
\begin{equation}\tag{8.6}
\resizebox{\textwidth}{!}{$
  (T_{r}^{s})^{p}(A_1,\ldots;B_1,\ldots)
  \times
  (T_{s}^{r})^{n-p}(B_1,\ldots;A_1,\ldots,A_r\otimes\Omega_X^{n})
  \longrightarrow H^{n}(X,\Omega_X^{n}).
$}
\end{equation}

\result{TEOREMA}{6} Si $X$ es una variedad proyectiva no singular, entonces
los acoplamientos (8.6) son dualidades.

Esto resulta de manera puramente formal del corolario del teorema 3. En efecto,
de este corolario resulta fácilmente que, si $K$ es un complejo de haces
algebraicos coherentes localmente libres, entonces la hipercohomología de $X$ con
respecto a $K$ está en dualidad con la hipercohomología de $X$ con respecto a
$K'\otimes\Omega_X^{n}$ mediante los acoplamientos naturales
\begin{equation}\tag{8.7}
  \mathbf R^{p}\Gamma(K)
  \times\mathbf R^{n-p}\Gamma(K'\otimes\Omega_X^{n})
  \longrightarrow\mathbf R^{n}\Gamma(\Omega_X^{n})
  =H^{n}(X,\Omega_X^{n}).
\end{equation}
Esto se ve utilizando la sucesión espectral de término inicial
$H^{p}(X,H^{q}(K))$\footnote{El texto impreso dice
$H^{p}(H^{q}(X,K))$; el término inicial es el de la sucesión espectral (8.2),
esto es, $H^{p}(X,H^{q}(K))$.} y la sucesión espectral análoga para
$K'\otimes\Omega_X^{n}$. Del resultado anterior se deduce el teorema 6
utilizando la definición (8.1).

\src{192}
\subsection*{OBSERVACIONES}

\noindent 1\textsuperscript{o} Para las definiciones anteriores al teorema 6,
no era necesario que $X$ fuese no singular, pues no era
indispensable trabajar únicamente con resoluciones finitas. Pero si $X$
es singular, ya no se podrá afirmar, a priori, que los
$(\underline{T}_{r}^{s})^{p}(A_1,\ldots;B_1,\ldots)$ sean haces
coherentes, pues en el complejo de haces
\[
  \Homsh_{\Osh_X}\bigl(L(A_1)\otimes\cdots,L(B_1)\otimes\cdots\bigr)
\]
habrá una infinidad de componentes de un grado total dado.

\medskip
\noindent 2\textsuperscript{o} Se comprueba fácilmente que, en las fórmulas
(8.1), se puede sustituir uno de los $L(B_i)$ por $B_i$. Teniendo en cuenta la
proposición 3, esto demuestra que se tiene
\begin{equation}\tag{8.8}
\left\{
\begin{aligned}
  \underline{T}_{1}^{1*}(A;B)&=\Extsh_{\Osh_X}^{*}(A,B),\\
  T_{1}^{1*}(A;B)&=\operatorname{Ext}_{\Osh_X}^{*}(X;A,B).
\end{aligned}
\right.
\end{equation}
En particular, tomando en (8.6) $r=s=1$ y $A_1=\Osh_X$, se recupera el
teorema 3. La fórmula (8.8) implica también
\[
  T_{0}^{1*}(B)=H^{*}(X,B),
  \qquad
  T_{1}^{0*}(A)=\operatorname{Ext}_{\Osh_X}^{*}(X;A,\Osh_X).
\]

\medskip
\noindent 3\textsuperscript{o} Se ve en la fórmula (8.1) que, en general, los
funtores $\underline{T}_{r}^{s*}$ y $T_{r}^{s*}$ tienen componentes de
grados positivos y negativos. Utilizando la observación 2, se ve que si la
dimensión de $X$ es $n$, entonces las componentes no nulas de
$\underline{T}_{r}^{s*}$ están comprendidas entre $-(s-1)n$ y $rn$, si $s>0$,
entre $0$ y $rn$, si $s=0$, y las componentes no nulas de $T_{r}^{s*}$
están comprendidas entre $-(s-1)n$ y $(r+1)n$, si $s>0$, entre $0$ y $(r+1)n$,
si $s=0$ (e incluso, salvo error, si $r>0$, entre $-(s-1)n$ y $rn$, resp. entre
$0$ y $rn$).

\section*{BIBLIOGRAFÍA}

\noindent[1] Cartier (Pierre). --- Les groupes
$\operatorname{Ext}^{s}(A,B)$, Séminaire A. Grothendieck : Algèbre homologique,
t. 1, 1957, n\textsuperscript{o} 3.

\noindent[2] Grothendieck (Alexandre). --- Sur quelques points d'algèbre
homologique, \emph{Tohoku Math. J.}, t. 9, 1957, p. 119--183.

\noindent[3] Serre (Jean-Pierre). --- Faisceaux algébriques cohérents,
\emph{Annals of Math.}, t. 61, 1955, p. 197--278.

\noindent[4] Serre (Jean-Pierre). --- Sur la dimension homologique des anneaux
et des modules noethériens, \emph{Proc. Intern. Symp. on Alg. Number Theory}
[1955, Tokyo et Nikko]. --- Tokyo, Science Council of Japan, 1956 ;
p. 175--189.

\clearpage
\src{193}
\section*{ADENDA}

Las dificultades señaladas en la observación de la página 13 están actualmente
completamente resueltas, gracias a una extensión de los teoremas de dualidad a
variedades cualesquiera (con singularidades arbitrarias). Para esta teoría, que puede
formularse en el marco general de la « teoría de esquemas », remitimos
a los « Éléments de Géométrie algébrique » de J. Dieudonné y A. Grothendieck,
actualmente en preparación. Se encontrarán algunas indicaciones en:

\begin{quote}
Grothendieck (A.). --- Cohomology theory of algebraic varieties,
\emph{Congrès int. des Mathématiciens}, 1958, Edinburgh (de próxima aparición).
\end{quote}

\begin{flushright}
[Abril de 1959]
\end{flushright}

\begin{center}\rule{36mm}{0.5pt}\end{center}

\def\FgaUnitKey{e149}\clearpage
\chapter*{Fe de erratas de la exposición n\textsuperscript{o} 149}
\addcontentsline{toc}{chapter}{Fe de erratas de la exposición n\textsuperscript{o} 149}
\setcounter{footnote}{0}
\src{298}

\noindent\textbf{1. Teorema de dualidad para los haces algebraicos
coherentes.}

\begin{center}
[Seminario Bourbaki, t. 9, 1956/57, n\textsuperscript{o} 149, 25 p.]
\end{center}

\noindent\textbf{Observación sobre la página 14.} --- Como señalé en mi
conferencia en el Congreso Internacional de Matemáticos de 1958,\footnote{
Grothendieck (Alexander). --- The cohomology theory of abstract algebraic
varieties, \emph{Proceedings of the International Congress of Mathematicians}
[1958, Edinburgh], p. 103--118. --- Cambridge, at the University Press, 1960.}
las cuestiones planteadas aquí están ahora completamente resueltas.

El lector encontrará más información sobre la dualidad de los haces
coherentes en \emph{loco citato}\footnotemark[1], p. 112--115, en EGA III,
2\textsuperscript{a} parte (en preparación) y, mientras tanto, en SGA 1962.
Un tratamiento más sistemático se encontrará en un capítulo posterior de EGA
(capítulo IX en el plan previsto).

\def\FgaUnitKey{182}\unit{Exposición 182 — Geometría formal y geometría algebraica}
\sid{182}
\src{193}
\seminar{(Mayo de 1959)}
\paperhead{GEOMETRÍA FORMAL Y GEOMETRÍA ALGEBRAICA}%
  {por Alexander GROTHENDIECK}

\fsection{1}{Esquemas}

Se sabe que un espacio algebraico afín definido sobre un cuerpo $k$ está
esencialmente determinado por su álgebra afín $A$ (anillo de las funciones
regulares definidas sobre $k$), pues los morfismos de espacios algebraicos
$X\longrightarrow Y$ corresponden biunívocamente a los homomorfismos de
$k$-álgebras $A(Y)\longrightarrow A(X)$. El álgebra afín correspondiente a un
espacio algebraico es una $k$-álgebra de tipo finito y, desde el punto de vista
«clásico», no tiene elementos nilpotentes; recíprocamente, toda álgebra de este
tipo se obtiene como álgebra afín de un espacio algebraico definido sobre $k$.
Se dispone así de un conocido diccionario que permite interpretar en términos
de álgebra conmutativa las situaciones relativas a espacios algebraicos afines.
Desde hace mucho se ha observado que de este modo se obtienen enunciados más
generales, pues por lo común ya no es necesario suponer que los anillos en juego
sean del tipo recién considerado, y la hipótesis noetheriana suele bastar. En
particular, haya o no un cuerpo de base dado, no hay razón para excluir el caso
en que estos anillos contengan elementos nilpotentes. Hasta ahora, los geómetras
se habían negado a tener en cuenta estas indicaciones y se han obstinado en
limitarse a considerar álgebras afines sin elementos nilpotentes, es decir,
espacios algebraicos cuyos haces estructurales no tienen elementos nilpotentes
(y, las más de las veces, incluso espacios algebraicos «absolutamente
irreducibles»). El conferenciante considera que esta actitud ha constituido un
serio obstáculo para el desarrollo de métodos verdaderamente naturales en
Geometría algebraica.

Sea $A$ un anillo conmutativo. Es bien sabido que el conjunto
$X=\operatorname{Spec}(A)$ de los ideales primos de $A$ está provisto de una
topología natural, la «topología de Zariski» o topología espectral. Por otra
parte, hay un haz de anillos conmutativos $\Osh_X$ sobre $X$, cuyo tallo en
$\mathfrak p\in X$ es el anillo localizado $A_{\mathfrak p}$ y cuyo anillo de
secciones se identifica con $A$. Así, $X$ se convierte en un \emph{espacio
anillado}, llamado \emph{espectro primo} de $A$. Un homomorfismo de anillos
$f:A\longrightarrow B$ define un morfismo de espacios anillados
$f':\operatorname{Spec}(B)\longrightarrow\operatorname{Spec}(A)$; la
aplicación subyacente de conjuntos no es otra que
$\mathfrak p\longmapsto f^{-1}(\mathfrak p)$. Los morfismos

\src{194}
de espacios anillados de $\operatorname{Spec}(B)$ en
$\operatorname{Spec}(A)$ obtenidos de este modo son exactamente aquellos para
los cuales los homomorfismos $\Osh_x\longrightarrow\Osh_y$ ($x=f'(y)$) son
locales (es decir, la imagen inversa del ideal maximal es el ideal maximal).

Se llama \emph{esquema afín} a un espacio anillado isomorfo a algún
$\operatorname{Spec}(A)$, y \emph{preesquema} a un espacio anillado localmente
afín, es decir, tal que cada punto tiene un entorno abierto que, con la
estructura inducida, es un esquema afín. Los \emph{morfismos} de preesquemas se
definen de manera evidente; localmente corresponden a homomorfismos de anillos.

Cuando se fija un preesquema $S$ y se consideran morfismos de preesquemas
$X\longrightarrow S$, el espacio $S$ desempeña el papel de un cuerpo o de un
anillo de base (o, mejor, de un espacio de base en una fibración). Se dice
entonces que $X$ es un $S$-preesquema; si $S=\operatorname{Spec}(A)$, esto
significa también que $\Osh_X$ es un haz de $A$-álgebras. Así, todo preesquema
puede considerarse de manera única como un $\mathbf Z$-preesquema. Naturalmente,
los $S$-preesquemas forman una categoría; además, se demuestra que en esta
categoría siempre existe el producto de dos objetos $X,Y$, que se denota por
$X\times_S Y$. Esta noción de producto permite definir el cambio de base de un
$S$-preesquema correspondiente a un morfismo $S'\longrightarrow S$: en efecto,
$X\times_S S'$ puede considerarse como un $S'$-preesquema.

Se dice que $X$ es \emph{separado} sobre $S$ si la diagonal de
$X\times_S X$ es cerrada. Se llama \emph{esquema} a un preesquema separado
sobre $\mathbf Z$; es entonces separado sobre cualquier base. Para simplificar,
en adelante solo hablaremos de esquemas y, además, los supondremos
\emph{noetherianos}, es decir, uniones finitas de abiertos afines que son
espectros de anillos noetherianos. Se dice que $X$ es \emph{de tipo finito
sobre} $S$ si, para todo abierto afín $U$ de $S$, su imagen inversa en $X$ es
una unión finita de abiertos afines cuyos anillos son álgebras de tipo finito
sobre el anillo de $U$. Son estos $S$-esquemas los que se prestan a un estudio
propiamente geométrico. En particular, para todo $s\in S$, la fibra
$f^{-1}(s)$ de $X$ sobre $s$ es un esquema algebraico sobre el cuerpo residual
$\kappa(s)$ del anillo local $\Osh_s$ de $s$ en $S$. Así, en cierta medida, $X$
puede considerarse como una familia de «espacios algebraicos» $f^{-1}(s)$, donde
el parámetro $s$ recorre $S$ (es decir, desde el punto de vista local, el
conjunto de los ideales primos de un anillo dado). Naturalmente, los
$\kappa(s)$ pueden tener características distintas. Si
$S=\operatorname{Spec}(k)$, donde $k$ es un cuerpo, se recupera esencialmente la
noción usual de «espacio algebraico», con la única diferencia de que ahora el
haz estructural puede tener elementos nilpotentes.

\src{195}
Inspirándose en nociones bien conocidas, se define la noción de
\emph{morfismo proyectivo} y, más generalmente, la de \emph{morfismo propio}.
Tal morfismo es de tipo finito; además, transforma partes cerradas en partes
cerradas y conserva esta propiedad bajo cualquier cambio de base.

Siendo $X$ un esquema (noetheriano, como siempre), el haz $\Osh_X$ es un
\emph{haz coherente de anillos} en el sentido de [2]. Por tanto, los haces
coherentes de módulos sobre $X$ son también los haces que localmente son
isomorfos al conúcleo de un homomorfismo
$\Osh_X^{m}\longrightarrow\Osh_X^{n}$.

\fsection{2}{Esquemas formales}

Sean $X$ un esquema y $X'$ una parte cerrada de $X$. Entonces existe un
subhaz coherente $J$ de $\Osh_X$ tal que
$X'=\operatorname{Supp}(\Osh_X/J)$ (e incluso existe uno máximo). Provisto de
$\Osh_X/J$, el espacio $X'$ se convierte en un esquema, denotado por $X_0$;
tal esquema se llama «subesquema cerrado de $X$». Para todo $n$ también se
puede considerar $X'$ provisto de $\Osh_X/J^{n+1}$, que se denota por $X_n$;
es un subpreesquema cerrado de $X$ cuyo conjunto subyacente sigue siendo $X'$,
pero que tiene otro haz estructural, a saber,
$\Osh_{X_n}=\Osh_X/J^{n+1}$. Evidentemente, los $\Osh_{X_n}$ forman un sistema
proyectivo de haces de anillos sobre $X'$, cuyo límite proyectivo
$\overline{\Osh}_X$ se llama \emph{completación formal de $\Osh_X$ a lo largo
de $X'$}. Provisto de este haz de anillos, $X'$ se llama \emph{completación
formal de $X$ a lo largo de $X'$}; es, pues, un espacio anillado, pero en
general no un esquema. Para todo haz coherente $F$ sobre $X$ se puede considerar
análogamente la completación formal $\overline F=\varprojlim_n F_n$ de $F$ a lo
largo de $X'$ (donde $F_n=F\otimes_{\Osh_X}\Osh_X/J^{n+1}$); es un haz de
módulos sobre $\overline X$. Sus secciones se llaman «secciones formales de $F$
a lo largo de $X'$» y se identifican con los elementos de
$\varprojlim_n\Gamma(X',F_n)$. Para $F=\Osh_X$ se obtienen las «funciones
holomorfas» de $X$ a lo largo de $X'$ en el sentido de Zariski (terminología que
no seguiremos por sus interferencias con la terminología clásica).

Se llama \emph{esquema formal} (se sobreentiende: noetheriano) a un espacio
topológico $\mathfrak X$, provisto de un haz de anillos topológicos
$\Osh_{\mathfrak X}$ que satisface la condición siguiente: hay un isomorfismo
de haces de anillos topológicos
$\Osh_{\mathfrak X}=\varprojlim_n\Osh_n$, donde los $\Osh_n$ forman un sistema
proyectivo de haces de anillos sobre $\mathfrak X$, cada uno de los cuales hace
de $\mathfrak X$ un esquema $\mathfrak X_n$, y tal que, para $m\geq n$, el
homomorfismo $\Osh_m\longrightarrow\Osh_n$ es sobreyectivo y tiene por núcleo
$J_m^{n+1}$, donde $J_m$ es

\src{196}
el núcleo de $\Osh_m\longrightarrow\Osh_0$. Se demuestra entonces que
$\Osh_{\mathfrak X}$ es un \emph{haz coherente de anillos locales
noetherianos}.

En virtud de las definiciones, una completación formal $\overline X$ como la
anterior es un esquema formal y, recíprocamente, todo esquema formal es
\emph{localmente} de este tipo. De hecho, dar un esquema formal \emph{afín} (es
decir, tal que $\mathfrak X_0$ sea afín, lo cual implica que todos los
$\mathfrak X_n$ lo son) equivale a dar un anillo topológico noetheriano,
separado y completo para la topología $J$-ádica.

Las definiciones habituales de morfismo, morfismo de tipo finito, morfismo
propio, etc., para los esquemas ordinarios se extienden sin dificultad a los
esquemas formales.

\fsection{3}{Los tres teoremas fundamentales}

Sea $f:X\longrightarrow Y$ un morfismo propio de esquemas (noetherianos, como
siempre), sea $Y'$ una parte cerrada de $Y$ y sea $X'$ su imagen inversa en
$X$; consideremos las completaciones formales correspondientes $\overline Y$ y
$\overline X$. Entonces $f$ induce un morfismo
$\overline f:\overline X\longrightarrow\overline Y$ de esquemas formales, que
además es propio. Si $F$ es un haz coherente sobre $X$, entonces $\overline F$
es un haz coherente sobre $\overline X$. En el teorema 1 se olvidan $X,Y,F$ y
solo se consideran el morfismo propio $\overline f$ de esquemas formales y el
haz coherente $\overline F$ sobre $\overline X$. (No obstante, el conferenciante
solo ha escrito una demostración completa para el caso en que se parte de
$X,Y,f,F$.)

\result{TEOREMA}{1 (teorema de finitud)}
\begin{enumerate}
  \item[\textit{i.}] Los $R^{q}\overline f_{*}(\overline F)$ son haces
    coherentes sobre $\overline Y$.
  \item[\textit{ii.}] Los homomorfismos naturales
    \[
      R^{q}\overline f_{*}(\overline F)
      \longrightarrow \varprojlim_n R^{q}f_{n*}(F_n)
    \]
    son isomorfismos.
\end{enumerate}

En este enunciado se supone elegido un subhaz coherente $J$ de $\Osh_Y$ que
define $Y'$; por imagen inversa se obtiene un subhaz coherente de $\Osh_X$ que
define $X'$ y, por consiguiente, se definen $F_n,X_n,Y_n$ y
$f_n:X_n\longrightarrow Y_n$ como en el n\textsuperscript{o} 2. Si se parte de
un morfismo propio cualquiera entre dos esquemas formales, son evidentes los
pequeños cambios que deben hacerse al explicitar estas notaciones.

El teorema 1 se refería únicamente a la «cohomología formal». El teorema siguiente

\src{197}
la relaciona con la «cohomología algebraica» y es análogo a un conocido teorema
de Serre [4] sobre la comparación entre cohomología algebraica y cohomología
analítica.

\result{TEOREMA}{2 (Primer teorema de comparación)} Los
$R^{q}f_{*}(F)$ son haces coherentes sobre $Y$ (lo cual es un caso particular
del teorema 1), y los homomorfismos naturales
\[
  \overline{R^{q}f_{*}(F)}
  \longrightarrow \varprojlim_n R^{q}f_{n*}(F_n)
\]
son isomorfismos.

\result{COROLARIO}{1} Se tienen isomorfismos canónicos:
\[
  \overline{R^{q}f_{*}(F)}=R^{q}\overline f_{*}(\overline F).
\]

Para $q=0$, este corolario es una generalización del «teorema fundamental de
las funciones holomorfas» de Zariski, del que deduciremos una generalización del
«teorema de conexión» de Zariski. Obsérvese además que, mientras el teorema 1
(\textit{ii}) es trivial para $q=0$, en modo alguno ocurre lo mismo con el
teorema 2 ni con su formulación equivalente (corolario 1). En efecto, la
demostración procede por inducción descendente sobre $q$ (es trivial para $q$
grande, pues entonces ambos miembros son nulos), de modo que el caso $q=0$
aparece como el último paso de la inducción y, por así decirlo, como el caso «más
difícil».

\result{COROLARIO}{2} Supongamos $Y=\operatorname{Spec}(A)$ y que $Y'$ está
definido por el ideal $J$ de $A$. Entonces, para todo haz coherente $F$ sobre
$X$, los $H^{q}(X,F)$ son $A$-módulos de tipo finito cuyas completaciones
$J$-ádicas son los $H^{q}(\overline X,\overline F)$.

Aplicando, por último, este corolario a
$H=\Homsh_{\Osh_X}(F,G)$, se obtiene:

\result{COROLARIO}{3} Supongamos $Y=\operatorname{Spec}(A)$ y que $Y'$ está
definido por el ideal $J$ de $A$. Sean $F,G$ dos haces coherentes sobre $X$;
entonces $\operatorname{Hom}(F,G)$ es un módulo de tipo finito sobre $A$, cuya
completación $J$-ádica se identifica con
$\operatorname{Hom}(\overline F,\overline G)$.

Naturalmente, la aplicación natural
$\operatorname{Hom}(F,G)\longrightarrow
\operatorname{Hom}(\overline F,\overline G)$ es la que asocia a un
homomorfismo $u:F\longrightarrow G$ su prolongación «por continuidad»
$\overline u:\overline F\longrightarrow\overline G$ (con lo cual
$\overline F$ se convierte en un funtor de $F$).

Supongamos ahora que $A$ es separado y completo para su topología $J$-ádica;
entonces los corolarios 2 y 3 precedentes dan:
\[
  H^{q}(X,F)=H^{q}(\overline X,\overline F),
  \qquad
  \operatorname{Hom}(F,G)
  =\operatorname{Hom}(\overline F,\overline G).
\]

\src{198}
Esta última identidad muestra que la categoría de los haces coherentes sobre
$X$ se identifica con una \emph{subcategoría} (cuyos morfismos son los
morfismos inducidos) de la categoría de los haces coherentes sobre
$\overline X$. De hecho, se tiene incluso:

\result{TEOREMA}{3} Para que un haz de módulos sobre $\overline X$ sea
coherente, es necesario y suficiente que sea isomorfo a un haz de la forma
$\overline F$, donde $F$ es un haz coherente sobre $X$ (determinado salvo
isomorfismo canónico en virtud del teorema 2, corolario 3). [Recordemos que
ahora $Y=\operatorname{Spec}(A)$, siendo $A$ un anillo topológico noetheriano,
separado y completo para la topología $J$-ádica.]

\result{COROLARIO}{1} Los subesquemas cerrados de $X$ corresponden
biunívocamente a los subesquemas formales cerrados de $\overline X$.

En efecto, corresponden a los subhaces coherentes de $\Osh_X$ y,
respectivamente, de $\overline{\Osh}_X$. Considerando los grafos de morfismos
como subesquemas cerrados, se deduce del corolario 1:

\result{COROLARIO}{2} Sean $X,Z$ dos esquemas propios sobre $A$ (anillo
noetheriano, separado y completo para la topología $J$-ádica). Entonces la
aplicación $g\longmapsto\overline g$ define una correspondencia biunívoca entre
los $Y$-morfismos de $X$ en $Z$ y los $\overline Y$-morfismos de $\overline X$
en $\overline Z$.

En otros términos, los esquemas algebraicos propios sobre $A$ aparecen como una
\emph{subcategoría} (cuyos morfismos son los morfismos inducidos) de la
categoría de los esquemas formales propios sobre $\overline Y$. Sin embargo,
debe advertirse que existen esquemas formales propios sobre $\overline Y$ que
no son «algebraizables», es decir, isomorfos a algún $\overline X$ con $X$
propio sobre $A$ (del mismo modo que existen variedades analíticas complejas
compactas que no provienen de variedades algebraicas definidas sobre el cuerpo
de los complejos). Tales esquemas formales aparecen de manera natural en la
«teoría de módulos». Señalemos, no obstante, un caso particular interesante en
el que un esquema formal es algebraizable:

\result{TEOREMA}{4} Sea $A$ un anillo noetheriano local completo de cuerpo
residual $k$, y sea $\mathfrak X$ un esquema formal propio sobre $A$ (provisto
de su topología $\mathfrak m(A)$-ádica). Se supone:
\begin{enumerate}
  \item[\textit{i.}] que los anillos locales de $\Osh_{\mathfrak X}$ son
    módulos planos sobre $A$ o, lo que equivale a lo mismo, si se dota a
    $\Osh_{\mathfrak X}$ y a $A$ de la filtración definida por las potencias del
    ideal maximal de $A$, para los graduados asociados se tiene la relación
    \[
      \operatorname{gr}(\Osh_{\mathfrak X})
      \simeq
      \operatorname{gr}^{0}(\Osh_{\mathfrak X})
      \otimes_k\operatorname{gr}(A).
    \]
  \item[\textit{ii.}] Considerando
    $\mathfrak X_0=\mathfrak X\otimes_A k$ como un esquema algebraico sobre
    $k$, se tiene
    \[
      H^2(\mathfrak X_0,\Osh_{\mathfrak X_0})=0.
    \]
  \item[\textit{iii.}] $\mathfrak X_0$ es proyectivo.
\end{enumerate}
Bajo estas condiciones, $\mathfrak X$ es algebraizable y, más precisamente, es
isomorfo a $\overline X$, donde $X$ es un $A$-esquema proyectivo.

Las condiciones \textit{(ii)} y \textit{(iii)} se satisfacen, en particular,
si $\mathfrak X_0$ es una curva simple sobre $k$, y el teorema 4 puede aplicarse
en particular en la «teoría de módulos» para las curvas de género dado\ldots{}
Indiquemos cómo se demuestra el teorema 4: se demuestra (cf. la proposición 3
más adelante) que \textit{(i)} y \textit{(ii)} implican que todo haz coherente
sobre $\mathfrak X_0$, localmente isomorfo al haz fundamental, se obtiene por
reducción a partir de un haz de la misma naturaleza sobre $\mathfrak X$. Se
parte entonces de un haz «amplio» sobre $\mathfrak X_0$ (existe uno por
\textit{(iii)}), se lo levanta a un haz invertible sobre $\mathfrak X$ y,
utilizando el teorema 1, se demuestra que un múltiplo de este último define una
inmersión de $\mathfrak X$ en la completación formal de un esquema
$\mathbf P_A^r$ («proyectivo tipo» de dimensión $r$ sobre $A$).

Para la demostración de los teoremas 1 a 3, véase [1].

\src{199}
\fsection{4}{Aplicación al teorema de conexión y al «Main theorem» de
Zariski}

Sea $f:X\longrightarrow Y$ un morfismo propio de esquemas. Por el teorema de
finitud, $f_*(\Osh_X)=\mathcal A$ es un haz coherente sobre $Y$; además, es un
haz de álgebras conmutativas, por lo que corresponde a un $Y$-esquema
$g:Y'\longrightarrow Y$ finito sobre $Y$ (definido por la condición de ser afín
sobre $Y$, es decir, que la imagen inversa de un abierto afín sea afín, y por
$g_*(\Osh_{Y'})=\mathcal A$). Es inmediato que $f$ se factoriza entonces
canónicamente como $f=gf'$, donde $f':X\longrightarrow Y'$ es un morfismo de
$X$ en $Y'$ que esta vez satisface $f'_*(\Osh_X)=\Osh_{Y'}$. Esta factorización
de $f$ se llama \emph{factorización de Stein} de $f$. Aplicando el primer
teorema de comparación, corolario 1, a $f'$ y a la parte de $Y'$ reducida a un
punto $y'$, se obtiene que $f'^{-1}(y')=X'$ es conexo (de otro modo, las
secciones formales de $X$ a lo largo de $X'$ no formarían un anillo local,
¡mientras que la completación de
$f'_*(\Osh_X)_{y'}=\Osh_{Y',y'}$ es local!). Hemos demostrado:

\result{TEOREMA}{5 («Teorema de conexión» de Zariski)} Sea
$f:X\longrightarrow Y$ un morfismo propio. Entonces $f$ se factoriza de manera
única (salvo isomorfismo) como $f=gf'$, donde $g:Y'\longrightarrow Y$ es finito
y $f':X\longrightarrow Y'$ satisface $f'_*(\Osh_X)=\Osh_{Y'}$ (de donde
$g_*(\Osh_{Y'})=f_*(\Osh_X)$). Las fibras de $f'$ son conexas; es decir, el
conjunto de las componentes conexas de una fibra $f^{-1}(y)$ de $f$ está en
correspondencia biunívoca con el conjunto de los puntos de $Y'$ sobre $y$, esto
es, con el conjunto de los ideales maximales de $f_*(\Osh_X)_y$.

De aquí se deducen inmediatamente las variantes habituales del teorema de
conexión. Señalemos aquí únicamente el

\src{200}
\result{COROLARIO}{1} Para que un punto $x$ de $X$ esté aislado en su fibra
$f^{-1}(y)$, es necesario y suficiente que la fibra $f'^{-1}(y')$ (donde
$y'=f'(x)$) se reduzca a $x$, o bien que $f'$ induzca un isomorfismo de un
entorno de $x$ sobre un entorno de $y'$. El conjunto de estos puntos es un
abierto $U$, y $f'$ induce un isomorfismo de $U$ sobre un abierto de $Y'$.

Para demostrar que $f'$ es un isomorfismo local en $x$, se observa que $f'$
induce un isomorfismo $\Osh_{y'}\longrightarrow\Osh_x$, como se ve gracias a
$f'_*(\Osh_X)=\Osh_{Y'}$ y al hecho de que (siendo $f'$ una aplicación cerrada
cuya fibra en $y'$ se reduce a $x$) los $f'^{-1}(V)$ recorren un sistema
fundamental de entornos de $x$ cuando $V$ recorre un sistema fundamental de
entornos de $y'$. De aquí se deduce enseguida el resultado siguiente, debido a
Chevalley en el caso «geométrico».

\result{COROLARIO}{2} Para que $f$ sea un morfismo finito, es necesario y
suficiente que sea propio y que sus fibras sean finitas.

Si así ocurre, $f'$ es efectivamente un isomorfismo por lo anterior.

Sea $f:X\longrightarrow Y$ un morfismo no necesariamente propio, pero
supongamos que $X$ está contenido como abierto en un $Y$-esquema propio
$\widetilde f:\widetilde X\longrightarrow Y$ (lo cual ocurre, en particular,
si $f$ es cuasiproyectivo). Aplicando el corolario 1, se ve que $\widetilde f'$
induce un isomorfismo del conjunto $U$ de los puntos de $X$ aislados en su fibra
sobre un abierto de $Y'$ (y que $U$ es, en efecto, abierto). Se deduce la
siguiente variante global del «Main Theorem» de Zariski.

\result{TEOREMA}{6 («Main Theorem» de Zariski)} Sea
$f:X\longrightarrow Y$ un morfismo de tipo finito. Entonces el conjunto $U$ de
los puntos de $X$ que están aislados en su fibra es abierto y, si $f$ es
separado,\footnote{La fe de erratas de 1962 sustituye aquí la hipótesis impresa
«$f$ cuasiproyectivo» por la hipótesis más débil «$f$ separado», en virtud de
SGA VIII, 6.2: todo morfismo cuasifinito y separado es cuasiproyectivo.} $U$ es
$Y$-isomorfo a una parte abierta de un esquema $Y'$ finito sobre $Y$.

Como un morfismo de tipo finito es localmente afín y, \textit{a fortiori},
localmente cuasiproyectivo, del teorema 6 se deducen enseguida las variantes
locales usuales del Main Theorem.

\fsection{5}{Aplicación al estudio cohomológico de los morfismos propios y
planos}

Sea $f:X\longrightarrow Y$ un morfismo propio, y sea $F$ un haz coherente sobre
$X$; se supone que $F$ es $Y$-plano, es decir, los $F_x$ son módulos planos sobre los
anillos $\Osh_y$ ($y=f(x)$). Esto significa también que, para todo $y\in Y$, si
filtramos $F$ a lo largo de la fibra $f^{-1}(y)$ mediante los
$\mathfrak m_y^nF$ (siendo $\mathfrak m_y$ el ideal maximal de $\Osh_y$), el
graduado asociado es isomorfo a
$(F/\mathfrak m_yF)\otimes_{\kappa(y)}\operatorname{gr}(\Osh_y)$; en otros
términos, se tiene
\[
  \mathfrak m_y^nF/\mathfrak m_y^{n+1}F
  =F_y\otimes_{\kappa(y)}
  (\mathfrak m_y^n/\mathfrak m_y^{n+1})
\]
para todo entero $n$, donde $X_y$ denota la fibra $f^{-1}(y)$ considerada como
esquema propio sobre el cuerpo residual $\kappa(y)$ de $y$, y $F_y$ el haz
$F/\mathfrak m_yF$ inducido por $F$ sobre $X_y$. Teniendo en cuenta este
isomorfismo y el teorema 2, se obtienen cotas y, en ocasiones, cálculos de los
$R^qf_*(F)$ en un entorno de $y$, conociendo la cohomología de $X_y$ con
coeficientes en $F_y$. El teorema 2 adopta aquí la forma
\[
  \widehat{R^qf_*(F)_y}
  =\varprojlim_n H^q(X_y,F/\mathfrak m_y^nF).
\]
Señalaremos únicamente la consecuencia siguiente:

\src{201}
\result{PROPOSICIÓN}{1} Sean $f:X\longrightarrow Y$ un morfismo propio y $F$
un haz sobre $X$, coherente y $Y$-plano. Sea $y\in Y$, sea $q$ un entero, y
supongamos que $H^q(X_y,F_y)=0$. Entonces $R^qf_*(F)$ es nulo en un entorno de
$y$ y, para todo $n$, el homomorfismo natural
\[
  R^{q-1}f_*(F)_y\longrightarrow
  H^{q-1}(X_y,F/\mathfrak m_y^nF)
\]
es sobreyectivo.

En particular, si $f$ es un morfismo plano (es decir, $\Osh_X$ es $Y$-plano),
entonces todo haz coherente localmente libre $F$ sobre $X$ es $Y$-plano. Sean
$F,G$ dos haces de este tipo; aplicando la proposición 1 a
$\Homsh_{\Osh_X}(F,G)$ y a $q=1$, se obtiene:

\result{TEOREMA}{7} Sea $f$ un morfismo propio y plano, sean $F,G$ dos haces
coherentes localmente libres sobre $X$, sea $y\in Y$ y supongamos que
\[
  H^1\bigl(X_y,\Homsh_{\Osh_{X_y}}(F_y,G_y)\bigr)=0,
\]
entonces todo homomorfismo $u_0:F_y\longrightarrow G_y$ está inducido por un
homomorfismo
\[
  u:F|V\longrightarrow G|V,
  \qquad V=f^{-1}(U),
\]
donde $U$ es un entorno de $y$.

\result{COROLARIO}{} Si $u_0$ es un isomorfismo (resp. un monomorfismo, un
epimorfismo), entonces lo mismo ocurre con $u$ para $U$ suficientemente
pequeño.

En particular:

\src{202}
\result{COROLARIO}{2} Sea $E_0$ un haz coherente localmente libre sobre $X_y$
tal que
\[
  H^1\bigl(X_y;\Homsh_{\Osh_{X_y}}(E_0,E_0)\bigr)=0,
\]
entonces dos haces localmente libres cuyas restricciones a $X_y$ sean isomorfas
a $E_0$ son isomorfos en un entorno de $X_y$.

Así:

\result{COROLARIO}{3} Supongamos $H^1(X_y,\Osh_{X_y})=0$. Entonces dos haces
invertibles sobre $X$ (es decir, localmente isomorfos a $\Osh_X$) cuyas
restricciones a $X_y$ sean isomorfas son isomorfos en un entorno de $X_y$.\footnote{El texto impreso no precisa el entorno. Se restituye aquí el alcance local del corolario 2, utilizado también en la demostración de la proposición 2.}

Se deduce:

\result{PROPOSICIÓN}{2} Sean $Y$ un esquema conexo y $E$ un haz coherente localmente libre de rango al menos $2$ sobre $Y$; consideremos el fibrado de espacios proyectivos
$X=\mathbf P(E)$ asociado a $E$, provisto de su conocido haz invertible
$\Osh_X(1)$. Entonces todo haz invertible $L$ sobre $X$ es isomorfo a un haz de
la forma
\[
  f^*(L')\otimes\Osh_X(n),
\]
donde $L'$ es un haz invertible sobre $Y$ y $n$ un entero. Este último está
determinado de manera única, y $L'$ está determinado salvo isomorfismo.

El corolario 3 precedente demuestra que $L$ es isomorfo a algún $\Osh_X(n)$ en
un entorno de cada fibra. El resto es casi formal.

La proposición 2 permite determinar los $Y$-morfismos de $X=\mathbf P(E)$ en
otro fibrado proyectivo. En particular, se obtiene:

\result{COROLARIO}{} Sea $u$ un automorfismo de $X=\mathbf P(E)$. Entonces
existen un haz invertible $L'$ sobre $Y$ y un isomorfismo $v$ de $E$ sobre
$E\otimes L'$ tales que $u$ es el isomorfismo correspondiente
\[
  \mathbf P(E)\xrightarrow{\sim}\mathbf P(E\otimes L')=\mathbf P(E);
\]
el par $(v,L')$ está determinado salvo isomorfismo.

Sea $\Gamma$ el conjunto de las clases de fibrados invertibles $L'$ sobre $Y$
tales que $E\otimes L'$ sea isomorfo a $E$. Sus elementos son de torsión, pues,
si $n$ es el rango de $E$, se ve (tomando las potencias exteriores $n$-ésimas)
que debe tenerse $L'^{\otimes n}\simeq\Osh_Y$. El corolario puede expresarse
entonces diciendo que se tiene una sucesión exacta de grupos
\[
  e\longrightarrow
  \operatorname{Aut}(E)/\Gamma(Y,\Osh_Y^*)
  \longrightarrow \operatorname{Aut}_Y(X)
  \longrightarrow \Gamma\longrightarrow e
\]
(deducida, por lo demás, de la sucesión exacta de cohomología obtenida a partir
de la siguiente sucesión exacta de haces de grupos
\[
  e\longrightarrow \Osh_Y^*
  \longrightarrow \underline{\operatorname{Aut}}(E)
  \longrightarrow \underline{\operatorname{Aut}}_Y(X)
  \longrightarrow e,
\]
donde $\Osh_Y^*$ es el haz de las «unidades» de $\Osh_Y$, identificado con el
centro de $\underline{\operatorname{Aut}}(E)$).

\src{203}
\fsection{6}{Aplicación a teoremas de existencia y unicidad para haces y
esquemas sobre un anillo completo para la topología $J$-ádica}

El teorema 7 daba un resultado de unicidad para haces coherentes localmente
libres utilizando los teoremas 1 y 2. Utilizando el teorema 3, deducimos ahora
teoremas de existencia de haces, de morfismos de esquemas o de esquemas. En lo
sucesivo, $A$ denota un anillo noetheriano local, separado y completo. El método
general consiste siempre en efectuar construcciones formales, lo cual consiste
esencialmente en hacer geometría algebraica sobre un anillo artiniano, y extraer
de ellas conclusiones de naturaleza «algebraica» mediante los tres teoremas
fundamentales.

\result{PROPOSICIÓN}{3} Sea $\mathfrak X$ un esquema formal propio y plano
sobre $A$, y sea $F_0$ un haz localmente libre sobre $X_0$ tal que
\[
  H^2\bigl(X_0,\Homsh_{\Osh_{X_0}}(F_0,F_0)\bigr)=0.
\]
Entonces existe un haz localmente libre $F$ sobre $\mathfrak X$ que induce
sobre $X_0$ un haz isomorfo a $F_0$. (Además, este $F$ es único salvo
isomorfismo si
$H^1(X_0,\Homsh_{\Osh_{X_0}}(F_0,F_0))=0$.)

Se construyen sucesivamente haces localmente libres $F_n$ sobre los $X_n$, cada
uno de los cuales induce el anterior. La construcción de $F_n$ encuentra una
obstrucción en
\[
  H^2\bigl(X_0,\Homsh_{\Osh_{X_0}}(F_0,F_0)\bigr)
  \otimes_{A/J}(J^n/J^{n+1}),
\]
que es, por tanto, nula por hipótesis. Utilizando ahora el teorema 3, se
obtiene:

\result{COROLARIO}{1} Sea $X$ un esquema propio y plano sobre $A$, y sea $F_0$
como arriba. Entonces existe un haz localmente libre $F$ sobre $X$ que induce
sobre $X_0$ un haz isomorfo a $F_0$. Además, este $F$ es único salvo
isomorfismo si
$H^1(X_0,\Homsh_{\Osh_{X_0}}(F_0,F_0))=0$.

Sea $X_0$ un esquema de tipo finito sobre el cuerpo $k$; se supone que $X_0$ es
simple (por lo cual entendemos absolutamente simple) sobre $k$, pero no
necesariamente propio sobre $k$. Sea $A$ un anillo artiniano local de cuerpo
residual $k$. Nos interesa hallar los esquemas $X$ planos sobre $A$ tales que
$X\otimes_A k=X_0$ (este es el punto de partida de la «teoría de módulos» o de
las «variaciones de estructura» de $X_0$). Es equivalente dar un tal $X$, o dar
sobre el espacio topológico $X_0$ un haz $\Osh_X$ provisto de las estructuras
siguientes:
\begin{enumerate}
  \item[\textit{i.}] Es un haz de $A$-álgebras.
  \item[\textit{ii.}] Está provisto de un homomorfismo de aumentación
    $\Osh_X\longrightarrow\Osh_{X_0}$ (compatible con las estructuras de
    $A$-álgebras); estos datos están sujetos a las condiciones siguientes: la
    aumentación induce un isomorfismo
    $\Osh_X\otimes_A k\xrightarrow{\sim}\Osh_{X_0}$; $\Osh_X$ es plano sobre
    $A$, es decir, el graduado asociado a $\Osh_X$, filtrado por las potencias
    del ideal maximal $\mathfrak m$ de $A$, es isomorfo a
    $\operatorname{gr}^0(\Osh_X)\otimes_k\operatorname{gr}(A)$; es decir, se
    tienen isomorfismos
    \[
      \mathfrak m^n\Osh_X/\mathfrak m^{n+1}\Osh_X
      =\Osh_{X_0}\otimes_k(\mathfrak m^n/\mathfrak m^{n+1}).
    \]
\end{enumerate}
El hecho fundamental es el siguiente:

\src{204}
\result{TEOREMA}{8} Sea $X_0$ un esquema de tipo finito y simple sobre el
cuerpo $k$, y supongamos que $X_0$ es afín. Sea $A$ un anillo local artiniano de
cuerpo residual $k$. Entonces existe un $A$-esquema $X$ plano sobre $A$ tal que
$X\otimes_A k=X_0$, y dos esquemas de este tipo son necesariamente isomorfos.

Obsérvese que el isomorfismo en cuestión no es canónico, pues en general $X$
tendrá $A$-automorfismos no triviales que inducen la identidad sobre $X_0$. Por
otra parte, en general no hay una elección «canónica» de un $X$ que satisfaga
las condiciones dadas, salvo cuando $A$ es una $k$-álgebra (caso de
características iguales), en cuyo caso puede tomarse $X=X_0\otimes_k A$, es
decir, $\Osh_X=\Osh_{X_0}\otimes_k A$ (sea o no afín $X_0$). En el caso de
características distintas, ignoro en general, cuando $X_0$ no es afín, si se
puede «levantar» $X_0$ a un $X$ definido sobre $A$. Sin embargo, sea $n$ un
entero $>0$, pongamos $A_{n-1}=A/\mathfrak m^n$, y supongamos que se ha
levantado $X_0$ a un $A_{n-1}$-esquema plano $X_{n-1}$; se pretende levantar
$X_{n-1}$ a un $A_n$-esquema plano $X_n$. Por el teorema 8 ya se sabe que esto
es posible localmente; por otra parte, se comprueba fácilmente que, si $U_n$
levanta un abierto $U_{n-1}$ de $X_{n-1}$, entonces el haz de los grupos de
automorfismos de $U_n$ (que inducen la identidad sobre $U_{n-1}$) es
canónicamente isomorfo a
\[
  \Theta_{X_0/k}\otimes_k(\mathfrak m^n/\mathfrak m^{n+1})
\]
restringido a $U_{n-1}$ y, en particular, es conmutativo
($\Theta_{X_0/k}$ denota el haz de gérmenes de $k$-derivaciones sobre $X_0$).
Se sigue fácilmente que hay una obstrucción para construir $X_n$ que levante
$X_{n-1}$, la cual se encuentra en
\[
  H^2(X_0,\Theta_{X_0/k})\otimes_k
  \mathfrak m^n/\mathfrak m^{n+1}.
\]
Por consiguiente:

\result{COROLARIO}{1} Sea $X_0$ un esquema de tipo finito simple sobre $k$, y
supongamos $H^2(X_0,\Theta_{X_0/k})=0$. Entonces, para todo anillo local
artiniano $A$ de cuerpo residual $k$, existe un $A$-esquema plano $X$ tal que
$X\otimes_A k=X_0$.

Además, si se ha podido hallar un $X$ plano sobre $A$ que levante $X_0$,
entonces, en virtud del teorema 8, el conjunto de las clases (salvo isomorfismo)
de $A$-esquemas planos que levantan $X_0$ se identifica con
\[
  H^1\bigl(X_0,\underline{\operatorname{Aut}}(X)\bigr),
\]
donde, naturalmente, $\underline{\operatorname{Aut}}(X)$ denota el haz de
gérmenes de automorfismos del haz de $A$-álgebras $\Osh_X$ compatibles con la
aumentación. La filtración de $\Osh_X$ define una filtración de
$\underline{\operatorname{Aut}}(X)$, y el cociente de este haz por el $n$-ésimo
subgrupo de la filtración es $\underline{\operatorname{Aut}}(X_n)$; el graduado
asociado a esta filtración es conmutativo y se identifica con
$\Theta_{X_0/k}\otimes_k\operatorname{gr}(A)$. En particular, si
$\mathfrak m^{n+1}$ es la primera potencia de $\mathfrak m$ que es nula,
entonces $F^n(\underline{\operatorname{Aut}}(X))$ (el último término de la
filtración) pertenece al centro de $\underline{\operatorname{Aut}}(X)$ y es
isomorfo a $\Theta_{X_0/k}\otimes_k\mathfrak m^n$; es también el haz de
gérmenes de automorfismos de $X$ que inducen la identidad sobre
$X_{n-1}=X\otimes_A A/\mathfrak m^n$. Utilizando estos resultados se obtienen
enseguida los enunciados siguientes:

\src{205}
\result{COROLARIO}{2} Sea $X_0$ un esquema de tipo finito simple sobre $k$, sea
$A$ un anillo local artiniano de cuerpo residual $k$ e ideal maximal
$\mathfrak m$, supongamos $\mathfrak m^{n+1}=0$, sea
$A_{n-1}=A/\mathfrak m^n$ y, por último, sea $X_{n-1}$ un
$A_{n-1}$-esquema plano tal que $X_{n-1}\otimes_A k=X_0$. Entonces el conjunto
de las clases (salvo un isomorfismo que induzca la identidad sobre $X_{n-1}$)
de $A$-esquemas planos $X_n$ tales que
$X_n\otimes_A A_{n-1}=X_{n-1}$ es vacío o es un espacio homogéneo principal
bajo
\[
  H^1(X_0,\Theta_{X_0/k})\otimes_k\mathfrak m^n.
\]

(Obsérvese que, en general, no hay un elemento origen privilegiado en este
último espacio, pues no hay una manera privilegiada de levantar $X_{n-1}$ a
$X_n$).

\result{COROLARIO}{3} Sea $X_0$ un esquema de tipo finito simple sobre $k$, y
supongamos $H^1(X_0,\Theta_{X_0/k})=0$. Entonces, para todo anillo local
artiniano $A$ de cuerpo residual $k$, existe a lo sumo (salvo isomorfismo) un
$A$-esquema plano $X$ tal que $X\otimes_A k=X_0$.

Los corolarios 1 y 3 implican inmediatamente los enunciados, en apariencia más
generales, que se obtienen suponiendo solamente que $A$ es un anillo local
noetheriano completo de cuerpo residual $k$, con tal de introducir allí $X$
como un esquema formal sobre $A$:

\result{TEOREMA}{9} Sean $k$ un cuerpo y $X_0$ un esquema de tipo finito simple
sobre $k$. Para todo anillo local noetheriano completo $A$ de cuerpo residual
$k$, sea $F(A)$ el conjunto de las clases (salvo un isomorfismo que induzca la
identidad sobre $X_0$) de esquemas formales $\mathfrak X$ sobre $A$, de tipo
finito y planos sobre $A$, tales que $\mathfrak X\otimes_A k=X_0$. Con estas
notaciones:
\begin{enumerate}
  \item[\textit{i.}] Si $H^1(X_0,\Theta_{X_0/k})=0$, entonces, para todo $A$,
    $F(A)$ tiene a lo sumo un elemento.
  \item[\textit{ii.}] Si $H^2(X_0,\Theta_{X_0/k})=0$, entonces, para todo $A$,
    $F(A)$ tiene al menos un elemento.
\end{enumerate}

\src{206}
\result{COROLARIO}{1} Supongamos que $X_0$ es propio sobre $k$. Bajo la
condición \textit{(i)}, para todo $A$ existe a lo sumo (salvo un isomorfismo que
induzca la identidad sobre $X_0$) un esquema $X$ propio y plano sobre $A$ tal
que $X\otimes_A k=X_0$.

Se utiliza el teorema 3, corolario 2. Por ejemplo:

\result{COROLARIO}{2} Si $X$ es un esquema propio y plano sobre $A$ tal que
$X\otimes_A k$ es isomorfo al esquema proyectivo tipo $\mathbf P_k^r$ de
dimensión $r$ sobre $k$, entonces $X$ es isomorfo a $\mathbf P_A^r$.

(También se puede deducir este resultado de la proposición 3, corolario 1).

\result{COROLARIO}{3} Sea $X_0$ un esquema simple y proyectivo sobre $k$, y
supongamos que se tiene
\[
  H^2(X_0,\Osh_{X_0})=H^2(X_0,\Theta_{X_0/k})=0.
\]
Entonces, para todo $A$, existe un esquema proyectivo y plano $X$ sobre $A$ tal
que $X\otimes_A k=X_0$.

Se combinan el teorema 9\textit{(ii)} y el teorema 4. En particular:

\result{COROLARIO}{4} Sea $X_0$ el esquema de una curva algebraica completa y
simple sobre $k$. Entonces, para todo anillo local noetheriano completo $A$ de
cuerpo residual $k$, existe un «esquema de curva simple» $X$ sobre $A$ tal que
$X\otimes_A k=X_0$.

\result{OBSERVACIONES}{}

1) Los corolarios 3 y 4 son especialmente interesantes si $k$ es de
característica $p\ne0$, tomando como $A$ un anillo de valoración discreta de
característica $0$ y cuerpo residual $k$; por ejemplo, el «menor $A$ posible»,
a saber, aquel cuyo ideal maximal está generado por $p$. (De hecho, por los
teoremas de Cohen, basta conocer los corolarios 3 y 4 para un anillo $A$ de este
tipo). Señalemos a este respecto que, según los especialistas, no se sabe si
existen esquemas sobre un cuerpo $k$ que no sean la reducción módulo $p$ de un
esquema plano definido sobre un tal anillo $A$. Al menos, los resultados de este
número proporcionan un medio de investigar sistemáticamente esta cuestión. Se
debería comenzar por examinar si la primera obstrucción que aparece en
$H^2(X_0,\Theta_{X_0/k})$ es necesariamente nula.

\emph{Complemento de 1962.} --- J.-P. Serre ha construido una variedad
proyectiva no singular, de dimensión 3, sobre un cuerpo algebraicamente cerrado
$k$ de característica $p>0$, que no se obtiene por reducción de un esquema
propio sobre un anillo local íntegro de cuerpo residual $k$ y cuerpo de
fracciones de característica cero. Mumford habría encontrado un ejemplo análogo
con una superficie proyectiva no singular.

2) Obsérvese que el teorema 3, y la técnica correspondiente, solo es válido para
un anillo de base completo (local, para fijar ideas).

Para pasar de resultados conocidos para la completación de un anillo local a
los resultados correspondientes para ese mismo anillo local, se necesitaría un
cuarto «teorema fundamental», cuyo enunciado definitivo aún está por hallar.

\src{207}
3) Compárense los resultados de este número (en particular, los corolarios 1 y
2 precedentes) y del siguiente con los resultados de Kodaira--Spencer sobre la
variación de las estructuras complejas. Utilizando el teorema conjetural al que
se acaba de aludir, bajo las condiciones del corolario 1, pero sin suponer ya
que $A$ sea completo, se debería poder concluir que existe un anillo $A'$ que
contiene a $A$, finito y no ramificado sobre $A$, tal que $X\otimes_A A'$ y
$X'\otimes_A A'$ son $A'$-isomorfos (siendo $X,X'$ dos $A$-esquemas propios y
planos dados, tales que $X\otimes_A k=X'\otimes_A k=X_0$). Esto es al menos lo
que puede demostrarse cuando $X_0=\mathbf P_k^r$, utilizando el corolario de la
proposición 2. En todo caso, cuando $H^1(X_0,\Theta_{X_0/k})=0$, puede
demostrarse que las fibras de $X,X'$ sobre todo punto $y$ de
$Y=\operatorname{Spec}(A)$ son isomorfas, al menos después de pasar a la
clausura algebraica del cuerpo residual $\kappa(y)$. (Se tiene un resultado
local, aparentemente más fuerte, sin suponer ya necesariamente que $A$ sea
local). En cuanto a las «variaciones de estructura» del espacio proyectivo,
señalemos además la cuestión siguiente, sugerida por un problema correspondiente
de Kodaira--Spencer. Sea $X$ un esquema propio y plano sobre el anillo local
íntegro $A$ de cuerpo de fracciones $K$ y cuerpo residual $k$, y supongamos que
$X\otimes_A K$ es isomorfo a $\mathbf P_K^r$; ¿es cierto que
$X\otimes_A k=X_0$ es isomorfo a $\mathbf P_k^r$ (al menos sobre la clausura
algebraica de $k$)? En este problema se puede suponer que $A$ es un anillo de
valoración discreta completo. Se plantea un problema análogo cuando $X_0$ es
una variedad abeliana.

\emph{Complemento de 1962.} --- El conferenciante es ahora menos optimista
respecto de los resultados conjeturados arriba. No obstante, Hironaka ha
resuelto afirmativamente el problema relativo a las variaciones de estructura
del espacio proyectivo, y Koizumi el problema análogo para las variedades
abelianas.

\fsection{7}{Aplicación a la «Teoría de módulos»}

Como el propio conferenciante apenas acaba de abordar esta teoría, nos veremos
obligados a limitarnos a algunas indicaciones. Para simplificar, nos situamos
sobre un cuerpo $k$, es decir, trabajamos en igual característica, aunque el
teorema 8 permite tratar también el caso general sin cambio esencial, al
parecer. Actualmente no hemos superado la etapa «formal»; sin embargo, el
conferenciante espera poder construir a partir de ella verdaderos esquemas de
módulos en ciertos casos y, en particular, construir para cada entero $g$ un
esquema sobre los enteros que desempeñe el papel de esquema de módulos universal
para las curvas simples de género $g$.

\emph{Complemento de 1962.} --- Mumford ha construido los esquemas de módulos
para las curvas de género $g$ (cf. SHMI). El teorema 10 demuestra además que
los «esquemas de Jacobi de nivel $n$» de la teoría de Módulos son no singulares,
e incluso simples sobre $\mathbf Z$.

\src{208}
Retomemos la situación y las notaciones del teorema 9, suponiendo ahora que $A$
es un álgebra local de rango finito sobre $k$, que para simplificar se supone
algebraicamente cerrado. Entonces $F(A)$ puede considerarse como un funtor
covariante de $A$, con valores en la categoría de conjuntos: un homomorfismo de
$k$-álgebras $A\longrightarrow B$ define una aplicación
$F(A)\longrightarrow F(B)$, pues todo $A$-esquema plano $X$ tal que
$X\otimes_A k=X_0$ da lugar a un $B$-esquema $X\otimes_A B$ con las mismas
propiedades. Supongamos que se puede hallar una $k$-álgebra local noetheriana
completa $\mathcal O$ y un isomorfismo funtorial
\begin{equation}
  \tag{*}
  \operatorname{Hom}(\mathcal O,A)\xrightarrow{\sim}F(A)
\end{equation}
(donde el primer miembro denota los homomorfismos de $k$-álgebras). Se ve
fácilmente que un tal $\mathcal O$ está determinado salvo isomorfismo canónico;
se llama entonces al espectro formal $\mathfrak Y$ de $\mathcal O$ (es decir,
al espacio topológico reducido a un punto y provisto de un haz de anillos
topológicos reducido a $\mathcal O$) el \emph{esquema formal de módulos} para
$X_0$. (Obsérvese que no existe necesariamente). Sea $\mathfrak r$ el ideal
maximal de $\mathcal O$; para cada $n$, sea
$\mathcal O_n=\mathcal O/\mathfrak r^{n+1}$ (por tanto $\mathcal O_0=k$).
Entonces el homomorfismo canónico $\mathcal O\longrightarrow\mathcal O_n$ es un
elemento de $\operatorname{Hom}(\mathcal O,\mathcal O_n)$, por lo que define un
elemento de $F(\mathcal O_n)$, es decir, un $\mathcal O_n$-esquema plano $X_n$
cuya reducción módulo $\mathfrak r$ es $X_0$. Estos $X_n$ se deducen unos de
otros por extensión de escalares (es decir, aquí por reducciones), de donde
resulta que provienen de un esquema formal $\mathfrak X$ bien determinado sobre
el esquema formal de módulos $\mathfrak Y$; $\mathfrak X$ es plano sobre
$\mathfrak Y$ y $\mathfrak X_0=X_0$. Como se ve enseguida, el isomorfismo (*)
se obtiene entonces asociando a todo homomorfismo de $k$-álgebras
$\mathcal O\longrightarrow A$ la clase del $A$-esquema
$\mathfrak X\otimes_{\mathcal O}A$ (es decir, a todo morfismo de $k$-esquemas
$Y'=\operatorname{Spec}(A)\longrightarrow\mathfrak Y$ se asocia el
$Y'$-esquema $\mathfrak X\otimes_{\mathfrak Y}Y'$ obtenido por cambio de base).
Además, se ve que el isomorfismo (*) y su descripción precedente siguen siendo
válidos si solo se supone que $A$ es una $k$-álgebra local noetheriana y
completa (no necesariamente artiniana). Naturalmente, como de costumbre,
$\mathcal O$ puede tener perfectamente \textit{a priori} elementos nilpotentes,
y parece probable que existan casos en que $\mathcal O$ sea él mismo artiniano
sin ser idéntico a $k$. Esto muestra hasta qué punto el punto de vista de
Kodaira--Spencer (que se limita a tomar como $A$ anillos regulares) es
inadecuado \textit{a priori} en el caso general.

\src{209}
Queda por dar condiciones suficientes para que exista un esquema formal de
módulos para $X_0$, que se supone propio sobre $k$. En general, es fácil dar
condiciones necesarias y suficientes sencillas sobre un funtor
$A\longmapsto F(A)$ (de $k$-álgebras locales de rango finito, con valores en
los conjuntos) para que pueda escribirse en la forma
$\operatorname{Hom}(\mathcal O,A)$ para un $\mathcal O$ adecuado. No las
detallaremos aquí. Señalemos solamente que, en el caso que nos ocupa, estas
condiciones imponen sobre $X_0$ condiciones cohomológicas no triviales, y parece
poco probable que se satisfagan siempre, aunque el conferenciante no ha
construido ningún contraejemplo. En cambio, parece plausible que la condición
$H^0(X_0,\Theta_{X_0/k})=0$ sea una condición suficiente (aunque en modo alguno
necesaria) para la existencia de un esquema formal de módulos. Nos limitamos a
enunciar aquí un teorema en un caso particularmente sencillo (cuyo análogo en
la teoría de espacios analíticos es bien conocido, cf. Kodaira--Spencer), que se
establece sin dificultad mediante los resultados del número anterior:

\result{TEOREMA}{10} Sea $X_0$ un esquema propio y simple sobre el cuerpo $k$
tal que
\[
  H^0(X_0,\Theta_{X_0/k})=H^2(X_0,\Theta_{X_0/k})=0.
\]
Entonces existe un esquema formal de módulos para $X_0$, correspondiente a un
anillo local regular $\mathcal O$ (es decir, un álgebra de series formales sobre
$k$).

Como ya hemos señalado, no es cierto en general que el esquema formal
$\mathfrak X$ sobre $\mathcal O$ sea algebraizable; pero se sabe que sí lo es
cuando $X_0$ es proyectivo y $H^2(X_0,\Osh_{X_0})=0$ (teorema 4), por ejemplo
cuando $X_0$ tiene dimensión 1. Esto da cierto fundamento a la esperanza de
construir un esquema de módulos sobre los enteros para las curvas de género
dado\ldots{}

Obsérvese también que métodos como los expuestos en este número pueden aplicarse
a la construcción y el estudio de las variedades de Picard, así como a muchas
otras construcciones. Volveremos sobre ello próximamente.

\fsection{8}{Aplicación al grupo fundamental}

Las técnicas expuestas permiten abordar el estudio sistemático del grupo
fundamental siguiendo el modelo de la teoría topológica. Los dos primeros
teoremas enunciados en este número son generalizaciones de resultados de un
trabajo reciente de Lang--Serre.

Sea $X$ un esquema. Un $X$-esquema $X'$ se llama revestimiento no ramificado de
$X$ si
\begin{enumerate}
  \item[\textit{i.}] $X'$ es finito sobre $X$, es decir, está definido por un
    haz coherente de álgebras $\mathcal A=\mathcal A(X')$ sobre $X$;
  \item[\textit{ii.}] $\mathcal A$ es un haz localmente libre sobre $X$;
  \item[\textit{iii.}] Para todo $x\in X$, el cociente
    \[
      \mathcal A_x/\mathfrak m_x\mathcal A_x
      =\mathcal A_x\otimes_{\Osh_x}\kappa(x)
    \]
    es un álgebra separable sobre $\kappa(x)$.
\end{enumerate}

\src{210}
Esta noción de revestimiento no ramificado (debida a Serre y al conferenciante)
posee todas las propiedades elementales que cabe razonablemente esperar, cuya
lista no daremos. Limitémonos a decir que da lugar a una teoría de Galois
modelada sobre la teoría de Galois clásica (y que la contiene; las
demostraciones son más bien más sencillas que las comúnmente aceptadas para
esta última) y sobre la teoría de Galois de los revestimientos topológicos. Más
precisamente, llamemos \emph{punto geométrico} de un esquema $X$ a un morfismo
$a$ del espectro $\xi$ de un cuerpo algebraicamente cerrado $\Omega$ en $X$, es
decir, al dato de una extensión algebraicamente cerrada del cuerpo residual
$\kappa(x)$ de un punto $x=|a|$ de $X$ (llamado \emph{localidad} del punto
geométrico $a$). Si $X'$ es entonces un revestimiento no ramificado de $X$, se
le puede asociar el conjunto $E_a(X')$ de los «puntos geométricos de $X'$ sobre
$a$», es decir, el conjunto de los pares formados por un $x'\in X'$ sobre $x$
y un $\kappa(x)$-homomorfismo $\kappa(x')\longrightarrow\Omega$. Se obtiene así
(para $(X,a)$ fijos) un funtor $F(X,a)$ de la categoría $R(X)$ de los
revestimientos no ramificados $X'$ de $X$ en la categoría de los conjuntos
finitos. Si $X$ es conexo, el par formado por $R(X)$ y $F(X,a)$ posee las
propiedades formales necesarias para ser isomorfo al par análogo definido por
un grupo topológico compacto totalmente discontinuo $\pi$ adecuado (es decir,
un límite proyectivo de grupos finitos): se toma la categoría $C(\pi)$ de los
conjuntos finitos $E$ sobre los cuales $\pi$ actúa continuamente, y el funtor
identidad $F(\pi)(E)=E$ de esta categoría en la categoría de los conjuntos
finitos. Además, el grupo $\pi$ está determinado salvo isomorfismo canónico por
la condición de que $(C(\pi),F(\pi))$ sea isomorfo a un par dado. En este caso,
$\pi$ se llama \emph{grupo fundamental} del esquema conexo $X$ en el punto
geométrico $a$, y se denota por $\pi_1(X,a)$. Si $X$ no es conexo, se lo
reemplaza por la componente conexa de $x=|a|$. Si, por el contrario, $X$ es
conexo, los grupos $\pi_1(X,a)$ correspondientes a dos puntos geométricos
$a',a''$ de $X$ son isomorfos (el isomorfismo está determinado salvo
automorfismos interiores), de modo que, como de costumbre, puede escogerse $a$
de la manera más conveniente, por ejemplo en el punto genérico de $X$, supuesto
irreducible. Naturalmente, $\pi_1(X,a)$ es un funtor covariante del esquema
punteado $(X,a)$. Todo enunciado relativo a la clasificación de los
revestimientos no ramificados\footnote{El texto impreso dice «revestimientos
inseparables»; el párrafo trata de la categoría $R(X)$ de los revestimientos no
ramificados y de su equivalencia con los conjuntos finitos dotados de una acción
continua de $\pi_1(X,a)$.} puede traducirse entonces al lenguaje de la teoría de
grupos siguiendo el diccionario bien conocido (salvo que hay que tener en cuenta
que aquí se trata de grupos topológicos).

\src{211}
Nuestro objetivo es desarrollar el análogo de la sucesión exacta de homotopía
de los espacios fibrados para un morfismo propio $f:X\longrightarrow Y$.
Evidentemente, al no saber qué son los grupos de homotopía superiores, solo
obtendremos resultados necesariamente incompletos. Para poder aplicar los
teoremas fundamentales del n\textsuperscript{o} 3, debemos explicitar primero
algunos lemas elementales relativos a los esquemas sobre cuerpos o anillos de
Artin (¡conforme al procedimiento general!).

\result{LEMA}{1} Sea $(X',a')$ un revestimiento no ramificado punteado asociado
a una representación punteada de $\pi_1(X,a)$ en un conjunto finito $E$
(provisto del punto marcado $e$). Entonces el morfismo canónico
$\pi_1(X',a')\longrightarrow\pi_1(X,a)$ identifica el primer grupo con el
estabilizador de $e$ en $\pi_1(X,a)$ (y es, por tanto, inyectivo).

\result{LEMA}{2} Sea $X$ un esquema algebraico sobre el cuerpo $k$, y sea $k'$
una extensión radicial de $k$. Entonces todo revestimiento no ramificado de
$X\otimes_k k'$ proviene por imagen inversa (es decir, por extensión de
escalares) de un revestimiento no ramificado de $X$, determinado salvo
isomorfismo canónico.

De estos dos lemas resulta, en particular, que para toda extensión algebraica
$K$ de $k$, y para todo punto geométrico $a'$ de $X'=X\otimes_k K$ que se
proyecta sobre el punto geométrico $a$ de $X$, el homomorfismo funtorial
$\pi_1(X',a')\longrightarrow\pi_1(X,a)$ es inyectivo.

\result{LEMA}{3} Sea $X$ un esquema completo sobre un anillo local artiniano
$A$, tal que $H^0(X,\Osh_X)=A$; sea $X'$ un revestimiento no ramificado de $X$,
y sea $A'=H^0(X',\Osh_{X'})$, que es, por tanto, un anillo finito sobre $A$
(que \textit{a priori} puede estar ramificado sobre $A$). Sean $X_0,X'_0$ los
subesquemas reducidos asociados a $X,X'$ (obtenidos dividiendo por los haces de
elementos nilpotentes de $\Osh_X,\Osh_{X'}$), y sea $k$ un subcuerpo de
$A/\mathfrak r(A)$ sobre el cual $A/\mathfrak r(A)$ sea finito (así, $X_0$ es
un esquema algebraico completo sobre $k$, y $X'_0$ es un revestimiento no
ramificado suyo). Sea, por último, $\Omega$ una extensión algebraicamente
cerrada de $k$, y consideremos el revestimiento no ramificado
$X'_0\otimes_k\Omega$ de $X_0\otimes_k\Omega$. Entonces las dos condiciones
siguientes son equivalentes:
\begin{enumerate}
  \item[\textit{i.}] $X'_0\otimes_k\Omega$ está completamente descompuesto
    sobre $X_0\otimes_k\Omega$.
  \item[\textit{ii.}] El morfismo natural
    $X'\longrightarrow X\otimes_A A'$ es un isomorfismo.
\end{enumerate}
Bajo estas condiciones, $A'$ es una extensión no ramificada de $A$. Por último,
si $X'$ es conexo, la condición \textit{(i)} equivale a la condición siguiente,
en apariencia más débil:
\begin{enumerate}
  \item[\textit{i bis.}] $X'_0\otimes_k\Omega$ admite una sección regular
    sobre $X_0\otimes_k\Omega$.
\end{enumerate}
Cuando se satisface la condición \textit{(ii)}, diremos que el revestimiento no
ramificado $X'$ de $X$ es \emph{geométricamente trivial}.

\src{212}
\result{LEMA}{4} Sea $f:X\longrightarrow Y$ un morfismo propio tal que
$f_*(\Osh_X)=\Osh_Y$. Sean $a$ un punto geométrico de $X$ y $b$ su proyección
sobre $Y$. Entonces $\pi_1(X,a)\longrightarrow\pi_1(Y,b)$ es sobreyectivo.

En efecto, hay que demostrar lo siguiente: si un revestimiento no ramificado
$Y'$ de $Y$ (correspondiente a un haz de álgebras localmente libre
$\mathcal A$) es tal que $X\otimes_Y Y'$ es disconexo, entonces $Y'$ también lo
es. En efecto, $f^*(\mathcal A)$ será suma directa de dos haces de anillos no
nulos y, por tanto, también lo será su imagen directa; ahora bien, esta última
no es otra que $\mathcal A\otimes f_*(\Osh_X)=\mathcal A$.

\result{LEMA}{5} Sea $X$ un esquema completo sobre un cuerpo $k$; se supone que
$H^0(X,\Osh_X)$ es un anillo local $A$ y que $A/\mathfrak r(A)$ es radicial
sobre $k$. Sea $\Omega$ una clausura algebraica de $k$, y sea
$\overline X=X\otimes_k\Omega$ (es conexo). Elijamos un punto geométrico
$\overline a$ de $\overline X$ que se proyecte sobre el punto geométrico $a$ de
$X$; entonces se tiene una sucesión exacta
\[
  e\longrightarrow\pi_1(\overline X,\overline a)
  \longrightarrow\pi_1(X,a)
  \longrightarrow\pi_1(k,b)\longrightarrow e
\]
(donde $\pi_1(k,b)$ es el grupo de Galois de $\Omega$ sobre $k$).

La inyectividad del primer homomorfismo se ha visto con los lemas 1 y 2, la
exactitud en el medio resulta del lema 3 y, por último, la sobreyectividad del
último homomorfismo (que es la única que utiliza el hecho de que
$A/\mathfrak r(A)$ sea radicial) resulta del lema 4.

\result{PROPOSICIÓN}{4} Sea $f:X\longrightarrow Y$ un morfismo propio y plano
tal que, para todo $y\in Y$, $H^0(f^{-1}(y),\Osh_{f^{-1}(y)})$ sea un álgebra
separable sobre el cuerpo residual $\kappa(y)$ (lo cual ocurre, por ejemplo, si
$f^{-1}(y)$ es un esquema separable sobre $\kappa(y)$, es decir, reducido y tal
que los cuerpos correspondientes a sus componentes irreducibles sean
extensiones separables de $\kappa(y)$). Entonces el revestimiento $Y'$ de $Y$
asociado a $f_*(\Osh_X)$ es no ramificado.

La demostración es fácil gracias al teorema 2.

\src{213}
Esta proposición, junto con el lema 1, reduce prácticamente el estudio
homotópico de los morfismos propios y planos (con fibras separables) al caso en
que $f_*(\Osh_X)=\Osh_Y$ (pues, utilizando la factorización de Stein, se
reemplaza $Y$ por $Y'$).

\result{OBSERVACIÓN}{} Un morfismo plano de tipo finito cuyas fibras son
esquemas separables (resp. simples) se llama \emph{separable} (resp.
\emph{simple}). Se demuestra que, si $f$ es plano y $f^{-1}(y)$ es separable
(resp. simple), existe un entorno de $f^{-1}(y)$ en el que $f$ es separable
(resp. simple). El mismo resultado es válido para «absolutamente normal»
(«teorema de Bertini»).

Sea $f:X\longrightarrow Y$ un morfismo propio tal que
\begin{equation}
  \tag{i}
  f_*(\Osh_X)=\Osh_Y,
\end{equation}
y sea $X'$ un esquema finito sobre $X$. Sea $Y'$ el revestimiento de $Y$
correspondiente a la factorización de Stein de $X'\longrightarrow Y$
(cf. teorema 5). Sea $y\in Y$; entonces el conjunto de las componentes conexas
de la fibra $F'$ de $X'$ sobre $y$ se identifica con el conjunto de los puntos
$y'$ de $Y'$ sobre $y$ (teorema 5). Consideremos el morfismo evidente
\begin{equation}
  \tag{*}
  X'\longrightarrow X\times_Y Y'
\end{equation}
deducido de los morfismos naturales $X'\longrightarrow X$ y
$X'\longrightarrow Y'$\footnote{El texto impreso omite la prima en el segundo
$X'$; el morfismo hacia $Y'$ es el de la factorización de Stein de
$X'\to Y$.}; será un isomorfismo siempre que $X'$ sea de la forma
$X\times_Y Y''$, donde $Y''$ es un revestimiento no ramificado de $Y$, y
entonces $Y'$ no será sino $Y''$ y (*) será la identidad. Queremos dar
precisamente condiciones bajo las cuales $X'$ sea de la forma que acabamos de
indicar, es decir, que $Y'$ sea no ramificado y (*) un isomorfismo. Para ello,
introduzcamos la fibra $F$ de $X$ en $y$; es un esquema propio sobre
$\kappa(y)$ del cual $F'$ es un revestimiento (no ramificado si $X'$ lo es).
Sea $F'_1$ una componente conexa de $F'$ correspondiente a un punto $y'_1$ de
$Y'$ sobre $y$. Supongamos además
\begin{enumerate}
  \item[\textit{(ii)}] $X'$ es no ramificado sobre $X$ en los puntos de $F'_1$
    (por tanto, $F'_1$ es un revestimiento no ramificado de $F$), y
  \item[\textit{(iii)}] $F'_1$ es un revestimiento geométricamente trivial de
    $F$ (cf. lema 3).
\end{enumerate}

\result{TEOREMA}{11} En estas condiciones, existe un entorno abierto $U'$ de
$y'_1$ en $Y'$ tal que (*) es un isomorfismo sobre $U'$. Además, $Y'$ es no
ramificado en $y'_1$ sobre $Y$ (pero puede ser ramificado en los demás puntos
$y'$ de $Y'$ sobre $y$).

Naturalmente, las condiciones \textit{(ii)} y \textit{(iii)} son también
necesarias para que sea válida la conclusión del teorema. La demostración del
teorema es fácil, con ayuda del lema 3 y del teorema 2.

\src{214}
\result{COROLARIO}{1} Supongamos siempre satisfecha \textit{(i)}. Para que un
revestimiento no ramificado sobre $X$ sea isomorfo a la imagen inversa de un
revestimiento no ramificado $Y'$ de $Y$, es necesario y suficiente que $X'$
induzca sobre cada fibra $f^{-1}(y)$ un revestimiento geométricamente trivial.

Como, según el teorema 11, el conjunto de los puntos de $Y$ para los cuales se
cumple esta condición es abierto, basta verificarla en los puntos $y$ que son
cerrados\ldots{} Señalemos el enunciado siguiente, equivalente al corolario 1:
el núcleo del homomorfismo (sobreyectivo según el lema 4)
$\pi_1(X)\longrightarrow\pi_1(Y)$ es el subgrupo invariante cerrado engendrado
por las imágenes en $\pi_1(X)$ de los $\pi_1(\overline{f^{-1}(y)})$, donde
\[
  \overline{f^{-1}(y)}
  =f^{-1}(y)\otimes_{\kappa(y)}\overline{\kappa(y)},
\]
$\overline{\kappa(y)}$ designa una clausura algebraica de $\kappa(y)$. Se
observará que, al no poder elegir el mismo punto base para todas las fibras,
los homomorfismos $\pi_1(\overline{f^{-1}(y)})\longrightarrow\pi_1(X)$ solo
quedan determinados en cualquier caso (una vez elegido un punto base para $X$
y, por tanto, para $Y$) módulo composición por un automorfismo interior en
$\pi_1(X)$.

\result{COROLARIO}{2} En las condiciones generales del teorema 11, supongamos
además que $Y,X,X'$ son íntegros, y sean $K,L,L'$ sus cuerpos. Entonces existe
una subextensión separable $K'$ de $K$ en $L'$, linealmente disjunta de $L$, tal
que $L'=LK'$ (de donde $L'=L\otimes_K K'$).

(Se aplica la última parte del lema 3 a la fibra geométrica de $X$). El caso de
aplicación más interesante del teorema 11 se obtiene cuando $f$ es un morfismo
separable. Entonces $X'$ es también separable sobre $Y$; por tanto, en virtud
de la proposición 4, $Y'$ es no ramificado sobre $Y$, y el segundo miembro
$X\times_Y Y'$ de (*) es no ramificado sobre $X$. Se concluye fácilmente:

\result{COROLARIO}{3} Supongamos, además de \textit{(i)}, que $f$ es separable.
Sea $X'$ un revestimiento no ramificado y conexo de $X$; para que
$X'$\footnote{El texto impreso dice aquí «$X$»; el sujeto del criterio es el
revestimiento $X'$ introducido en la frase anterior.} sea imagen inversa de un
revestimiento no ramificado $Y'$ de $Y$, es necesario y suficiente que el
revestimiento $\overline F'$ inducido sobre una fibra geométrica
$\overline F=f^{-1}(y)$ admita una sección regular.

Se observará que no era necesario suponer que $\overline F'$ fuera
geométricamente trivial sobre $\overline F$ (lo cual será cierto \textit{a
posteriori}, aunque \textit{a priori} esta condición sea mucho más fuerte). El
corolario 3 es, de hecho, equivalente al enunciado siguiente:

\src{215}
\result{COROLARIO}{4} Sea $f:X\longrightarrow Y$ un morfismo propio y separable
tal que $f_*(\Osh_X)=\Osh_Y$, sea $\overline F$ la fibra geométrica de un punto
$y\in Y$, y elijamos un punto geométrico en $\overline F$, del cual se obtienen,
mediante los morfismos $\overline F\longrightarrow X\longrightarrow Y$, puntos
geométricos en $X,Y$, que se tomarán como puntos base para los grupos
fundamentales de $\overline F,X,Y$. En estas condiciones se tiene la sucesión
exacta
\[
  \boxed{
    \pi_1(\overline F)\longrightarrow\pi_1(X)
    \longrightarrow\pi_1(Y)\longrightarrow e
  }.
\]

De ello se deducen fácilmente los dos enunciados siguientes de Serre--Lang,
liberados aquí de toda hipótesis de normalidad:

\result{COROLARIO}{5} Sean $X,Y$ esquemas conexos sobre un cuerpo $k$;
supongamos que el esquema reducido $X_{\mathrm{réd}}$ es separable sobre $k$
(lo cual es automáticamente cierto si $k$ es perfecto), y que $X$ o $Y$ es
propio sobre $k$.\footnote{Condición añadida en la fe de erratas de 1962.}
Elijamos un punto geométrico $a$ (resp. $b$) en $X$ (resp. $Y$), de donde se
obtienen un punto geométrico $c=(a,b)$ en $X\times_kY$ y un morfismo natural
\[
  \pi_1(X\times_kY,c)\longrightarrow
  \pi_1(X,a)\times\pi_1(Y,b)
\]
(deducido de los morfismos funtoriales de $\pi_1(X\times Y,c)$ en $\pi_1(X,a)$
y $\pi_1(Y,b)$). Este morfismo es inyectivo, e incluso biyectivo si $k$ es
algebraicamente cerrado.

(La sobreyectividad en este último caso es casi trivial). De ello se deduce,
con Serre--Lang:

\result{COROLARIO}{6} Sea $X$ un esquema algebraico conexo y propio sobre un cuerpo
algebraicamente cerrado $k$, y sea $K$ una extensión algebraicamente cerrada
de $k$; entonces los grupos fundamentales de $X$ y $X\otimes_kK$ son los
mismos, es decir, todo revestimiento no ramificado de este último esquema
proviene por extensión de escalares de un revestimiento no ramificado (único
salvo isomorfismo) de $X$.

\result{OBSERVACIONES}{}

1) Utilizando la proposición 4, se ve que la hipótesis
$f_*(\Osh_X)=\Osh_Y$ del corolario 4 no es esencial. En el caso general, en vez
de colocar el grupo unidad $e$ después de $\pi_1(Y)$, hay que continuar con
\[
  \pi_0(\overline F)\longrightarrow\pi_0(X)
  \longrightarrow\pi_0(Y)\longrightarrow e.
\]
como en Topología algebraica.

\src{216}
2) En general, por el momento no se puede decir nada sobre el núcleo de
$\pi_1(\overline F)\longrightarrow\pi_1(X)$, en el que debería intervenir un
$\pi_2(Y)$. Parece, sin embargo, que debería poder demostrarse que
$\pi_1(\overline F)\longrightarrow\pi_1(X)$ es inyectivo si $Y$ es el espectro
de un anillo local $A$, apoyándose en el teorema 12 de más abajo (que afirma que
así es si $A$ es completo).

El teorema 11 solo utilizaba los teoremas 1 y 2. Vamos a aplicar ahora el
teorema 3, apoyándonos en el lema elemental siguiente:

\result{LEMA}{6} Sea $X$ un esquema, y $X_0$ el esquema reducido correspondiente
(es decir, en el que se han anulado los elementos nilpotentes). Entonces todo
revestimiento no ramificado $X'_0$ de $X_0$ está inducido por un revestimiento
no ramificado $X'$ de $X$, determinado salvo isomorfismo canónico.

Este lema, de naturaleza puramente local, desempeña aquí un papel análogo al del
teorema 8 en teoría de módulos. Combinándolo con el teorema de existencia
(teorema 3), se deduce:

\result{TEOREMA}{12} Sea $A$ un anillo local noetheriano y completo de cuerpo
residual $k$. Sea $X$ un esquema propio sobre $A$. Entonces todo revestimiento
no ramificado $X'_0$ de $X_0=X\otimes_A k$ está inducido por un revestimiento no
ramificado $X'$ de $X$, único salvo isomorfismo.

En otros términos:

\result{COROLARIO}{1} Elijamos un punto geométrico en $X_0$ como punto base para
los grupos fundamentales de $X_0$ y $X$. Entonces el homomorfismo canónico
$\pi_1(X_0)\longrightarrow\pi_1(X)$ es un isomorfismo.

Apliquemos ahora a $X_0$ el lema 5 (suponiendo, para simplificar, que
$H^0(X_0,\Osh_{X_0})=k$), y observemos que, por ser $A$ completo, sus
extensiones no ramificadas corresponden a las extensiones no ramificadas de su
cuerpo residual, es decir, se tiene $\pi_1(Y)=\pi_1(k)$ (donde
$Y=\operatorname{Spec}(A)$). Se obtiene la sucesión exacta:
\[
  e\longrightarrow\pi_1(\overline X_0)
  \longrightarrow\pi_1(X)\longrightarrow\pi_1(Y)\longrightarrow e.
\]

\result{COROLARIO}{2} Sea $f:X\longrightarrow Y$ un morfismo propio y plano,
sean $y_1$ un punto de $Y$ e $y_0$ una especialización de $y_1$, y consideremos
las fibras «geométricas» correspondientes $\overline X_1$ y $\overline X_0$;
se supone que $\overline X_0$ es separable y conexa (lo que implica que
$\overline X_1$ satisface las mismas condiciones). Entonces se puede hallar un
homomorfismo de grupos
$\pi_1(\overline X_1)\longrightarrow\pi_1(\overline X_0)$, definido salvo un
automorfismo interior, y este homomorfismo es sobreyectivo.

\src{217}
Podría esperarse que este homomorfismo fuera siempre biyectivo; por desgracia,
en general no es así si $\kappa(y_0)$ tiene característica $>0$. Obtendremos,
sin embargo, más abajo una cota para el núcleo de este homomorfismo (al menos si
$\overline X_0$ es simple), de la que se sigue que, si $\kappa(y_0)$ tiene
característica $0$, el homomorfismo anterior es biyectivo (resultado que también
podría demostrarse por vía trascendente). En todo caso, obtenemos ya una cota
para el $\pi_1$ de una fibra especial por medio del de una fibra genérica.
Utilizando, por ejemplo, el hecho de que una curva algebraica en característica
$p$ se puede levantar a una curva en característica $0$ (teorema 9, corolario
4), se obtiene por vía trascendente:

\result{COROLARIO}{3} Sea $X_0$ el esquema de una curva completa simple sobre
un cuerpo algebraicamente cerrado de característica cualquiera, y sea $g$ el
género de $X_0$; entonces $\pi_1(X_0)$ admite $2g$ generadores topológicos,
ligados por la relación bien conocida.

De ello se deduce, mediante una técnica bien conocida de secciones
hiperplanas:

\result{COROLARIO}{4} Sea $X$ un esquema proyectivo simple sobre un cuerpo
algebraicamente cerrado de característica cualquiera; entonces $\pi_1(X)$
admite un número finito de generadores topológicos.

Busquemos el núcleo del homomorfismo
$\pi_1(\overline X_1)\longrightarrow\pi_1(\overline X_0)$. Para ello se puede
suponer que $Y$ es el espectro de un anillo de valoración discreta completo
$V=\Osh_y$ ($y=y_0$). La cuestión equivale a la siguiente: dado un revestimiento
no ramificado $\overline X'_1$ de $\overline X_1$ (que, si se quiere, puede
suponerse galoisiano), determinar en qué condiciones proviene de un
revestimiento no ramificado de $X_0$. \textit{A priori}, el revestimiento dado
proviene, por extensión de escalares, de un revestimiento no ramificado $X'_1$
de $X_1\otimes_K K'$, donde $K'$ es una subextensión finita de la clausura
algebraica $\overline K$ del cuerpo de fracciones $K$ de $V$; si
$\overline X'_1$ fuera galoisiano de grupo $G$, se puede elegir $X'_1$
galoisiano de grupo $G$. Se tiene entonces: para que
$X'_1\otimes_{K'}\overline K=\overline X'_1$ provenga de un revestimiento no
ramificado de $X_0$, es necesario y suficiente que se pueda hallar una extensión
finita $K''$ de $K'$ en $\overline K$ tal que
$X''_1=X'_1\otimes_{K'}K''$ sea de la forma $X''\otimes_{V''}K''$, donde $V''$
es la clausura normal de $V$ en $K''$ y $X''$ es un revestimiento no ramificado
de $X\otimes_VV''$. Supongamos, por ejemplo, que $X_0$ es absolutamente normal;
de ello resulta que $X\otimes_VV''$ es normal (pues es plano sobre $V''$ y tiene
fibra especial normal), y su cuerpo de funciones es idéntico al $K''(X_1)$ de
\[
  X_1\otimes_KK''
  =(X\otimes_VK)\otimes_KK''
  =X\otimes_VK''
  =(X\otimes_VV'')\otimes_{V''}K''.
\]
Sea $L''=K''(X'_1)$ el cuerpo de funciones de $X'_1\otimes_{K'}K''$; es una
extensión finita separable de $K''(X_1)$, y la condición anterior significa
también que $L''$ es una extensión no ramificada del cuerpo de funciones de
$X\otimes_VV''$ (es decir, la normalización de $X\otimes_VV''$ en $L''$ es no
ramificada sobre $X\otimes_VV''$). Basta, por lo demás, verificar que $L''$ es
no ramificado en los puntos de la fibra especial de $X\otimes_VV''$ (pues es no
ramificado sobre la fibra genérica $X_1\otimes_KK''$). Si ahora $X_0$ es simple,
del «teorema de pureza» de Nagata--Zariski resulta incluso que basta verificar
que $L''$ es no ramificado sobre el anillo local $O''$ del punto genérico de la
fibra especial de $X\otimes_VV''$, que es un anillo de valoración discreta,
normalización en $K''(X_1)$ del anillo local $O\subset K(X_1)$ del punto
genérico de la fibra especial de $X$. Por tanto, se ha obtenido:

\src{218}
\result{COROLARIO}{5} En las condiciones y con las notaciones anteriores, para
que el revestimiento no ramificado
$X'_1\otimes_{K'}\overline K$ de
$\overline X_1=X_1\otimes_K\overline K$ provenga de un revestimiento no
ramificado de $\overline X_0$, es necesario y suficiente que exista una
subextensión finita $K''$ de $\overline K/K'$ tal que $K''(X'_1)$ sea no
ramificado sobre el anillo de valoración discreta $O''\subset K''(X_1)$.

Observemos ahora que $O''$ es la normalización en $K''(X_1)$ del anillo de
valoración discreta $O'\subset K'(X_1)$ (normalización de $O$ en $K'(X_1)$), y
que $O'$ contiene la normalización $V'$ de $V$ en $K'$; un uniformizante $u$ de
$V'$ es también un uniformizante de $O'$. Supongamos ahora que $X'_1$ es
galoisiano de grupo de Galois $G$, de orden $n$ primo con la característica $p$
de $\kappa(y_0)$ (que es también la característica del cuerpo residual de
$O'$). Por tanto, $K'(X'_1)$ está «moderadamente ramificado» sobre $O'$, de lo
que resulta fácilmente («lema de Abhyankar») que, si se le adjunta una raíz
$n$-ésima $v$ de un uniformizante de $O'$, se vuelve no ramificado sobre la
normalización de $O'$ en $K'(X_1)(v)$. Ahora bien, se puede tomar como $v$ una
raíz $n$-ésima de un uniformizante de $V'$, lo que prueba que se satisface la
condición del corolario 5. (Esta posibilidad de utilizar el lema de Abhyankar y
el teorema de pureza me la ha vendido Serre). Para expresar el resultado
obtenido, introduzcamos, para todo grupo compacto totalmente discontinuo $\pi$,
el grupo cociente $\widetilde\pi$ de $\pi$ por el subgrupo cerrado engendrado
por sus $p$-subgrupos de Sylow, es decir, el límite proyectivo de los grupos
cocientes discretos de $\pi$ que son de orden primo con $p$. Con esta notación,
se obtiene:

\result{TEOREMA}{13} Sea $f:X\longrightarrow Y$ un morfismo propio y plano,
sean $y_1$ un punto de $Y$ e $y_0$ una especialización de $y_1$, y supongamos
$\overline X_0$ conexo y simple. Entonces el homomorfismo
$\widetilde\pi_1(\overline X_1)\longrightarrow
\widetilde\pi_1(\overline X_0)$ deducido del homomorfismo sobreyectivo del
teorema 12, corolario 2, es un isomorfismo.

En otros términos:

\src{219}
\result{COROLARIO}{1} La clasificación de los revestimientos galoisianos no
ramificados, de grupo de Galois de orden primo con la característica $p$ de
$\kappa(y_0)$, es la misma para $\overline X_0$ y para $\overline X_1$.

En particular, si $\kappa(y_0)$ tiene característica cero, se obtiene por vía
algebraica que $\pi_1(\overline X_1)\longrightarrow\pi_1(\overline X_0)$ es
biyectivo.

Señalemos, por último, que las técnicas utilizadas proporcionan también el
resultado siguiente, más general que el teorema 13:

\result{TEOREMA}{14} Sea $f:X\longrightarrow Y$ un morfismo propio y simple,
sea $D$ un subesquema cerrado de $X$, simple sobre $Y$, de codimensión 1 en
todos sus puntos. Para una fibra $Z=f^{-1}(z)$ de $f$, sea $Z'=Z-Z\cap D$ y sea
$\pi_1^t(Z')$ el cociente del grupo fundamental $\pi_1(Z')$ que clasifica los
revestimientos no ramificados de $Z'$ que están «moderadamente ramificados»
sobre $\overline{Z\cap D}$. Sean $y_0,y_1$ como en el teorema 13. Entonces se
tiene un homomorfismo sobreyectivo (definido salvo un automorfismo interior)
\[
  \pi_1^t(\overline X'_1)\longrightarrow\pi_1^t(\overline X'_0),
\]
y el homomorfismo correspondiente
\[
  \widetilde\pi_1^t(\overline X'_1)
  \longrightarrow\widetilde\pi_1^t(\overline X'_0)
\]
es un isomorfismo.

De ello se deducen las variantes correspondientes de los corolarios del
teorema 13 y del corolario 4 del teorema 12. Del mismo modo, utilizando el
teorema 9, corolario 3, se obtiene por vía trascendente:

\result{COROLARIO}{} Sea $X_0$ el esquema de una curva completa simple sobre un
cuerpo algebraicamente cerrado de característica cualquiera, y sea
$S=(s_i)_{1\le i\le n}$ una parte finita con $n$ elementos de $X_0$. Entonces
$\pi_1^t(X_0-S)$ admite $2g+n$ generadores topológicos
$x_i,y_i$ ($1\le i\le g$) y $\sigma_j$ ($1\le j\le n$), ligados por la relación
\[
  \left(\prod_i x_i y_i x_i^{-1}y_i^{-1}\right)
  \sigma_1\cdots\sigma_n=1,
\]
donde los $\sigma_j$ son generadores de los grupos de inercia correspondientes
a los $s_j$. Para todo grupo finito $G$ de orden primo con la característica,
engendrado por elementos $x_i,y_i,\sigma_j$ que satisfacen la relación anterior,
existe un revestimiento galoisiano no ramificado de $X_0-S$, de grupo $G$, cuyos
grupos de inercia en los puntos $s_j$ están engendrados por los $\sigma_j$.

Cuando $X_0$ es de género 0 y $n=3$, se tiene una solución del «problema de los
tres puntos», al menos para los revestimientos galoisianos de orden primo con
la característica. (Aquí, por lo demás, el teorema 9 es inútil, y parece, por
otra parte, que el corolario anterior puede deducirse del caso particular
considerado en el problema de los tres puntos).

\src{220}
\result{OBSERVACIONES}{}

1) Un estudio más completo, en el que sin duda intervendrían revestimientos
galoisianos generalizados de $X,X_0,X_1$ (con grupo de Galois eventualmente
infinitesimal), debería permitir recuperar el núcleo del teorema 12, corolario
2. En cambio, el estudio de los revestimientos que admiten ramificación no
«moderada» parece tener que ser mucho más difícil.

2) El lema 6, junto con un resultado de Grauert relativo a la completación
formal de un esquema proyectivo no singular a lo largo de una sección
hiperplana (o con el teorema, no demostrado por el momento, mencionado en la
observación 2 después del teorema 11), permite también demostrar en geometría
algebraica «abstracta» el teorema clásico de Lefschetz sobre el grupo
fundamental.

\begin{center}
  \textbf{BIBLIOGRAFÍA}
\end{center}

\begin{description}
  \item[{[1]}] Dieudonné (J.) y Grothendieck (A.). --- \emph{Éléments de
    géométrie algébrique}, Publications mathématiques de l'Institut des Hautes
    Études Scientifiques (de próxima aparición).
  \item[{[2]}] Grothendieck (Alexander). --- \emph{The cohomology theory of
    abstract algebraic varieties}. International Congress of Mathematicians
    [1958. Edinburgh] (de próxima aparición).
  \item[{[3]}] Serre (Jean-Pierre). --- \emph{Faisceaux algébriques
    cohérents}, Annals of Math., Series 2, t. 61, 1955, p. 197--278.
  \item[{[4]}] Serre (Jean-Pierre). --- \emph{Géométrie algébrique et
    géométrie analytique}, Ann. Institut Fourier Grenoble, t. 6, 1955--1956,
    p. 1--42.
\end{description}

\def\FgaUnitKey{e182}\clearpage
\chapter*{Fe de erratas de la exposición n\textsuperscript{o} 182}
\addcontentsline{toc}{chapter}{Fe de erratas de la exposición n\textsuperscript{o} 182}
\setcounter{footnote}{1}
\src{298}

\noindent\textbf{2. Geometría formal y geometría algebraica.}

\begin{center}
[Seminario Bourbaki, t. 11, 1958/59, n\textsuperscript{o} 182, 28 p.]
\end{center}

\noindent\textbf{Página 8, teorema 6.} --- La hipótesis «$f$
cuasiproyectivo» puede sustituirse por la hipótesis más débil «$f$
separado», en virtud del resultado siguiente (cf. SGA VIII, 6.2): todo morfismo
$f:X\longrightarrow Y$ que es cuasifinito y separado es cuasiproyectivo.

\src{299}
\noindent\textbf{Página 12, y página 14, observación 1.} --- Señalemos que J.-P.
Serre ha construido una variedad proyectiva no singular, de dimensión 3, sobre
un cuerpo algebraicamente cerrado $k$, de característica $p>0$, que no proviene
por reducción de un esquema propio sobre un anillo local íntegro de cuerpo
residual $k$ y con cuerpo de fracciones de característica
cero.\footnote{Serre (Jean-Pierre). --- Exemples de variétés projectives en
caractéristique $p$ non relevables en caractéristique zéro, \emph{Proc. Nat.
Acad. Sc. U.S.A.}, t. 47, 1961, p. 108--109.} Mumford habría hallado un ejemplo
análogo, con una superficie proyectiva no singular.

\noindent\textbf{Página 15, observaciones 2 y 3.} --- Actualmente soy menos
optimista respecto de los resultados aquí conjeturados. Sin embargo, el
problema relativo a las variaciones de estructura del espacio proyectivo,
señalado al final de la observación 3, ha sido resuelto afirmativamente por
Hironaka, y el problema análogo para las variedades abelianas ha sido resuelto
por Koizumi.

\noindent\textbf{Página 16, línea 3.} --- Señalemos que Mumford ha construido
recientemente los esquemas de módulos para las curvas de género $g$ (cf. SHMI).
El teorema 10 prueba, por lo demás, que los «esquemas de Jacobi de nivel $n$»
de la teoría de Módulos son no singulares (e incluso simples sobre $\mathbf Z$).

\noindent\textbf{Página 23, línea 12, primera línea del corolario
5.} --- Añadir: «siendo $X$ o $Y$ propio sobre $k$».

La sustancia de los §§ 1 a 5 está contenida en la parte publicada de EGA III, y
la de los §§ 6 y 7 está contenida en SGA III. Para el estudio del grupo
fundamental, véanse SGA V, IX, X, XI, así como SGA 1962 (exposiciones X, XII y
XIII), para los teoremas de tipo Lefschetz y numerosas cuestiones abiertas.
Solo la teoría de los revestimientos moderadamente ramificados (cf. teorema 14)
no ha sido todavía objeto de una redacción formal. El corolario del teorema 14,
que da la determinación completa de los revestimientos galoisianos de orden
primo con la característica de una curva algebraica sobre un cuerpo
algebraicamente cerrado, ha sido utilizado de manera esencial en tres ocasiones:

1\textsuperscript{o} en la demostración por Igusa de la desigualdad de Picard
para las superficies proyectivas no singulares en característica cualquiera;

2\textsuperscript{o} en el estudio (desarrollado independientemente por Ogg y
Šafarevič) del grupo de los fibrados principales homogéneos bajo una variedad
abeliana definida sobre un cuerpo de funciones de una variable, en
característica cualquiera;

3\textsuperscript{o} en la demostración reciente, por Artin, de ciertos
teoremas clave sobre la «cohomología de Weil» de las variedades algebraicas.

\def\FgaUnitKey{190}\unit{Exposé 190 -- Técnica de descenso I}
\sid{190}
\src{299}
\seminar{12\textsuperscript{o} año, 1959/60, n\textsuperscript{o} 190\\Diciembre de 1959}
\paperhead{TÉCNICA DE DESCENSO Y TEOREMAS DE EXISTENCIA EN GEOMETRÍA ALGEBRAICA\\
I. GENERALIDADES. DESCENSO POR MORFISMOS FIELMENTE PLANOS}%
  {por Alexander GROTHENDIECK}

Desde un punto de vista técnico, el presente exposé y los que le seguirán
pueden considerarse variaciones sobre el célebre «teorema 90» de Hilbert.
Parece que A. Weil introdujo el método de descenso en geometría algebraica con
el nombre de «descenso del cuerpo base». Weil se restringe, por lo demás, al
caso de extensiones finitas separables de cuerpos. P. Cartier trató después el
caso de extensiones radiciales de altura 1. Por falta del lenguaje de los
esquemas y, más concretamente, por falta de elementos nilpotentes en los
anillos que se permitía considerar, Cartier no había podido formular la
identidad esencial entre estos dos casos.

En la actualidad, parece que la técnica general de descenso que se expondrá
(unida, cuando proceda, a los teoremas fundamentales de la «geometría formal»,
cf. [3]) está en la base de la mayoría de los teoremas de existencia en
geometría algebraica.\footnote{La fe de erratas de 1962 considera excesiva esta
formulación: es válida en gran medida para las técnicas no proyectivas de los
dos primeros exposés de la serie TDTE I a VI, pero no para las técnicas
proyectivas (TDTE IV, V y VI).} Conviene señalar, además, que esta técnica puede
sin duda trasladarse a la «geometría analítica», y cabe esperar que en un
futuro próximo los especialistas sepan demostrar los análogos «analíticos» de
los teoremas de existencia en geometría formal que se darán en el exposé II.
En cualquier caso, los trabajos recientes de Kodaira--Spencer, cuyos métodos
parecen inadecuados para definir y estudiar «variedades de módulos» en un
entorno de sus puntos singulares, muestran con bastante claridad la necesidad
de métodos más próximos a la teoría de esquemas (que deberán completar, por
supuesto, las técnicas trascendentes).

En el presente exposé I trataremos el caso de descenso más elemental (indicado
en el título). Sin embargo, las aplicaciones de los teoremas 1, 2 y 3 que
siguen son ya muy numerosas. Nos hemos limitado a dar algunas a modo de
ejemplo, sin intentar conferirles aquí la máxima generalidad posible.

Utilizaremos libremente el lenguaje de los esquemas, para el cual remitimos al
exposé ya citado, así como a [2]. Señalemos, no obstante, expresamente que los
preesquemas considerados en el presente exposé no son necesariamente
noetherianos, y los

\src{300}
morfismos no son necesariamente de tipo finito. Así, si $A$ es un anillo local
noetheriano, con completado $\widehat A$, habrá que considerar el anillo no
noetheriano $\widehat A\otimes_A\widehat A$, y los morfismos de esquemas afines
correspondientes a las inclusiones entre los anillos considerados.

\fsection{A}{Preliminares sobre las categorías}

\subsection*{1. Categorías fibradas, datos de descenso, morfismos de
$\mathcal F$-descenso}
\addcontentsline{toc}{subsection}{1. Categorías fibradas y datos de descenso}

\result{a. DEFINICIÓN}{1.1} Una \emph{categoría fibrada} $\mathcal F$ de base
$\mathcal C$ consta de una categoría $\mathcal C$; para cada
$X\in\mathcal C$, de una categoría $\mathcal F_X$; para cada
$\mathcal C$-morfismo $f:X\longrightarrow Y$, de un funtor
$f^*:\mathcal F_Y\longrightarrow\mathcal F_X$, denotado también por
\[
  f^*(\xi)=\xi\times_Y X
\]
para $\xi\in\mathcal F_Y$ (considerando $X$ como un «objeto de
$\mathcal C$ sobre $Y$», es decir, provisto del morfismo $f$); y, por último,
para dos morfismos componibles
$X\xrightarrow{f}Y\xrightarrow{g}Z$, de un isomorfismo de funtores
\[
  c_{f,g}:(gf)^*\longrightarrow f^*g^*,
\]
sujetos estos datos a las condiciones siguientes:

\begin{enumerate}
\item[(i)] $\operatorname{id}_X^*=\operatorname{id}$;
\item[(ii)] $c_{f,g}$ es el isomorfismo identidad si $f$ o $g$ es un
isomorfismo identidad;
\item[(iii)] dados tres morfismos componibles
$X\xrightarrow{f}Y\xrightarrow{g}Z\xrightarrow{h}T$, el siguiente diagrama de
isomorfismos, formado mediante isomorfismos de la forma $c_{u,v}$, es
conmutativo:
\[
\begin{array}{ccc}
  (h(gf))^*=((hg)f)^* & \longrightarrow & f^*(hg)^* \\
  \big\downarrow && \big\downarrow \\
  (gf)^*h^* & \longrightarrow & (f^*g^*)h^*=f^*(g^*h^*).
\end{array}
\]
\end{enumerate}

\result{EJEMPLO}{1} Sea $\mathcal C$ una categoría en la que existen los
productos fibrados. Se define entonces una categoría fibrada $\mathcal F$ de
base $\mathcal C$ tomando $\mathcal F_X=$ categoría de los objetos de
$\mathcal C$ sobre $X$; el funtor
$f^*:\mathcal F_Y\longrightarrow\mathcal F_X$ correspondiente a un morfismo
$f:X\longrightarrow Y$ se define mediante el producto fibrado
$Z\longmapsto Z\times_YX$.

\result{EJEMPLO}{2} Sea $\mathcal C$ la categoría de los preesquemas y, para
$X\in\mathcal C$, sea $\mathcal F_X$ la categoría de los haces
cuasi-coherentes de módulos sobre $X$. Si $f:X\longrightarrow Y$ es un morfismo
de preesquemas, $f^*:\mathcal F_Y\longrightarrow\mathcal F_X$ es el funtor
imagen inversa

\src{301}
de haces de módulos. Se obtiene así una categoría fibrada de base
$\mathcal C$.

\result{b. DEFINICIÓN}{1.2} Un diagrama
\[
  E\xrightarrow{u}E'
  \overset{v_1}{\underset{v_2}{\rightrightarrows}}E''
\]
de aplicaciones de conjuntos se llama \emph{exacto} si $u$ es una biyección de
$E$ sobre la parte de $E'$ formada por los $x'\in E'$ tales que
$v_1(x')=v_2(x')$.

\result{DEFINICIÓN}{1.3} Sea $\mathcal F$ una categoría fibrada de base
$\mathcal C$; consideremos un diagrama
\[
  \tag{+}
  S\xleftarrow{\alpha}S'
  \overset{\beta_1}{\underset{\beta_2}{\rightrightarrows}}S''
\]
de morfismos en $\mathcal C$ tales que
$\alpha\beta_1=\alpha\beta_2$; este diagrama se llama
$\mathcal F$-\emph{exacto} si, para cada par $\xi,\eta$ de elementos de
$\mathcal F_S$, el siguiente diagrama de aplicaciones de conjuntos
\[
  \tag{+'}
  \operatorname{Hom}(\xi,\eta)
  \xrightarrow{\alpha^*}
  \operatorname{Hom}(\alpha^*\xi,\alpha^*\eta)
  \overset{\beta_1^*}{\underset{\beta_2^*}{\rightrightarrows}}
  \operatorname{Hom}(\gamma^*\xi,\gamma^*\eta),
\]
donde $\gamma=\alpha\beta_1=\alpha\beta_2$, es exacto. En este último diagrama
hemos identificado $\beta_i^*\alpha^*$ con
$(\alpha\beta_i)^*=\gamma^*$ mediante $c_{\beta_i,\alpha}$, para simplificar.

\result{DEFINICIÓN}{1.4} Sea $\mathcal F$ una categoría fibrada de base
$\mathcal C$; consideremos dos morfismos $\beta_1,\beta_2$ de $S''$ en $S'$,
en $\mathcal C$. Sea
$\xi'\in\mathcal F_{S'}$. Se llama \emph{dato de pegado} sobre $\xi'$
(relativo al par $(\beta_1,\beta_2)$) a un isomorfismo de
$\beta_1^*(\xi')$ sobre $\beta_2^*(\xi')$. Si $\xi',\eta'\in\mathcal F_{S'}$
están provistos cada uno de un dato de pegado, un morfismo
$u:\xi'\longrightarrow\eta'$ en $\mathcal F_{S'}$ se dice
\emph{compatible con los datos de pegado} si el diagrama siguiente es
conmutativo:
\[
\begin{array}{ccc}
  \beta_1^*(\xi') & \longrightarrow & \beta_2^*(\xi') \\
  \big\downarrow && \big\downarrow \\
  \beta_1^*(\eta') & \longrightarrow & \beta_2^*(\eta').
\end{array}
\]
De este modo, los objetos de $\mathcal F_{S'}$ provistos de un dato de pegado
(relativo a $\beta_1,\beta_2$) forman una categoría. Si
$\alpha:S'\longrightarrow S$ es un morfismo tal que
$\alpha\beta_1=\alpha\beta_2$, entonces, para todo $\xi\in\mathcal F_S$, el
objeto $\xi'=\alpha^*(\xi)$

\src{302}
de $\mathcal F_{S'}$ está provisto de un dato de pegado canónico, puesto que
\[
  \beta_i^*\alpha^*(\xi)\simeq(\alpha\beta_i)^*(\xi)=\gamma^*(\xi),
\]
tomando de nuevo $\gamma=\alpha\beta_1=\alpha\beta_2$; además, si
$u:\xi\longrightarrow\eta$ es un morfismo en $\mathcal F_S$, entonces
\[
  \alpha^*(u):\alpha^*(\xi)\longrightarrow\alpha^*(\eta)
\]
es un morfismo en $\mathcal F_{S'}$ compatible con los datos de pegado
canónicos. Se obtiene así un funtor canónico de la categoría $\mathcal F_S$ a
la categoría de los objetos de $\mathcal F_{S'}$ provistos de un dato de pegado
relativo al par $(\beta_1,\beta_2)$. Dicho esto, la definición 1.3 también
puede expresarse diciendo que el diagrama (+) de dicha definición es
$\mathcal F$-exacto si el funtor anterior es «plenamente fiel», es decir, si
define una equivalencia de la categoría $\mathcal F_S$ con una subcategoría de
la categoría de los objetos de $\mathcal F_{S'}$ provistos de un dato de pegado
relativo a $(\beta_1,\beta_2)$.

\result{DEFINICIÓN}{1.5} Se dice que un dato de pegado sobre
$\xi'\in\mathcal F_{S'}$ es \emph{efectivo} (respecto de $\alpha$) si
$\xi'$, provisto de este dato, es isomorfo a un $\alpha^*(\xi)$, con
$\xi\in\mathcal F_S$.

Cuando el diagrama (+) es $\mathcal F$-exacto, el objeto $\xi$ de la definición
anterior queda, por tanto, determinado salvo isomorfismo único, y la categoría
$\mathcal F_S$ es equivalente a la categoría de los objetos de
$\mathcal F_{S'}$ provistos de un dato de pegado efectivo.

\medskip\noindent
c. El caso más importante es aquel en que
\[
  S''=S'\times_S S',
\]
siendo las $\beta_i$ las dos proyecciones $p_1$ y $p_2$ de
$S'\times_S S'$ sobre sus dos factores (suponemos ahora que en
$\mathcal C$ existen los productos fibrados). Se tienen entonces dos
condiciones necesarias inmediatas para que sea efectivo un dato de pegado
$\varphi:p_1^*(\xi')\longrightarrow p_2^*(\xi')$ sobre un
$\xi'\in\mathcal F_{S'}$. Se tiene
\[
  \tag{i}
  \Delta^*(\varphi)=\operatorname{id}_{\xi'},
\]
donde $\Delta:S'\longrightarrow S'\times_S S'$ designa el morfismo diagonal,
y donde se ha identificado $\Delta^*p_i^*(\xi')$ con
$(p_i\Delta)^*(\xi')=\xi'$. Se tiene
\[
  \tag{ii}
  p_{31}^*(\varphi)=p_{32}^*(\varphi)\,p_{21}^*(\varphi),
\]

\src{303}
donde $p_{ij}$ designa la proyección canónica de
$S'\times_S S'\times_S S'$ sobre el producto parcial de sus factores
$i$-ésimo y $j$-ésimo.

\result{DEFINICIÓN}{1.6} Se llama \emph{dato de descenso} sobre
$\xi'\in\mathcal F_{S'}$, relativo al morfismo
$\alpha:S'\longrightarrow S$, a un dato de pegado sobre $\xi'$ relativo al par
$(p_1,p_2)$ de proyecciones canónicas
$S'\times_S S'\longrightarrow S'$ que satisface las condiciones (i) y (ii)
anteriores.

\result{DEFINICIÓN}{1.7} Un morfismo $\alpha:S'\longrightarrow S$ se llama
\emph{morfismo de $\mathcal F$-descenso} si el diagrama de morfismos
\[
  S\xleftarrow{\alpha}S'
  \overset{p_1}{\underset{p_2}{\rightrightarrows}}S'\times_S S'
\]
es $\mathcal F$-exacto (definición 1.3); se dice que $\alpha$ es un
\emph{morfismo de $\mathcal F$-descenso estricto} si, además, todo dato de
descenso (definición 1.6) sobre un objeto de $\mathcal F_{S'}$ es efectivo.

Esta última condición también puede enunciarse, de modo más gráfico, así:
«Da lo mismo» dar un objeto de $\mathcal F_S$ o un objeto de
$\mathcal F_{S'}$ provisto de un dato de descenso.

Obsérvese que si un morfismo $\alpha:S'\longrightarrow S$ admite una «sección»
$s:S\longrightarrow S'$ (es decir, un morfismo $s$ tal que
$\alpha s=\operatorname{id}_S$), entonces es un morfismo de
$\mathcal F$-descenso estricto.\footnote{La fe de erratas de 1962 precisa que
no es necesario suponer previamente que $\alpha$ sea un morfismo de
$\mathcal F$-descenso.} En efecto, si $\xi'\in\mathcal F_{S'}$ está provisto de
un dato de descenso relativo a $\alpha$, «proviene» de
$\xi=s^*(\xi')$.

\medskip\noindent
d. Las nociones anteriores pueden presentarse de manera más intuitiva
introduciendo, para un objeto parámetro $T$ de $\mathcal C$ sobre $S$, el
conjunto
\[
  \operatorname{Hom}_S(T,S')=S'(T),
\]
cuyos elementos se denotarán por $t,t'$, etc. Dado un objeto
$\xi'\in\mathcal F_{S'}$, a cada $t\in S'(T)$ le corresponde entonces un objeto
$t^*(\xi')$ de $\mathcal F_T$, que también se denotará por $\xi'_t$. Un dato de
pegado sobre $\xi'$ (relativo a $p_1,p_2$) queda definido entonces dando, para
todo $T$ sobre $S$ y todo par de puntos $t,t'\in S'(T)$, un isomorfismo
\[
  \varphi_{t',t}:\xi'_t\longrightarrow\xi'_{t'}
\]
(que satisface condiciones evidentes de funtorialidad al variar $T$). Las
condiciones (i) y (ii) de 1.c se escriben entonces
\[
  \tag{i bis}
  \varphi_{t,t}=\operatorname{id}_{\xi'_t}
  \qquad\text{para todo $T$ y todo $t\in S'(T)$},
\]

\src{304}
\[
  \tag{ii bis}
  \varphi_{t,t''}=\varphi_{t,t'}\varphi_{t',t''}
  \qquad\text{para todo $T$ y todos $t,t',t''\in S'(T)$}.
\]
Se ve, además, que (ii bis), tomando $t=t'=t''$, implica
$\varphi_{t,t}=\varphi_{t,t}^2$; puesto que $\varphi_{t,t}$ es un isomorfismo
por hipótesis, de aquí se obtiene la relación (i bis), que es de hecho una
consecuencia de (ii bis) (por tanto, (i) es consecuencia de (ii)). Pero si ya
no se supone a priori que los $\varphi_{t,t'}$ sean isomorfismos (es decir, que
$\varphi:p_1^*(\xi')\longrightarrow p_2^*(\xi')$ sea un isomorfismo), entonces
(ii bis) no implica necesariamente (i bis); sin embargo, la conjunción de
(ii bis) y (i bis) implica que los $\varphi_{t,t'}$ son isomorfismos (pues se
tendrá
$\varphi_{t,t'}\varphi_{t',t}=\varphi_{t,t}=\operatorname{id}_{\xi'_t}$).

\subsection*{2. Diagramas exactos y epimorfismos estrictos, morfismos de
descenso. Ejemplos}
\addcontentsline{toc}{subsection}{2. Diagramas exactos y epimorfismos estrictos}

\result{a. DEFINICIÓN}{2.1} Sea $\mathcal C$ una categoría. Un diagrama
\[
  T\xrightarrow{\alpha}T'
  \overset{\beta_1}{\underset{\beta_2}{\rightrightarrows}}T''
\]
de morfismos en $\mathcal C$ se llama \emph{exacto} si, para todo
$Z\in\mathcal C$, el diagrama correspondiente de aplicaciones de conjuntos
\[
  \operatorname{Hom}(Z,T)\longrightarrow
  \operatorname{Hom}(Z,T')
  \rightrightarrows\operatorname{Hom}(Z,T'')
\]
es exacto (definición 1.2). Se dice entonces que $(T,\alpha)$ (o, por abuso de
lenguaje, que $T$) es un \emph{núcleo} del par de morfismos
$(\beta_1,\beta_2)$.

Evidentemente, este núcleo está determinado salvo isomorfismo único. Si
$\mathcal C$ es la categoría de los conjuntos, la definición anterior es
compatible con la definición 1.2. Se define de manera dual la exactitud de un
diagrama
\[
  S\xleftarrow{\alpha}S'
  \overset{\beta_1}{\underset{\beta_2}{\rightrightarrows}}S''
\]
de morfismos en $\mathcal C$; se dice entonces que $(S,\alpha)$ es un
\emph{conúcleo} del par de morfismos $(\beta_1,\beta_2)$.

\result{DEFINICIÓN}{2.2} Un morfismo $\alpha:S'\longrightarrow S$ se llama
\emph{epimorfismo estricto} si es un epimorfismo y si, para todo morfismo
$u:S'\longrightarrow Z$, la siguiente condición necesaria es también
suficiente para que $u$ se factorice como
$S'\longrightarrow S\longrightarrow Z$: para todo $S''\in\mathcal C$ y todo
par de morfismos $\beta_1,\beta_2:S''\longrightarrow S'$ tal que
$\alpha\beta_1=\alpha\beta_2$, se tiene también $u\beta_1=u\beta_2$.

Si existe el producto fibrado $S'\times_S S'$, equivale a decir que el diagrama
\[
  S\xleftarrow{\alpha}S'
  \overset{p_1}{\underset{p_2}{\rightrightarrows}}S'\times_S S'
\]

\src{305}
es exacto, es decir, que $S$ es un conúcleo del par $(p_1,p_2)$. En cualquier
caso, un morfismo conúcleo es un epimorfismo estricto. Señalemos también que un
epimorfismo estricto que sea un monomorfismo es un isomorfismo. Dejamos al
lector desarrollar la noción dual de \emph{monomorfismo estricto}.

Para precisar las relaciones entre la noción de morfismo de
$\mathcal F$-descenso y la de epimorfismo estricto, introducimos además las
definiciones siguientes:

\result{DEFINICIÓN}{2.3} Un morfismo $\alpha:S'\longrightarrow S$ se llama
\emph{epimorfismo universal} (resp. \emph{epimorfismo estricto universal}) si,
para todo $T$ sobre $S$, existe el producto fibrado $T'=S'\times_ST$ y la
proyección $T'\longrightarrow T$ es un epimorfismo (resp. un epimorfismo
estricto).

En las categorías muy buenas (como la categoría de los conjuntos, la
categoría de los haces de conjuntos sobre un espacio topológico, las categorías
abelianas, etc.), las cuatro nociones de epiyectividad así introducidas
coinciden; por el contrario, todas ellas son distintas en una categoría como
la de los preesquemas, o la de los preesquemas sobre un preesquema no vacío
dado $S$, incluso si nos limitamos a $S$-esquemas finitos sobre $S$.

\result{DEFINICIÓN}{2.4} Un morfismo $\alpha:S'\longrightarrow S$ se llama
\emph{morfismo de descenso} (resp. \emph{morfismo de descenso estricto}) si es
un morfismo de $\mathcal F$-descenso (resp. un morfismo de
$\mathcal F$-descenso estricto) (cf. definición 1.7), donde $\mathcal F$ denota
la categoría fibrada de base $\mathcal C$ de los objetos de $\mathcal C$ sobre
objetos de $\mathcal C$ (n\textsuperscript{o} 1, ejemplo 1).

\result{PROPOSICIÓN}{2.1} Si en $\mathcal C$ existen los productos finitos y
los productos fibrados (finitos), entonces, en $\mathcal C$, los morfismos de
descenso coinciden con los epimorfismos estrictos universales.

\result{b. EJEMPLOS}{} Sea $\mathcal C$ la categoría de los preesquemas. Sea
$S\in\mathcal C$, y sean $S'$ y $S''$ dos preesquemas finitos sobre $S$, es
decir, correspondientes a haces de álgebras $\mathcal A'$, $\mathcal A''$
sobre $S$ que, como haces de módulos, son cuasi-coherentes y de tipo finito (es
decir, coherentes si $S$ es localmente noetheriano). Sea
$\alpha:S'\longrightarrow S$ el morfismo estructural de $S'$, y sean
$\beta_1,\beta_2$ dos $S$-morfismos de $S''$ en $S'$, definidos por
homomorfismos de álgebras $\mathcal A'\longrightarrow\mathcal A''$, denotados
también por $\beta_1,\beta_2$. Utilizando el hecho de que un morfismo finito es
cerrado (primer teorema de Cohen--Seidenberg), se demuestra fácilmente que el
diagrama
\[
  \tag{+}
  S\xleftarrow{\alpha}S'
  \overset{\beta_1}{\underset{\beta_2}{\rightrightarrows}}S''
\]
en $\mathcal C$ es exacto si y solo si lo es el diagrama de haces

\src{306}
\[
  \mathcal O_S=\mathcal A\longrightarrow\mathcal A'
  \overset{\beta_1}{\underset{\beta_2}{\rightrightarrows}}\mathcal A''
\]
sobre $S$. En particular, si $\alpha:S'\longrightarrow S$ es un morfismo
finito correspondiente a un haz $\mathcal A'$ de álgebras sobre $S$, entonces
$\alpha$ es un epimorfismo estricto si y solo si el diagrama de haces
\[
  \mathcal O_S=\mathcal A\longrightarrow\mathcal A'
  \overset{p_1}{\underset{p_2}{\rightrightarrows}}
  \mathcal A'\otimes_{\mathcal A}\mathcal A'
\]
es exacto (es un epimorfismo si y solo si
$\mathcal A\longrightarrow\mathcal A'$ es inyectivo). Si $S$ es afín de anillo
$A$, y por tanto $S'$ es afín de anillo $A'$ finito sobre $A$, entonces
$S'\longrightarrow S$ es un epimorfismo estricto si y solo si
$A\longrightarrow A'$ es un isomorfismo de $A$ sobre el subanillo de $A'$
formado por los $x'\in A'$ tales que
\[
  1_{A'}\otimes_Ax'-x'\otimes_A1_{A'}=0
\]
(es un epimorfismo si y solo si $A\longrightarrow A'$ es inyectivo). Como ya
hemos señalado, incluso si $S$ es el esquema de un anillo local artiniano, un
morfismo finito $S'\longrightarrow S$ que sea un epimorfismo no es
necesariamente un epimorfismo estricto. Sin embargo, puede demostrarse que, si
$S$ es un preesquema noetheriano, todo morfismo finito
$S'\longrightarrow S$ que sea un epimorfismo es la composición de una sucesión
finita de epimorfismos estrictos (también finitos). Esto muestra, además, que
la composición de dos epimorfismos estrictos no es necesariamente un
epimorfismo estricto.

\result{c}{} Si $(+)$ es un diagrama exacto de morfismos finitos, entonces,
para todo morfismo plano $T\longrightarrow S$ de preesquemas, el diagrama
obtenido de $(+)$ mediante el cambio de base $T\longrightarrow S$ sigue siendo
exacto. De ello se deduce que, si $X,Y$ son dos $S$-preesquemas, siendo $X$
plano sobre $S$, entonces el siguiente diagrama de aplicaciones de conjuntos
(donde $X',Y'$ son las imágenes inversas de $X,Y$ sobre $S'$, y $X'',Y''$ sus
imágenes inversas sobre $S''$)
\[
  \operatorname{Hom}_S(X,Y)\longrightarrow
  \operatorname{Hom}_{S'}(X',Y')
  \rightrightarrows
  \operatorname{Hom}_{S''}(X'',Y'')
\]
es exacto. En particular, si $\mathcal F$ denota la categoría fibrada cuya base
es la categoría $\mathcal C$ de los preesquemas y tal que, para
$X\in\mathcal C$, $\mathcal F_X$ es la categoría de los $X$-preesquemas
planos, entonces el diagrama $(+)$ es $\mathcal F$-exacto. (Este resultado deja
de ser cierto si no se impone la hipótesis de planitud; en particular, un
epimorfismo estricto finito no es necesariamente un morfismo de descenso.) Se
ve asimismo que $(+)$ es $\mathcal F$-exacto si $\mathcal F$ denota la
categoría fibrada para la cual $\mathcal F_X$ es la categoría de los haces
cuasi-coherentes y planos sobre el preesquema $X$ (también aquí la hipótesis

\src{307}
de planitud es esencial). En uno y otro caso, la cuestión de la efectividad de
un dato de pegado (y, más concretamente, de un dato de descenso cuando
$S''=S'\times_SS'$) sobre un objeto plano sobre $S'$ es delicada (y su
respuesta en diversos casos particulares es uno de los objetos principales de
los presentes exposés). El conferenciante ignora si, para todo epimorfismo
estricto finito $S'\longrightarrow S$, todo dato de descenso sobre un haz
cuasi-coherente plano sobre $S'$ es efectivo (incluso suponiendo que $S$ es el
espectro de un anillo local artiniano y limitándose a los haces localmente
libres de rango 1). Más generalmente, sean $A$ un anillo y $A'$ una
$A$-álgebra (todo es conmutativo) tal que el diagrama de aplicaciones
\[
  A\longrightarrow A'\rightrightarrows A'\otimes_AA'
\]
sea exacto, lo que equivale también a que el morfismo
$S'\longrightarrow S$ correspondiente entre los espectros de $A,A'$ sea un
morfismo de $\mathcal F$-descenso, donde $\mathcal F$ es la categoría fibrada
de los haces cuasi-coherentes planos. Sea $M'$ un $A'$-módulo plano provisto de
un dato de descenso a $A$, es decir, de un isomorfismo
\[
  \varphi:M'\otimes_AA'\xrightarrow{\sim}A'\otimes_AM'
\]
de $A'\otimes_AA'$-módulos que satisface las condiciones (i) y (ii) de 1.c
(cuya explicitación en términos de módulos dejamos al lector). ¿Es efectivo
este dato (respecto de la categoría fibrada de los haces cuasi-coherentes
planos)? Sea $M$ el subconjunto de $M'$ formado por los $x'\in M'$ tales que
\[
  \varphi(x'\otimes_A1_{A'})=1_{A'}\otimes_Ax',
\]
que es un sub-$A$-módulo de $M'$. La inyección canónica
$M\longrightarrow M'$ define un homomorfismo de $A'$-módulos
$M\otimes_AA'\longrightarrow M'$. La efectividad de $\varphi$ significa
entonces lo siguiente: $M$ es un $A$-módulo plano y el homomorfismo anterior es
un isomorfismo.

\result{OBSERVACIÓN}{} En las consideraciones anteriores no habíamos impuesto
ninguna hipótesis de planitud a los morfismos del diagrama $(+)$, lo que nos
obligaba, para disponer de una técnica de descenso, a imponer hipótesis de
planitud a los objetos sobre $S,S'$ considerados. En el parágrafo 2 impondremos
una hipótesis de planitud a $\alpha:S'\longrightarrow S$, lo que nos permitirá
disponer de una técnica de descenso para objetos sobre $S,S'$ que ya no están
sometidos a condición alguna de planitud. En todos los casos interviene una
hipótesis de planitud. Esta es una de las principales razones de la importancia
de la noción de planitud en geometría algebraica (¡cuyo papel no podía
manifestarse mientras uno se limitara a cuerpos base, sobre los cuales, en
efecto, todo es plano!).

\src{308}
\subsection*{3. Aplicación a los morfismos étale}
\addcontentsline{toc}{subsection}{3. Aplicación a los morfismos étale}

Sean $A$ un anillo local y $B$ un álgebra local sobre $A$ cuyo ideal maximal
induce el de $A$. Diremos que $B$ es \emph{étale} sobre $A$ (en lugar de
«no ramificada», expresión usada en otros lugares) si satisface las condiciones
siguientes:
\begin{enumerate}
\item[(i)] $B$ es plana sobre $A$;
\item[(ii)] $B/\mathfrak mB$ es una extensión finita separable de
$A/\mathfrak m=k$ (donde $\mathfrak m$ denota el ideal maximal de $A$).
\end{enumerate}
Cuando $A$ y $B$ son noetherianos y $k$ es algebraicamente cerrado, esto
significa que el homomorfismo $\widehat A\longrightarrow\widehat B$ entre los
completados que prolonga $A\longrightarrow B$ es un isomorfismo. Un morfismo de
tipo finito $f:T\longrightarrow S$ se dice \emph{étale en $x\in T$}, o también
que $T$ es \emph{étale sobre $S$ en $x$}, si $\mathcal O_x$ es étale sobre
$\mathcal O_{f(x)}$; se dice que $f$ es \emph{étale}, o se llama a $f$ un
\emph{morfismo étale}, o se dice que $T$ es \emph{étale sobre $S$}, si $f$ es
étale en todo $x\in T$. Señalemos, además, que si $S$ es localmente
noetheriano, el conjunto de puntos de $T$ en los que $f$ es étale es abierto, y
el uso del «main theorem» de Zariski permite precisar la estructura de $T/S$
en un entorno de tal punto (mediante una ecuación de tipo bien conocido).

Cuando $S$ es un esquema de tipo finito sobre el cuerpo de los números
complejos, le corresponde un espacio analítico $\overline S$ en el sentido de
Serre [5], salvo que $\overline S$ puede tener elementos nilpotentes en su haz
estructural, lo que no cambia nada esencial en [5]. Se ve entonces fácilmente
que $f$ es étale si y solo si
$\overline f:\overline T\longrightarrow\overline S$ lo es, es decir, si todo
punto de $\overline T$ admite un entorno en el que $\overline f$ induce un
isomorfismo sobre un abierto de $\overline S$. En particular, a todo
\emph{recubrimiento étale} $T$ de $S$ (es decir, un morfismo finito étale
$f:T\longrightarrow S$) le corresponde un recubrimiento étale $\overline T$ de
$\overline S$, que es conexo si y solo si $T$ lo es [5]. Se ve además
fácilmente que, si $T,T'$ son dos esquemas étale sobre $S$, entonces la
aplicación natural
\[
  \operatorname{Hom}_S(T,T')\longrightarrow
  \operatorname{Hom}_{\overline S}(\overline T,\overline{T'})
\]
es biyectiva, es decir, el funtor $T\longmapsto\overline T$ de la categoría de
los esquemas étale sobre $S$ a la categoría de los espacios analíticos étale
sobre $\overline S$ es «plenamente fiel» y, por tanto, define una equivalencia
de la primera categoría con una subcategoría de la segunda. Un teorema de
Grauert--Remmert [2] implica que, si $S$ es normal, se obtiene así una
equivalencia entre la categoría de los \emph{recubrimientos étale} de $S$ y la
categoría de los recubrimientos étale (finitos) de $\overline S$, es decir, que
todo recubrimiento étale de $\overline S$ es $\overline S$-isomorfo a un
$\overline T$, donde $T$ es un recubrimiento étale de $S$. Demostremos que el
\emph{teorema de Grauert--Remmert sigue siendo válido sin hipótesis de
normalidad sobre $S$}. Sea, en efecto, primero $S'\longrightarrow S$ un
epimorfismo estricto finito; supongamos demostrado el

\src{309}
teorema para $S'$ y demostremos que será cierto para $S$. Sea
$\mathfrak T$ un recubrimiento étale de $\overline S$ y consideremos su imagen
inversa $\mathfrak T'$ sobre $\overline{S'}$, que corresponde a un haz
analítico coherente $\mathcal A'$ de álgebras sobre $\overline{S'}$, imagen
inversa del haz de álgebras $\mathcal A$ sobre $\overline S$ que define
$\mathfrak T$. Por hipótesis, sobre $\overline{S'}$, $\mathfrak T'$ proviene de
un recubrimiento étale $T'$ de $S'$, es decir, $\mathcal A'$ proviene de un haz
coherente de álgebras $A'$ sobre $S'$. Por otra parte, $\mathcal A'$ está
provisto de un dato de descenso canónico relativo a
$\overline{S'}\longrightarrow\overline S$, es decir, de un isomorfismo entre
sus dos imágenes inversas sobre
\[
  \overline{S'}\times_{\overline S}\overline{S'}
  =\overline{(S'\times_SS')},
\]
(que satisface las condiciones (i), (ii)), y este isomorfismo proviene, según
lo dicho, de un isomorfismo entre los haces algebraicos correspondientes, es
decir, de un dato de descenso sobre $A'$ relativo a
$S'\longrightarrow S$. Se comprueba fácilmente que este último es efectivo
(pues el de $\mathcal A'$ lo es, y la efectividad de un dato de descenso, tal
como se explicitó en el número anterior, se reconoce localmente en los
\emph{completados} de los módulos que intervienen). Se obtiene así un haz
coherente de álgebras $A$ sobre $S$, que define un recubrimiento $T$ de $S$,
el recubrimiento buscado. El resultado anterior sigue siendo entonces
manifiestamente válido si $S'\longrightarrow S$ es solo una composición de un
número finito de epimorfismos estrictos finitos, es decir, si es un
epimorfismo finito cualquiera (por el resultado de factorización señalado en
el parágrafo 2). Se sigue que el teorema de Grauert--Remmert continúa siendo
válido si $S$ es un esquema \emph{reducido}, es decir, si $\mathcal O_S$ no
tiene elementos nilpotentes, como se ve introduciendo su normalización $S'$.
De aquí se pasa fácilmente al caso general.

Una demostración enteramente análoga, que utiliza de nuevo el resultado de
factorización para los epimorfismos estrictos finitos y la naturaleza «formal»
de la efectividad de los datos de descenso, permite demostrar el resultado
siguiente: sea $S$ un preesquema localmente noetheriano y sea
$S'\longrightarrow S$ un morfismo \emph{finito, sobreyectivo y radicial} (o,
lo que equivale a lo mismo, un morfismo de tipo finito tal que, para todo $T$
sobre $S$, el morfismo $T'=S'\times_ST\longrightarrow T$ sea un homeomorfismo,
lo cual también se expresa diciendo que $S'\longrightarrow S$ es un
\emph{homeomorfismo universal}). Para todo $T$ étale sobre $S$, consideremos
su imagen inversa $T'=T\times_SS'$, que es étale sobre $S'$. Entonces el
funtor $T\longmapsto T'$ es una \emph{equivalencia de categorías entre la
categoría de los preesquemas $T$ étale sobre $S$ y la categoría de los
preesquemas $T'$ étale sobre $S'$}. (Se utiliza la biyectividad de
\[
  \operatorname{Hom}_S(T_1,T_2)\longrightarrow
  \operatorname{Hom}_{S'}(T'_1,T'_2)
\]
para dos preesquemas $T_1,T_2$ étale sobre $S$, hecho cuya comprobación directa
es fácil, y el hecho de que el teorema enunciado es cierto si
$S'=(S,\mathcal O_S/\mathcal J)$,

\src{310}
donde $\mathcal J$ es un haz coherente nilpotente de ideales de $\mathcal O_S$
([4], lema 6)). Obsérvese, además, que no suponemos aquí que los $T$
considerados sean finitos sobre $S$. Este resultado implica, en particular,
que el morfismo $S'\longrightarrow S$ induce un isomorfismo del grupo
fundamental $[\,]$ de $S'$ sobre el de $S$ («\emph{invariancia topológica del
grupo fundamental de un preesquema}»).

\subsection*{4. Relaciones con la 1-cohomología}
\addcontentsline{toc}{subsection}{4. Relaciones con la 1-cohomología}

\medskip\noindent
a. Sea $\mathcal C$ una categoría en la que siempre existe el producto de dos
objetos, y sea $T\in\mathcal C$. Para todo conjunto finito
$I\ne\varnothing$, se puede considerar $T^I$; al variar $I$, se obtiene así un
funtor covariante de la categoría de los conjuntos finitos no vacíos a
$\mathcal C$, es decir, lo que puede llamarse un \emph{objeto simplicial} de
$\mathcal C$, denotado por $K_T$. Este último depende covariantemente de $T$;
además, si $u,v$ son dos morfismos $T\longrightarrow T'$, los morfismos
correspondientes $K_T\longrightarrow K_{T'}$ son homotópicos. Diremos que $T$
\emph{domina} a $T'$ si $\operatorname{Hom}(T,T')\ne\varnothing$; esta es una
relación de preorden filtrante creciente en $\mathcal C$. De lo anterior se
deduce que, si $T$ domina a $T'$, existe una clase canónica (salvo homotopía)
de homomorfismos de objetos simpliciales $K_T\longrightarrow K_{T'}$; en
particular, si $K_T$ y $K_{T'}$ son tales que cada uno domina al otro, entonces
$K_T$ y $K_{T'}$ son homotópicamente equivalentes. Sea ahora $F$ un funtor
(contravariante, para fijar ideas) de $\mathcal C$ a una categoría abeliana
$\mathcal C'$; entonces
\[
  C^*(T,F)=F(K_T)
\]
es un objeto cosimplicial de $\mathcal C'$ y, por tanto, define de la manera
habitual un complejo (de cocadenas) en $\mathcal C'$, del cual puede tomarse la
cohomología
\[
  H^*(T,F)=H^*(C^*(T,F))=H^*(F(K_T))
\]
(se podrá poner una $\mathcal C$ como subíndice de $H^*$ si existe posibilidad
de confusión). Este es un funtor cohomológico en $F$, cuya varianza al variar
$T$ se deduce de lo dicho sobre los $K_T$; más precisamente, para $F$ fijo y
$T$ variable en $\mathcal C$ (preordenada por la relación de dominación), los
$H^*(T,F)$ forman un sistema inductivo de objetos graduados de $\mathcal C'$;
en particular, si $T$ y $T'$ son tales que cada uno domina al otro, entonces
$H^*(T,F)$ y $H^*(T',F)$ son canónicamente isomorfos.

Supongamos que en $\mathcal C$ existen los productos fibrados. Para
$S\in\mathcal C$ fijo, se puede aplicar lo anterior a la categoría
$\mathcal C_S$ de los objetos de $\mathcal C$ sobre $S$; escribiremos
$C^*(T/S,F)$ y $H^*(T/S,F)$ en lugar de $C^*(T,F)$ y $H^*(T,F)$ cuando
queramos precisar que trabajamos en la categoría $\mathcal C_S$; así,

\src{311}
$C^*(T/S,F)$ es un complejo de cocadenas en $\mathcal C'$ que, en dimensión
$n$, es igual a
\[
  F(T\times_ST\times_S\cdots\times_ST)
\]
(donde el paréntesis contiene $n+1$ factores).

Obsérvese que, como de costumbre, puede definirse $H^0(T/S,F)$ sin suponer
abeliana la categoría $\mathcal C'$: es el núcleo (definición 2.1), si existe,
del par de morfismos $F(p_i)$ ($i=1,2$)
\[
  F(T)\rightrightarrows F(T\times_ST)
\]
correspondientes a las dos proyecciones
$p_1,p_2:T\times_ST\rightrightarrows T$. En particular, se tendrá un morfismo
natural (llamado de \emph{aumento})
\[
  F(S)\longrightarrow H^0(T/S,F)
\]
que será un isomorfismo en los casos favorables (en particular, si
$T\longrightarrow S$ es un epimorfismo estricto y $F$ es «exacto por la
izquierda»). De igual modo, cuando $F$ toma valores en la categoría de los
grupos en una categoría $\mathcal C''$, también puede definirse $H^1(T/S,F)$;
cuando $\mathcal C''$ es la categoría de los conjuntos (es decir, $F$ toma
valores en la categoría de los grupos ordinarios, no necesariamente
conmutativos), $H^1(T/S,F)$ es el cociente del subgrupo $Z^1(T/S,F)$ de
$C^1(T/S,F)=F(T\times_ST)$ formado por los $g$ tales que
\[
  F(p_{31})(g)=F(p_{32})(g)\,F(p_{21})(g)
\]
por el grupo de operadores $F(T)$, que actúa sobre $C^1(T/S,F)$ y, en
particular, sobre el subconjunto $Z^1(T/S,F)$ mediante
\[
  \rho(g')\cdot g=F(p_2)(g')\,g\,F(p_1)(g')^{-1}.
\]

\medskip\noindent
b. Sea, por ejemplo, $\mathcal F$ una categoría fibrada de base $\mathcal C$.
Sean $\xi,\eta\in\mathcal F_S$ y, para todo $S'$ sobre $S$, sea
\[
  F_{\xi,\eta}(S')=
  \operatorname{Hom}(\xi\times_SS',\eta\times_SS').
\]
Así, $F_{\xi,\eta}$ es un funtor contravariante de $\mathcal C_S$ a la
categoría de los conjuntos. Dicho esto, \emph{afirmar que el morfismo de
aumento}
\[
  F_{\xi,\eta}(S)\longrightarrow H^0(S'/S,F_{\xi,\eta})
\]
\emph{es un isomorfismo para todo par de elementos
$\xi,\eta\in\mathcal F_S$ significa que
$\alpha:S'\longrightarrow S$ es un morfismo de $\mathcal F$-descenso}
(definición 1.7).

\medskip\noindent
c. Pongamos asimismo, para $\xi\in\mathcal F_S$ y todo objeto $S'$ de
$\mathcal C$ sobre $S$,
\[
  G_\xi(S')=\operatorname{Aut}(\xi\times_SS') ;
\]
se ha definido así un funtor contravariante $G_\xi$ de $\mathcal C_S$ a la
categoría de los

\src{312}
grupos. Dicho esto, se constata que $Z^1(S'/S,G_\xi)$ \emph{se identifica
canónicamente con el conjunto de los datos de descenso sobre
$\xi'=\xi\times_SS'$ relativos a $S'\longrightarrow S$} (definición 1.6), y
$H^1(S'/S,G_\xi)$ \emph{se identifica con el conjunto de las clases (salvo
isomorfismo) de objetos de $\mathcal F_{S'}$ provistos de un dato de descenso
relativo a $\alpha:S'\longrightarrow S$ que, como objetos de
$\mathcal F_{S'}$, son isomorfos a $\xi'=\xi\times_SS'$}. Por tanto, si
$\alpha:S'\longrightarrow S$ es un morfismo de $\mathcal F$-descenso
(cf. (b)), entonces $H^1(S'/S,G_\xi)$ \emph{contiene como subconjunto el
conjunto de las clases (salvo isomorfismo) de objetos $\eta$ de
$\mathcal F_S$ tales que $\eta\times_SS'$ sea isomorfo en
$\mathcal F_{S'}$ a $\xi\times_SS'$; y esta inclusión es una identidad si y
solo si todo dato de descenso sobre $\xi'=\xi\times_SS'$ relativo a
$\alpha:S'\longrightarrow S$ es efectivo}. (Este será el caso, en particular,
si $\alpha:S'\longrightarrow S$ es un morfismo de $\mathcal F$-descenso
estricto.)

\result{OBSERVACIÓN}{} Los complejos de cocadenas del tipo $C^*(T/S,F)$
contienen como casos particulares la mayoría de los complejos estándar
conocidos (cohomología de Čech, cohomología de grupos, etc.) y desempeñan un
papel importante en geometría algebraica, en particular en la «cohomología de
Weil» de los preesquemas.

\medskip\noindent
d. \result{EJEMPLO}{1} Sea $S'$ un objeto sobre $S\in\mathcal C$, y sea
$\Gamma$ un grupo de automorfismos de $S'$ tal que $S'$ sea «formalmente
principal sobre $S$, de grupo $\Gamma$», es decir, tal que el morfismo natural
\[
  \Gamma\times S'\longrightarrow S'\times_SS',
\]
donde $\Gamma\times S'$ denota la suma directa de $\Gamma$ copias de $S'$, sea
un isomorfismo. (Se supone que en $\mathcal C$ existen las sumas directas que
intervienen aquí.) Sea $F$ un funtor contravariante de $\mathcal C$ a la
categoría de los grupos abelianos. Entonces $C^*(S'/S,F)$ es canónicamente
isomorfo al grupo simplicial estándar de cocadenas homogéneas
$C^*(\Gamma,F(S'))$ y, por tanto, $H^*(S'/S,F)$ es canónicamente isomorfo a
$H^*(\Gamma,F(S'))$.

\medskip\noindent
e. \result{EJEMPLO}{2} Sea $\mathcal C$ la categoría de los preesquemas. Se
denota por $G_a$ («grupo aditivo») el funtor contravariante de $\mathcal C$ a
la categoría de los grupos abelianos definido por
\[
  G_a(X)=H^0(X,\mathcal O_X).
\]
Se define de igual modo el funtor $G_m$ («grupo multiplicativo») mediante
\[
  G_m(X)=H^0(X,\mathcal O_X)^*
\]
(grupo de los elementos invertibles del anillo $H^0(X,\mathcal O_X)$), y, más
generalmente, el funtor $\operatorname{Gl}(n)$ («grupo lineal de orden $n$»)
mediante

\src{313}
\[
  \operatorname{Gl}(n)(X)=
  \operatorname{Gl}\bigl(n,H^0(X,\mathcal O_X)\bigr),
\]
que es un funtor de $\mathcal C$ a la categoría de los grupos (no
necesariamente conmutativos si $n>1$; para $n=1$ se recupera $G_m$). Además,
puede interpretarse $\operatorname{Gl}(n)$ como un funtor de automorfismos
(cf. (c)) considerando la categoría fibrada $\mathcal F$ de base $\mathcal C$
tal que, para $X\in\mathcal C$, $\mathcal F_X$ sea la categoría de los haces
localmente libres sobre $X$: en efecto, se tiene
\[
  \operatorname{Gl}(n)(X)=
  \operatorname{Aut}_{\mathcal F_X}(\mathcal O_X^n).
\]
De (b) se sigue que, si $\alpha:S'\longrightarrow S$ es un morfismo de
$\mathcal F$-descenso (cf. parágrafo 2.c),
$H^1(S'/S,\operatorname{Gl}(n))$ contiene el conjunto de las clases (salvo
isomorfismo) de haces localmente libres sobre $S$ cuya imagen inversa sobre
$S'$ es isomorfa a $\mathcal O_{S'}^n$, y esta inclusión es una igualdad si y
solo si todo dato de descenso sobre $\mathcal O_{S'}^n$ (relativo a
$\alpha:S'\longrightarrow S$) es efectivo. Cuando $S$ es el espectro de un
anillo local, esto significa, por tanto,
$H^1(S'/S,\operatorname{Gl}(n))=(e)$, pues todo haz localmente libre sobre $S$
es entonces trivial.

Señalemos la equivalencia de las condiciones siguientes para un morfismo
$\alpha:S'\longrightarrow S$:
\begin{enumerate}
\item[(i)] el homomorfismo de aumento
\[
  H^0(S,\mathcal O_S)=G_a(S)\longrightarrow H^0(S'/S,G_a)
\]
es un isomorfismo;
\item[(ii)] $S'\longrightarrow S$ es un morfismo de $\mathcal F$-descenso
(siendo $\mathcal F$ la categoría fibrada de base $\mathcal C$ considerada
arriba).
\end{enumerate}
Si $S'\longrightarrow S$ es finito, estas condiciones equivalen también a:
\begin{enumerate}
\item[(iii)] $S'\longrightarrow S$ es un epimorfismo estricto
(cf. parágrafo 2.c).
\end{enumerate}

Supongamos ahora que $S=\operatorname{Spec}(A)$,
$S'=\operatorname{Spec}(A')$; entonces
\[
  C^n(S'/S,G_a)=C^n(A'/A,G_a)=\bigotimes_A^{n+1}A',
\]
siendo el operador coborde
$C^n(A'/A,G_a)\longrightarrow C^{n+1}(A'/A,G_a)$ la suma alternada de los
operadores de cara
\[
  \partial_i(x_0\otimes x_1\otimes\cdots\otimes x_n)
  =x_0\otimes\cdots\otimes x_{i-1}\otimes1_{A'}\otimes
   x_i\otimes\cdots\otimes x_n.
\]
Del mismo modo,
\[
  C^n(S'/S,G_m)=C^n(A'/A,G_m)
  \simeq\left(\bigotimes_A^{n+1}A'\right)^*,
\]
y las operaciones simpliciales en $C^*(A'/A,G_m)$ están inducidas por las de
$C^*(S'/S,G_a)$. Se explicitan análogamente las operaciones simpliciales en
$C^*(A'/A,\operatorname{Gl}(n))$. \emph{En todos los casos que conoce el
conferenciante, se tiene}
\[
  H^i(A'/A,G_a)=0\qquad\text{para $i>0$},
\]
\emph{y si $A$ es local, se tiene}
\[
  H^1(A'/A,G_m)=0
  \quad\text{y, más generalmente,}\quad
  H^1(A'/A,\operatorname{Gl}(n))=(e)
\]
(cuando $S'\longrightarrow S$ es un morfismo de $\mathcal F$-descenso, es
decir, cuando el diagrama
$A\longrightarrow A'\rightrightarrows A'\otimes_AA'$ es exacto; compárese con
el parágrafo 2.c). Obsérvese que \emph{el «teorema 90» de Hilbert no es sino
la relación}
\src{314}
\[
  H^1(S'/S,G_m)=0
\]
\emph{cuando $A$ es un cuerpo y $A'$ una extensión de Galois finita de este
(cf. ejemplo 1), y también puede expresarse diciendo que, en el caso
considerado, $S'\longrightarrow S$ es un morfismo de descenso estricto para la
categoría fibrada de los haces localmente libres de rango 1.} Es bajo esta
última forma como conviene generalizar el teorema de Hilbert, variando tanto
las hipótesis sobre el morfismo $S'\longrightarrow S$ como las relativas a los
haces cuasi-coherentes considerados.

Señalemos, por último, la equivalencia de las propiedades siguientes cuando
$A$ es un anillo local artiniano de ideal maximal $\mathfrak m$ y $A'$ una
$A$-álgebra (denotando, para todo entero $k>0$, por $A_k$ (resp. $A'_k$) los
anillos $A/\mathfrak m^{k+1}$ (resp. $A'/\mathfrak m^{k+1}A'$)):
\begin{enumerate}
\item[(i)] $H^1(A'_k/A_k,G_a)=0$ para todo $k$.
\item[(ii)] $H^1(A'_k/A_k,G_m)=0$ para todo $k$.
\item[(iii)] $H^1(A'_k/A_k,\operatorname{Gl}(n))=(e)$ para todo $k$ y todo
$n$.
\end{enumerate}
Si $S'\longrightarrow S$ es un epimorfismo estricto, las condiciones
anteriores implican incluso que es un morfismo de descenso estricto para los
módulos libres (de tipo finito o no) sobre $A'$.

\result{OBSERVACIÓN}{} La definición de los grupos $H^i(S'/S,G_m)$, cuando
$S,S'$ son esquemas de cuerpos $A,A'$, se debe a Amitsur. El grupo
$H^2(S'/S,G_m)$ es particularmente interesante como variante «global» del
grupo de Brauer; para esta variante se podrá consultar [1], capítulo VII.

\fsection{B}{Descenso por morfismos fielmente planos}

\subsection*{1. Enunciado de los teoremas de descenso}
\addcontentsline{toc}{subsection}{1. Enunciado de los teoremas de descenso}

\result{DEFINICIÓN}{1.1} Un morfismo
$\alpha:S'\longrightarrow S$ de preesquemas se dice \emph{plano} si, para todo
$x'\in S'$, $\mathcal O_{x'}$ es un módulo plano sobre el anillo
$\mathcal O_{\alpha(x')}$ (es decir, si
$\mathcal O_{x'}\otimes_{\mathcal O_{\alpha(x')}}M$ es un funtor exacto en el
$\mathcal O_{\alpha(x')}$-módulo $M$). Un morfismo se dice
\emph{fielmente plano} si es plano y sobreyectivo.

Por ejemplo, si $S=\operatorname{Spec}(A)$ y
$S'=\operatorname{Spec}(A')$, entonces $S'$ es plano sobre $S$ si y solo si
$A'$ es un $A$-módulo plano, y $S'$ es fielmente plano sobre $S$ si y solo si
$A'$ es un $A$-módulo fielmente plano (es decir, el funtor
$A'\otimes_A M$ en el $A$-módulo $M$ es exacto y fiel); esto significa también,
en la terminología de Serre [5], que el par $(A,A')$ es plano. Si $S'$ es
fielmente plano sobre $S$, entonces el funtor imagen inversa de haces
cuasi-coherentes sobre $S$ es exacto y fiel; en otras palabras, para que una
sucesión de homomorfismos de haces

\src{315}
cuasi-coherentes sobre $S$ sea exacta, es necesario y suficiente que lo sea su
imagen inversa sobre $S'$ (en particular, para que un homomorfismo de haces
cuasi-coherentes sobre $S$ sea un monomorfismo, resp. un epimorfismo, resp. un
isomorfismo, es necesario y suficiente que lo sea su imagen inversa sobre
$S'$). Esta propiedad sigue siendo cierta si nos restringimos sobre un abierto
cualquiera de $S$ y, bajo esta forma, caracteriza a los morfismos fielmente
planos.

\result{DEFINICIÓN}{1.2} Un morfismo $\alpha:S'\longrightarrow S$ se dice
\emph{cuasicompacto} si la imagen inversa de toda parte abierta cuasicompacta
$U$ de $S$ es cuasicompacta (es decir, unión finita de abiertos afines).

Evidentemente, basta comprobar esta propiedad para los abiertos afines de $S$.
Por ejemplo, un morfismo afín (es decir, tal que la imagen inversa de un abierto
afín sea afín) es cuasicompacto.

La clase de los morfismos planos, resp. fielmente planos, resp. cuasicompactos,
es estable por composición y por «cambio de base», y contiene, por supuesto,
los isomorfismos.

\result{TEOREMA}{1} Sea $\alpha:S'\longrightarrow S$ un morfismo de
preesquemas, fielmente plano y cuasicompacto. Entonces $\alpha$ es un
\emph{morfismo de descenso estricto} (A, definición 1.7) para la categoría
fibrada $\mathcal F$ de los haces cuasi-coherentes (A, parágrafo 1, ejemplo 2).

Este enunciado significa dos cosas:
\begin{enumerate}
\item[(i)] Si $F$ y $G$ son dos haces cuasi-coherentes sobre $S$, y $F'$ y
$G'$ sus imágenes inversas sobre $S'$, entonces el homomorfismo natural
\[
  \operatorname{Hom}(F,G)\longrightarrow\operatorname{Hom}(F',G')
\]
es una biyección del primer miembro sobre el subgrupo del segundo formado por
los homomorfismos $F'\longrightarrow G'$ compatibles con los datos de descenso
canónicos sobre estos haces, es decir, aquellos cuyas imágenes inversas por
las dos proyecciones de $S''=S'\times_SS'$ sobre $S'$ dan un mismo
homomorfismo $F''\longrightarrow G''$.
\item[(ii)] Todo haz cuasi-coherente $F'$ sobre $S'$, provisto de un dato de
descenso relativo al morfismo $\alpha:S'\longrightarrow S$
(A, definición 1.6), es isomorfo (provisto de este dato) a la imagen inversa de
un haz cuasi-coherente $F$ sobre $S$.
\end{enumerate}
Tomando $F=\mathcal O_S$ en (i), se obtiene:

\result{COROLARIO}{1} Sea $G$ un haz cuasi-coherente sobre $S$; sean
$G'$ y $G''$ sus imágenes inversas sobre $S'$ y sobre
$S''=S'\times_SS'$, y sean $p_1,p_2$ las dos proyecciones de $S''$ sobre
$S'$. Entonces el siguiente diagrama de aplicaciones de conjuntos

\src{316}
\[
  \Gamma(G)\xrightarrow{\alpha^*}\Gamma(G')
  \underset{p_2^*}{\overset{p_1^*}{\rightrightarrows}}\Gamma(G'')
\]
es exacto (A, definición 1.(a)).

Por otra parte, la conjunción de (i), (ii) y de la definición 1.1 da el

\result{COROLARIO}{2} Sea $G$ como en el corolario 1. Entonces existe una
correspondencia biunívoca entre los subhaces cuasi-coherentes de $G$ y los
subhaces cuasi-coherentes de $G'$ cuyas imágenes inversas sobre $S''$, por las
dos proyecciones $p_1,p_2$, dan el mismo subhaz de $G''$.

Por supuesto, se tiene un enunciado equivalente en términos de haces
cociente. Como se ha visto (A, parágrafo 4, (e)), el teorema 1 debe considerarse
una generalización del «teorema 90» de Hilbert, e implica como casos
particulares diversas formulaciones en términos de 1-cohomología. Para la
demostración, se reduce fácilmente al caso en que $S=\operatorname{Spec}(A)$,
$S'=\operatorname{Spec}(A')$, y para (i) se reduce fácilmente a demostrar el
corolario 1, es decir, la exactitud del diagrama
\[
  M=A\otimes_A M\longrightarrow A'\otimes_A M
  \rightrightarrows A'\otimes_AA'\otimes_A M
\]
para todo $A$-módulo $M$, lo cual se deduce del lema más general:

\result{LEMA}{1.1} Sea $A'$ una $A$-álgebra fielmente plana. Entonces, para
todo $A$-módulo $M$, el complejo aumentado por $M$
$C^*(A'/A,G_a)\otimes_A M$ (cf. A, parágrafo 4, (e)) es una
\emph{resolución} de $M$.

Basta demostrar que el complejo aumentado deducido del anterior mediante el
cambio de base de $A$ a $A'$ satisface las mismas conclusiones. Esto conduce a
comprobar el enunciado cuando se reemplazan $A$ por $A'$ y $A'$ por
$A'\otimes_AA'$, por lo que nos reduce a un caso en el que existe un
homomorfismo de $A$-álgebras $A'\longrightarrow A$ (o, en términos
geométricos, al caso en que $S'$ sobre $S$ admite una sección). En este caso,
esto se deduce de las generalidades de A, parágrafo 4, (a). Señalemos de paso
el siguiente corolario, que generaliza un enunciado bien conocido en
cohomología de Galois (compárese con A, parágrafo 4, (e)):

\result{COROLARIO}{} Si $A'$ es fielmente plana sobre $A$, se tiene
$H^0(A'/A,G_a)=A$ y
\[
  H^i(A'/A,G_a)=0\qquad\text{para $i\geq1$}.
\]

Para demostrar la parte (ii) del teorema 1, se procede como para (i),
reduciéndose al caso en que $S'$ sobre $S$ admite una sección, donde se deduce
de (i) (cf. A, parágrafo 1, (c)).

Evidentemente, se puede variar ad libitum el teorema 1 y sus corolarios
introduciendo diversas estructuras adicionales en los haces (o sistemas de
haces) cuasi-coherentes considerados. Por ejemplo, dar sobre $S$ un haz
\src{317}
cuasi-coherente de álgebras conmutativas «equivale» a dar sobre $S'$ un haz de
este tipo provisto de un dato de descenso relativo a
$\alpha:S'\longrightarrow S$. Teniendo en cuenta la correspondencia funtorial
entre tales haces cuasi-coherentes sobre $S$ y los preesquemas afines sobre
$S$, se obtiene la segunda afirmación del teorema siguiente:

\result{TEOREMA}{2} Sea $\alpha:S'\longrightarrow S$ como en el teorema 1.
Entonces $\alpha$ es un \emph{morfismo de descenso} (no estricto en general)
(A, definición 2.4), y es un \emph{morfismo de descenso estricto} para la
categoría fibrada de los esquemas afines sobre preesquemas
(A, definición 1.7).

La primera afirmación del teorema significa lo siguiente: sean $X,Y$ dos
preesquemas sobre $S$; sean $X',Y'$ sus imágenes inversas sobre $S'$, y
$X'',Y''$ sus imágenes inversas sobre $S''=S'\times_SS'$. Entonces el siguiente
diagrama de aplicaciones naturales
\[
  \operatorname{Hom}_S(X,Y)\xrightarrow{\alpha^*}
  \operatorname{Hom}_{S'}(X',Y')
  \underset{p_2^*}{\overset{p_1^*}{\rightrightarrows}}
  \operatorname{Hom}_{S''}(X'',Y'')
\]
es exacto; es decir, $\alpha^*$ es una biyección de
$\operatorname{Hom}_S(X,Y)$ sobre la parte de
$\operatorname{Hom}_{S'}(X',Y')$ formada por los homomorfismos compatibles con
los datos de descenso canónicos sobre $X',Y'$ (es decir, aquellos cuyas
imágenes inversas por las dos proyecciones de $S''$ sobre $S'$ son iguales).
Esto se deduce fácilmente del teorema 1, corolario 1, cuando uno se restringe
a $Y$ afín sobre $S$; en el caso general, hay que combinar el teorema 1 con el
resultado siguiente:

\result{LEMA}{1.2} Sea $\alpha:S'\longrightarrow S$ un morfismo fielmente plano
y cuasicompacto. Entonces $S$ se identifica con un \emph{espacio topológico
cociente} de $S'$; es decir, toda parte $U$ de $S$ tal que
$\alpha^{-1}(U)$ sea abierta es abierta.

Para completar el teorema 2, hay que dar criterios de efectividad para un dato
de descenso sobre un $S'$-preesquema $X'$ (cuando no se supone que $X'$ sea
afín sobre $S'$). Señalemos primero que \emph{tal dato de descenso no es
necesariamente efectivo}, incluso si $S$ es el espectro de un cuerpo $k$, $S'$
el espectro de una extensión cuadrática $k'$ de este y $X'$ un esquema
algebraico propio de dimensión 2 sobre $S'$ (como puede verse, según Serre,
utilizando la superficie no proyectiva de Nagata). \emph{Para que un dato de
descenso sobre $X'/S'$ relativo a $\alpha:S'\longrightarrow S$ (fielmente
plano y cuasicompacto) sea efectivo, es necesario y suficiente que $X'$ sea
unión de abiertos $X'_i$, afines sobre $S'$, que sean «estables» por el dato de
descenso sobre $X'$.} Esto se cumple con certeza (cualesquiera que sean $X'/S'$
y el dato de descenso sobre $X'$) si el morfismo
$\alpha:S'\longrightarrow S$ es \emph{radicial} (es decir, inyectivo y con
extensiones residuales que

\src{318}
son radiciales). Puede demostrarse también que sigue cumpliéndose si
$\alpha:S'\longrightarrow S$ es \emph{finito} y toda parte finita de $X'$,
contenida en una fibra de $X'$ sobre $S$, está contenida en un abierto de $X'$
afín sobre $S$ (este es el \emph{criterio de Weil}). En particular, esto se
cumple si $X'/S'$ es cuasiproyectivo, y en ese caso puede demostrarse que el
preesquema «descendido» $X/S$ también es cuasiproyectivo (y proyectivo si
$X'/S'$ lo es). En resumen:

\result{TEOREMA}{3} Sea $\alpha:S'\longrightarrow S$ un morfismo de
preesquemas, fielmente plano y cuasicompacto. Si $\alpha$ es \emph{radicial},
es un \emph{morfismo de descenso estricto}. Si $\alpha$ es finito, es un
morfismo de descenso estricto respecto de la categoría fibrada de los
preesquemas cuasiproyectivos (o proyectivos) sobre preesquemas.

\result{OBSERVACIONES}{} Ignoro si, en la segunda afirmación anterior, la
hipótesis de que $\alpha$ sea un morfismo \emph{finito} es realmente necesaria;
en cualquier caso, se comprueba de manera «formal» que puede sustituirse por la
siguiente hipótesis, aparentemente más débil: \emph{para todo punto de $S$
existe un entorno abierto $U$, un $U'$ finito y fielmente plano sobre $U$, y un
$S$-morfismo de $U'$ en $S'$.} Un caso típico que no entra en el anterior es
aquel en que $S=\operatorname{Spec}(A)$,
$S'=\operatorname{Spec}(\widehat A)$, donde $A$ es un anillo local noetheriano
y $\widehat A$ su completado; o también aquel en que $S'$ es cuasifinito sobre
$S$ (es decir, localmente isomorfo a un abierto de un $S$-esquema finito) y no
finito. En estos dos casos, el conferenciante también ignora la respuesta a la
pregunta siguiente: sea $X$ un $S$-esquema tal que $X'=X\times_SS'$ sea
proyectivo sobre $S'$; ¿es cierto que $X$ es proyectivo sobre
$S$?\footnote{La fe de erratas de 1962 responde negativamente: un morfismo
$S'\longrightarrow S$ cuasifinito, étale y sobreyectivo, o un morfismo
$\operatorname{Spec}(\widehat A)\longrightarrow\operatorname{Spec}(A)$, no es
siempre un morfismo de descenso estricto, incluso si $A$ es el anillo local de
una curva algebraica sobre un cuerpo algebraicamente cerrado $k$ y
$S=\operatorname{Spec}(A)$. Se puede encontrar así un morfismo propio y liso
$f:X\longrightarrow S$, que hace de $X$ un fibrado principal de base $S$ bajo
una curva elíptica $E$, tal que $f':X'\longrightarrow S'$ sea proyectivo, pero
$f$ no lo sea. Por tanto, también es un ejemplo de fibrado principal homogéneo
no isotrivial bajo un esquema abeliano.}

\subsection*{2. Aplicación al descenso de ciertas propiedades de los morfismos}
\addcontentsline{toc}{subsection}{2. Descenso de propiedades de los morfismos}

Sea $P$ una clase de morfismos de preesquemas. Sea
$\alpha:S'\longrightarrow S$ un morfismo de preesquemas, sea
$f:X\longrightarrow Y$ un morfismo de $S$-preesquemas y sea
$f':X'\longrightarrow Y'$ la imagen inversa de $f$ por $\alpha$. Cabe
preguntarse entonces si la relación «$f'\in P$» implica «$f\in P$». Resulta que
la respuesta es afirmativa en muchos casos importantes cuando se supone que
$\alpha$ es \emph{fielmente plano y cuasicompacto} (siendo a veces superflua
esta última hipótesis). Esto se ve directamente y sin dificultad si $P$ es la
clase de los morfismos sobreyectivos, resp. radiciales (estos casos se deducen
de la sobreyectividad de $\alpha$), resp. planos, resp. fielmente planos, resp.
lisos (estos casos se deducen de la planitud fiel de $\alpha$), resp. de tipo
finito. Utilizando los teoremas 1 y 2 y el lema 1.2, se ve también que ocurre
lo mismo si $P$ es una de las clases siguientes: isomorfismos, inmersiones
abiertas, inmersiones cerradas, inmersiones (si $f$ es de tipo finito e $Y$
localmente noetheriano), morfismos afines, morfismos finitos, morfismos
cuasifinitos, morfismos abiertos, morfismos cerrados, homeomorfismos, morfismos
separados,

\src{319}
morfismos propios. El único caso importante no dilucidado es el de los
morfismos proyectivos y cuasiproyectivos, ya señalado en la observación del
parágrafo 1.

\subsection*{3. Descenso por morfismos finitos fielmente planos}
\addcontentsline{toc}{subsection}{3. Descenso por morfismos finitos fielmente planos}

Sea $\alpha:S'\longrightarrow S$ un morfismo finito correspondiente a un haz de
álgebras $\mathcal A'$ sobre $S$ que, como haz de módulos, es localmente libre
de tipo finito y es $\ne0$ en todas partes; entonces $\alpha$ es un morfismo
fielmente plano y cuasicompacto, al que pueden aplicarse los resultados
anteriores. Dar un haz cuasi-coherente $F'$ sobre $S'$ equivale a dar el haz
cuasi-coherente $\alpha_*(F')$ sobre $S$, provisto de su estructura de
$\mathcal A'$-módulo (observando que
$\mathcal A'=\alpha_*(\mathcal O_{S'})$). Para simplificar, este haz sobre $S$
también se denotará por $F'$. Las dos imágenes inversas $p_i^*(F')$ de $F'$
sobre $S'\times_SS'$ corresponden asimismo a los haces cuasi-coherentes de
$(\mathcal A'\otimes_{\mathcal O_S}\mathcal A')$-módulos
\[
  F'\otimes_{\mathcal O_S}\mathcal A'
  \quad\text{y}\quad
  \mathcal A'\otimes_{\mathcal O_S}F'.
\]
Dar un $(S'\times_SS')$-homomorfismo del primero al segundo equivale a dar un
homomorfismo de $(\mathcal A'\otimes_{\mathcal O_S}\mathcal A')$-módulos y,
teniendo en cuenta que $\mathcal A'$ es localmente libre, esto equivale a dar
un homomorfismo de $(\mathcal A'\otimes\mathcal A')$-módulos:
\[
  \mathcal U=
  \operatorname{Hom}_{\mathcal O_S}(\mathcal A',\mathcal A')
  =\mathcal A'\otimes\check{\mathcal A}'
  \longrightarrow
  \operatorname{Hom}_{\mathcal O_S}(F',F'),
\]
es decir, a dar, para cada sección $\xi$ de $\mathcal U$ sobre un abierto $V$,
un homomorfismo de $\mathcal O_S$-módulos
$T_\xi:F'|_V\longrightarrow F'|_V$ que satisface las condiciones
\begin{equation}
  T_{f\xi}(x)=fT_\xi(x),
  \qquad T_{\xi f}(x)=T_\xi(fx),
  \tag{3.1}
\end{equation}
donde $f,x$ son, respectivamente, secciones de $\mathcal A',F'$ sobre un
abierto de $S$ contenido en $V$. Las condiciones (i) y (ii) de un dato de
descenso (A, parágrafo 1, (c)) se escriben entonces, respectivamente:
\begin{equation}
  T_{1_{\mathcal U}}x=x,
  \qquad\text{es decir,}\qquad
  T_{1_{\mathcal U}}=\operatorname{id}_{F'},
  \tag{3.2}
\end{equation}
\begin{equation}
  T_{\xi\eta}=T_\xi T_\eta.
  \tag{3.3}
\end{equation}
En otras palabras, \emph{un dato de descenso sobre $F'$ equivale a una
representación del haz de $\mathcal O_S$-álgebras
$\mathcal U=\operatorname{Hom}_{\mathcal O_S}(\mathcal A',\mathcal A')$ en el
haz de $\mathcal O_S$-álgebras
$\operatorname{Hom}_{\mathcal O_S}(F',F')$ que satisface las dos condiciones de
linealidad (3.1).} Si se tiene un apareamiento de haces cuasi-coherentes sobre
$S'$:
\[
  F'_1\times F'_2\longrightarrow F'_3
\]
(que puede interpretarse como un apareamiento de haces de
$\mathcal A'$-módulos sobre $S$),

\src{320}
y datos de pegado sobre los $F'_i$, definidos por homomorfismos
\[
  T_i:\mathcal U\longrightarrow
  \operatorname{Hom}_{\mathcal O_S}(F'_i,F'_i)
  \qquad (i=1,2,3),
\]
entonces estos datos son compatibles con el apareamiento dado, en el sentido
evidente del término, si y solo si se verifica la condición siguiente.

Para toda sección $\xi$ de $\mathcal U$ sobre un abierto, denotando por
\[
  \Delta\xi=\sum_i\xi'_i\otimes_{\mathcal A'}\xi''_i
\]
la sección de $\mathcal U\otimes_{\mathcal A'}\mathcal U$ (considerando
$\mathcal U$ como $\mathcal A'$-módulo mediante su estructura izquierda)
definida por la fórmula
\[
  \xi\mathbin{\cdot}(fg)=\sum_i\xi'_i(f)\xi''_i(g)
\]
(donde $f$ y $g$ son dos secciones de $\mathcal A'$ sobre un abierto más
pequeño), se tiene la fórmula
\begin{equation}
  T^{(3)}_\xi(x\mathbin{\cdot}y)
  =\sum_iT^{(1)}_{\xi'_i}x\mathbin{\cdot}T^{(2)}_{\xi''_i}y
  \tag{3.4}
\end{equation}
para todo par de secciones $x,y$ de $F'_1,F'_2$, respectivamente, sobre un
abierto más pequeño.\footnote{El texto impreso dice «de $\mathcal A'$». Los
dominios de $T^{(1)}$ y $T^{(2)}$, así como el apareamiento
$F'_1\times F'_2\longrightarrow F'_3$, obligan a que $x$ e $y$ sean secciones
de $F'_1$ y $F'_2$, respectivamente.} (Esta propiedad podrá expresarse
diciendo que los homomorfismos $T^{(i)}$ son compatibles con la aplicación
diagonal de $\mathcal U$ respecto del apareamiento dado.) En particular, las
fórmulas (3.1) a (3.4) permiten interpretar, en términos de representaciones de
álgebras con aplicaciones diagonales, los datos de descenso sobre un haz
cuasi-coherente de álgebras sobre $S'$ y, por tanto, también (restringiéndose a
las álgebras conmutativas) los datos de descenso sobre un $S'$-esquema afín.

De aquí se pasa a una interpretación análoga de los datos de descenso sobre un
$S'$-preesquema $X'$ cualquiera: dar tal $X'$ equivale a dar un preesquema
$X'$ sobre $S$, provisto de un homomorfismo de $\mathcal O_S$-álgebras
\[
  \mathcal A'\longrightarrow\mathcal O_{X'},
\]
y un dato de descenso sobre $X'$ equivale a dar un homomorfismo de haces
(compatible con el morfismo $h:X'\longrightarrow S'$)
\[
  \mathcal U\longrightarrow
  \operatorname{Hom}_{h^{-1}(\mathcal O_S)}
  (\mathcal O_{X'},\mathcal O_{X'}),
\]
que satisface condiciones análogas a las condiciones (3.1) a (3.4) anteriores.

\textsc{Ejemplo 1 («Weil»).} -- Supongamos que $S'/S$ es un recubrimiento étale
de Galois de grupo de Galois $\Gamma$ (cf. A, parágrafos 3 y 4.d). Entonces, un
dato de descenso sobre un haz cuasi-coherente $F'$ sobre $S'$ (resp. sobre un
$S'$-preesquema $X'$) equivale a dar una representación de $\Gamma$ por
automorfismos de $(S',F')$ (resp. de $(S',X')$) compatible con las operaciones
de $\Gamma$ sobre $S'$. Este resultado

\src{321}
es «formal», es decir, se demuestra en términos de categorías, pero desde el
punto de vista de este número se deduce también de la estructura explícita de
$\mathcal U$ (provisto de su estructura de anillo, del homomorfismo de anillos
$\mathcal A'\longrightarrow\mathcal U$ y de la aplicación diagonal), conocida
por completo gracias al resultado siguiente: $\mathcal U$ admite, como
$\mathcal A'$-módulo izquierdo, una base formada por las secciones de
$\mathcal U$ que corresponden a los elementos de $\Gamma$.

\textsc{Ejemplo 2 («Cartier»).} -- Sea $p$ un número primo; supongamos que
$p\mathcal O_S=0$ (es decir, $\mathcal O_S$ es de característica $p$),
$(\mathcal A')^p\subset\mathcal O_S=\mathcal A$ (es decir, $S'/S$ es radicial
de altura 1) y que el haz de álgebras $\mathcal A'$ sobre $\mathcal A$ admite
localmente una $p$-base, es decir, una familia $(x_i)$ de secciones tal que
$\mathcal A'$ está generada como álgebra por los $x_i$, sometidos a las únicas
condiciones $x_i^p=0$. Suponemos que el conjunto de los $i$ es finito, de
cardinal $n$. Sea $\mathfrak g$ el haz de las $\mathcal A$-derivaciones de
$\mathcal A'$; es un haz de $\mathcal A'$-módulos localmente libre de rango
$n$ y, además, un haz de $p$-álgebras de Lie sobre $\mathcal A$ (pero no sobre
$\mathcal A'$) que satisface la condición
\begin{equation}
  [X,fY]=X(f)Y+f[X,Y].
  \tag{3.5}
\end{equation}

\textsc{Lema.} -- $\mathcal U=
\operatorname{Hom}_{\mathcal O_S}(\mathcal A',\mathcal A')$ está generada,
como $\mathcal O_S$-álgebra provista de un homomorfismo de álgebras
$\mathcal A'\longrightarrow\mathcal U$, por el sub-$\mathcal A'$-módulo
izquierdo $\mathfrak g$, con las relaciones adicionales
\begin{equation}
  \left\{
  \begin{aligned}
    Xf-fX&=X(f),\\
    XY-YX&=[X,Y],\\
    X^p&=X^{(p)}.
  \end{aligned}
  \right.
  \tag{3.6}
\end{equation}

Del lema anterior se deduce que un dato de descenso sobre el haz
cuasi-coherente $F'$ sobre $S'$ equivale a dar, para todo $X\in\mathfrak g$, un
$\mathcal O_S$-endomorfismo $\overline X$ de $F'$ que satisface las condiciones
\begin{align}
  \overline{fX}&=f\overline X, \tag{3.7}\\
  \overline X(fx)&=X(f)x+f\overline X(x), \tag{3.8}\\
  \overline{[X,Y]}&=[\overline X,\overline Y], \tag{3.9}\\
  \overline{X^{(p)}}&=\overline X^p. \tag{3.10}
\end{align}
(Esto es lo que podría llamarse una conexión lineal sobre $F'$, sin curvatura
y compatible con la potencia $p$-ésima.) Se explicita de igual modo la noción
de dato de descenso sobre un $S'$-preesquema $X'$; la relación (3.4) se
sustituye aquí por la condición de que los $\overline X$ sean derivaciones de
$\mathcal O_{X'}$. Como el morfismo $S'\longrightarrow S$ es radicial, el
teorema 3 garantiza que todo dato de descenso de este tipo es

\src{322}
efectivo y, por tanto, define un $S$-preesquema $X$.

Obsérvese que no ha sido necesario imponer ninguna hipótesis de planitud, de no
singularidad ni de finitud de ningún tipo sobre $F'$ resp. $X'$.

\subsection*{4. Aplicación a criterios de racionalidad}
\addcontentsline{toc}{subsection}{4. Aplicación a criterios de racionalidad}

Sea $X$ un $S$-preesquema tal que la imagen directa de $\mathcal O_X$ sobre
$S$ sea $\mathcal O_S$; esta propiedad seguirá siendo cierta tras todo cambio
plano $S'\longrightarrow S$ de la base $S$. Si $F$ es un haz invertible (es
decir, localmente libre de rango 1) sobre $X$, los automorfismos de $F$, que se
identifican con las secciones invertibles de $\mathcal O_X$, corresponden
biunívocamente a las secciones invertibles de $\mathcal O_S$. Sea entonces $s$
una sección de $X$ sobre $S$; llamaremos, de manera gráfica, sección de $F$
sobre $s$ a una sección del haz invertible $s^*(F)$ sobre $S$. De lo anterior
se deduce que, si $F_i$ ($i=1,2$) son dos haces invertibles sobre $X$, cada uno
provisto de una sección sobre $s$, y si $F_1$ y $F_2$ son isomorfos, existe un
único isomorfismo de $F_1$ sobre $F_2$ compatible con las secciones en cuestión
(es decir, que transforma la primera en la segunda). Por otra parte, e
independientemente de la sección $s$, convengamos en considerar equivalentes
dos haces invertibles $F_1$ y $F_2$ sobre $X$ tales que todo punto de $S$ tenga
un entorno abierto $U$ para el cual las restricciones de $F_1$ y $F_2$ a
$X|U$ sean isomorfas. Entonces todo haz invertible $F$ sobre $X$ es equivalente
a un haz invertible $F_1$ provisto de una sección distinguida sobre $s$ (se
toma $F_1=F\,s^*(F)^{-1}$), y $F_1$ está determinado salvo isomorfismo. En
otras palabras, la clasificación de los haces invertibles sobre $X$ salvo
equivalencia es la misma que la clasificación, salvo isomorfismo, de los haces
invertibles provistos de una sección distinguida.

Como estas propiedades siguen siendo ciertas tras un cambio plano
$\alpha:S'\longrightarrow S$ de la base (reemplazando la sección $s$ por su
imagen inversa $s'$ por $\alpha$), se concluye, teniendo en cuenta el teorema 1:

\emph{Sea el preesquema $X/S$ como arriba y supongamos que admite una sección
$s$; sea $\alpha:S'\longrightarrow S$ un morfismo fielmente plano y
cuasicompacto; sea $F'$ un haz invertible sobre $X'=X\times_SS'$. Para que $F'$
sea equivalente a la imagen inversa sobre $X'$ de un haz invertible $F$ sobre
$X$, es necesario y suficiente que sus imágenes inversas $p_1^*(F')$ y
$p_2^*(F')$ sobre $X'\times_XX'=X\times_S(S'\times_SS')$ sean equivalentes. Si
es así, $F$ está determinado salvo equivalencia.} (Se dirá entonces que $F'$
es racional sobre $S$.)

Inspirándose en este principio cuando $\alpha:S'\longrightarrow S$ es como en
el ejemplo 1 o el ejemplo 2 del número anterior, se obtienen los criterios de
racionalidad de Weil o de Cartier. (Obsérvese que estos autores se limitan al
caso en que $S$ y $S'$

\src{323}
son espectros de cuerpos; a fortiori, $S$ es entonces el espectro de un anillo
local, y la relación de equivalencia introducida arriba no es sino la relación
de isomorfía.) En el primer caso, $F'$ es racional sobre $S$ si y solo si sus
transformados por $\Gamma$ son equivalentes a $F'$. Para expresar el criterio
de racionalidad en el segundo caso, se consideran, en general, el morfismo
diagonal
\[
  X'\longrightarrow X''=X'\times_XX'
\]
de $X'/X$, el haz de ideales $\mathcal I$ correspondiente sobre
$X'\times_XX'$ y el haz $\mathcal I/\mathcal I^2$, que se identifica con su
imagen inversa $\Omega^1_{X'/X}$ sobre $X'$. Como las restricciones de los
$F''_i=p_i^*(F')$ ($i=1,2$) a la diagonal son isomorfas (pues son isomorfas a
$F'$), es decir, $F''_1(F''_2)^{-1}=F''$ tiene una restricción a la diagonal
que es trivial, se sigue que la restricción de $F''$ a
$(X'',\mathcal O_{X''}/\mathcal I^2)$ viene dada, salvo isomorfismo, por un
elemento bien determinado $\xi$ de
\[
  H^1(X'',\mathcal I/\mathcal I^2)
  =H^1(X',\Omega^1_{X'/X}).
\]
Además, en este caso se tiene
\[
  \Omega^1_{X'/X}
  =\Omega^1_{S'/S}\otimes_{\mathcal O_S}\mathcal O_X,
\]
y, por tanto, si $\Omega^1_{S'/S}$ es localmente libre sobre $S$ (como en el
caso de Cartier), $\xi$ define una sección de
\[
  R^1f'_*(\mathcal O_{X'})\otimes\Omega^1_{S'/S}
\]
sobre $S'$ (llamada clase de Atiyah--Cartier del haz invertible $F'$ sobre
$X'/S$), cuya anulación es necesaria y suficiente para que sean equivalentes
las imágenes inversas de $F'$ por las dos proyecciones de
\[
  (X'',\mathcal O_{X''}/\mathcal I^2)
  =X\times_S(S'',\mathcal O_{S''}/\mathcal J^2)
\]
sobre $X'$ (donde $\mathcal J$ es el haz de ideales sobre
$S''=S'\times_SS'$ definido por el morfismo diagonal
$S'\longrightarrow S'\times_SS'$). Esta anulación es, pues, trivialmente
necesaria para que sean equivalentes las imágenes inversas de $F'$ sobre el
propio $X''=X\times_SS''$ y, por tanto, también para que $F'$ sea equivalente
a la imagen inversa de un haz invertible $F$ sobre $X$. Además, la clase de
Atiyah--Cartier puede interpretarse también como la obstrucción a la existencia,
localmente sobre $S'$, de una conexión de $F'$ relativa a las derivaciones de
$X'/X$; tal conexión queda, además, determinada cuando se conocen las
derivaciones de $F'$ correspondientes a las prolongaciones naturales a $X'$ de
las derivaciones de $S'/S$. De esto y de los desarrollos del número anterior se
concluye fácilmente que, en el caso del ejemplo 2 de dicho número y cuando
$X/S$ admite una sección, la anulación de la clase de Atiyah--Cartier es también
suficiente para que $F'$ sea racional sobre $S$.

\subsection*{5. Aplicación a la restricción de la base en un esquema abeliano}
\addcontentsline{toc}{subsection}{5. Restricción de la base en un esquema abeliano}

Sea $S$ un preesquema. Se llama esquema abeliano sobre $S$ a un esquema $X$
liso y propio

\src{324}
sobre $S$ cuyas fibras en los puntos $x\in S$ son esquemas de variedades
abelianas sobre los $\kappa(x)$. Supongamos que $S$ es noetheriano y regular
(es decir, que sus anillos locales son regulares); entonces, utilizando el
teorema de conexidad de Murre [4] (al menos en el caso de «características
iguales», en el que el teorema citado está demostrado actualmente), puede
demostrarse que toda sección racional de $X$ sobre $S$ está definida en todas
partes (es decir, es una sección), lo que generaliza un teorema clásico de
Weil. Se deduce, más generalmente, que si $X'$ es un esquema liso sobre $S$,
entonces toda $S$-aplicación racional de $X'$ en $X$ está definida en todas
partes. De ello se obtiene lo siguiente, que generaliza un resultado de
Chow--Iang:\footnote{El texto impreso dice «Igusa--Lang»; la fe de erratas de
1962 prescribe «Chow--Iang».} Sea $S$ noetheriano y regular, y denote $K$ su
anillo de funciones racionales (producto directo de cuerpos); sea $X$ un
esquema abeliano sobre $K$. Si $X$ es isomorfo a un $K$-esquema de la forma
$X_0\times_S\operatorname{Spec}(K)$, donde $X_0$ es un esquema abeliano sobre
$S$, entonces $X_0$ está determinado salvo isomorfismo único.

Utilizando el resultado de unicidad anterior, se ve que la cuestión de la
restricción de la base en $X$ es local sobre $S$ (y, por tanto, que basta saber
efectuar la restricción a los $\operatorname{Spec}(\mathcal O_x)$, con
$x\in S$). Se ve de igual manera que si
$\alpha:S'\longrightarrow S$ es un morfismo liso de tipo finito, si $K'$ es el
anillo de funciones racionales de $S'$, y si $X\otimes_KK'$ es de la forma
$X'_0\times_{S'}\operatorname{Spec}(K')$, entonces $X'_0$ está provisto de un
dato de descenso canónico relativo a $\alpha$. Teniendo en cuenta el teorema 3,
se concluye:

\result{PROPOSICIÓN}{5.1} Sea $S$ un preesquema noetheriano, regular e
irreducible, de cuerpo de funciones racionales $K$; sea $K'$ una extensión
finita de $K$, no ramificada sobre $S$; sea $S'$ la normalización de $S$ en
$K'$ (que es, por tanto, un recubrimiento étale de $S$); y sea $X$ un esquema
abeliano sobre $K$ tal que $X\otimes_KK'$ sea de la forma
$X'_0\times_{S'}\operatorname{Spec}(K')$, donde $X'_0$ es un esquema abeliano
proyectivo sobre $S'$. Entonces $X$ es de la forma
$X_0\times_S\operatorname{Spec}(K)$, donde $X_0$ es un esquema abeliano
proyectivo sobre $S$.

\result{OBSERVACIONES}{} El conferenciante ignora si puede sustituirse la
hipótesis de que $S'\longrightarrow S$ sea un recubrimiento étale sobreyectivo
(que permite utilizar el teorema 3) por la hipótesis de que sea un morfismo
liso, sobreyectivo y de tipo finito (incluso si se supone que es un morfismo
étale), o si la proposición sigue siendo válida sin suponer que $X'_0$ sea
proyectivo sobre $S'$ (condición que tal vez se cumpla automáticamente).

\subsection*{6. Aplicación a criterios de trivialidad local e isotrivialidad}
\addcontentsline{toc}{subsection}{6. Criterios de trivialidad local e isotrivialidad}

Sean $S$ un preesquema, $G$ un «preesquema en grupos» sobre $S$, y $P$ un
preesquema sobre $S$ sobre el que «actúa $G$» (por la derecha). Se dice que $P$
es formalmente principal homogéneo bajo $G$ si el morfismo bien conocido
\[
  G\times_SP\longrightarrow P\times_SP
\]

\src{325}
inducido por la acción de $G$ sobre $P$ es un isomorfismo. En lo sucesivo,
supondremos que $G$ es plano sobre $S$ (y, por tanto, fielmente plano sobre
$S$), y reservaremos el nombre de fibrado principal homogéneo bajo $G$ para un
fibrado formalmente principal homogéneo $P$ que sea fielmente plano y
cuasicompacto sobre $S$. Es inmediato que esto equivale a decir que puede
encontrarse un cambio fielmente plano y cuasicompacto
$S'\longrightarrow S$ de la base $S$, tal que el fibrado formalmente principal
homogéneo
\[
  P'=P\times_SS'
\]
bajo $G'=G\times_SS'$ sea trivial, es decir, isomorfo a $G'$ (es decir, que
admita una sección); en particular, se podrá tomar $S'=P$. Obsérvese además
que, si $S$ es localmente noetheriano, la hipótesis de planitud fiel sobre $P$
equivale a la hipótesis de que
\[
  \widehat P_s=P\times_S\operatorname{Spec}(\widehat{\mathcal O}_s)
\]
sea fielmente plano sobre $\widehat{\mathcal O}_s$ para todo $s\in S$ (donde
$\widehat{\mathcal O}_s$ denota el completado del anillo local
$\mathcal O_s$), como se deduce del hecho de que $\widehat{\mathcal O}_s$ es
fielmente plano sobre $\mathcal O_s$. Además, si $P$ es de tipo finito sobre
$S$ localmente noetheriano, el conjunto de los puntos $s$ que satisfacen la
condición anterior es constructible; por tanto, si $S$ es un «preesquema de
Jacobson» (por ejemplo, un esquema de tipo finito sobre un cuerpo o, más
generalmente, el espectro de un anillo de Jacobson), basta comprobar la
condición en cuestión en los puntos cerrados de $S$. Esto nos reduce al caso en
que la base es el espectro de un anillo local completo $A$. Cuando
$S=\operatorname{Spec}(A)$ ($A$ un anillo local noetheriano completo) y $P$ es
de tipo finito sobre $S$, la planitud fiel de $P/S$ equivale también a la
existencia de un $S'$ finito y plano sobre $S$ tal que $P'$ sea trivial; y, si
además $G$ es liso sobre $S$, puede suponerse que $S'$ es étale sobre $S$. Por
consiguiente, si además el cuerpo residual de $A$ es algebraicamente cerrado
(«caso geométrico»), $P$ es fielmente plano sobre $A$ si y solo si es trivial.
Así pues, si $S$ es un preesquema algebraico sobre un cuerpo algebraicamente
cerrado y $G$ es liso y de tipo finito sobre $S$, se ve que la condición de
planitud fiel sobre $S$ equivale a la condición de trivialidad analítica (SLF)
de Serre ([6], pp. 1--12).

Pueden introducirse otros tipos, más fuertes, de condiciones sobre $P$ que
tengan la naturaleza de una «trivialidad local». Diremos, en particular, que
$P$ es isotrivial (resp. estrictamente isotrivial) si, para todo $s\in S$,
existen un entorno abierto $U$ de $s$ y un morfismo finito y fielmente plano
(resp. un recubrimiento étale sobreyectivo) $U'\longrightarrow U$ tales que
$P'=P\times_SU'$ sea trivial. (Nuestra terminología difiere de la de Serre [6],
quien llama localmente isotrivial a lo que nosotros llamamos estrictamente
isotrivial.) La isotrivialidad estricta es especialmente útil si $G$ es liso
sobre $S$, pero, en cambio, es una noción inadecuada en los demás casos.

Si $G$ es afín sobre $S$, todo fibrado principal homogéneo $P$ bajo $G$ es afín
por el parágrafo 2, de donde, gracias al teorema 2, resulta la posibilidad de
«descender»

\src{326}
tales fibrados por morfismos fielmente planos y cuasicompactos. Tomemos, en
particular, $G=\operatorname{Gl}(n)_S$, definido por la condición de que el
funtor de $S$-preesquemas
\[
  S'\longmapsto\operatorname{Hom}_S(S',G)
\]
(con valores en la categoría de los grupos) se identifique con el funtor
\[
  \operatorname{Gl}(n)(S')
  =\operatorname{Gl}\bigl(n,H^0(S',\mathcal O_{S'})\bigr)
\]
de A, parágrafo 4.e. Utilizando los hechos (i) de que todo fibrado principal
homogéneo bajo $G$ (resp. todo haz localmente libre de rango $n$ sobre $S$) se
vuelve isomorfo al objeto «trivial» $G$ (resp. $\mathcal O_S^n$) tras un cambio
fielmente plano y cuasicompacto adecuado de $S$; (ii) de que pueden descenderse
los objetos del tipo considerado (fibrados principales homogéneos bajo $G$,
resp. haces localmente libres de rango $n$) por tales morfismos; y, por último,
(iii) de que el grupo de automorfismos del fibrado trivial sobre un $S'/S$ es
funtorialmente isomorfo al grupo de automorfismos del haz localmente libre
trivial de rango $n$ sobre $S'$, se concluye «formalmente» que «da lo mismo»
dar sobre $S$ (o sobre un $S'/S$) un fibrado principal homogéneo de grupo $G$
que dar allí un haz localmente libre de rango $n$. (Más precisamente, se tiene
una equivalencia de categorías fibradas.) Se concluye, en particular:

\result{PROPOSICIÓN}{6.1} Todo fibrado principal homogéneo de grupo
$\operatorname{Gl}(n)_S$ es localmente trivial.

Mediante argumentos conocidos se obtiene el mismo resultado para grupos
estructurales tales como $\operatorname{Sl}(n)_S$, $\operatorname{Sp}(n)_S$ y
productos de tales grupos. Se deduce también que, si $F$ es un subgrupo cerrado
de $G=\operatorname{Gl}(n)_S$, plano sobre $S$, tal que existe el cociente
$G/F$, y si $G$ es un fibrado principal homogéneo isotrivial (resp.
estrictamente isotrivial) sobre $G/F$, de grupo $F\times_S(G/F)$, entonces todo
fibrado principal homogéneo de grupo $F$ es isotrivial (resp. estrictamente
isotrivial). Esto se aplica a todos los «grupos lineales» sobre $S$ utilizados
hasta ahora y, en particular, cuando $G=S\times_k\Gamma$, donde $S$ es un
preesquema sobre el cuerpo $k$ y $\Gamma$ un grupo lineal en el sentido clásico
y, en particular, liso, sobre $k$. Esto resuelve, pues, para tales grupos, una
cuestión de Serre (loc. cit.).

Señalemos asimismo que, para la mayoría de los grupos lisos sobre $S$ conocidos,
sean lineales o no, y en cualquier caso para los de la forma
$S\times_k\Gamma$ anterior, puede demostrarse que todo fibrado principal
homogéneo isotrivial es estrictamente isotrivial. Esto resuelve, en particular,
otra cuestión de Serre (loc. cit., pp. 1--14), dado que un fibrado principal
homogéneo obtenido por un descenso a la Cartier (cf. parágrafo 3, ejemplo 2) es
manifiestamente isotrivial.

\result{OBSERVACIÓN}{} Una de las dificultades esenciales en estas cuestiones
(aparte de la cuestión de la existencia de los esquemas cociente $G/F$) es la
falta de criterios de efectividad para un dato de descenso por un morfismo
fielmente plano no finito.

\src{327}
\begin{center}
  \textbf{BIBLIOGRAFÍA}
\end{center}

\begin{description}
  \item[{[1]}] Dieudonné (J.) y Grothendieck (Alexander). --- \emph{Éléments
    de géométrie algébrique}, de próxima aparición en las Publications
    mathématiques de l'Institut des Hautes Études Scientifiques.
  \item[{[2]}] Grauert (H.) und Remmert (R.). --- \emph{Komplexe Räume}, Math.
    Annalen, t. 136, 1958, pp. 245--318.
  \item[{[3]}] Grothendieck (Alexander). --- \emph{Géométrie formelle et
    géométrie algébrique}. Séminaire Bourbaki, t. 11, 1958/59, no 182.
  \item[{[4]}] Murre (J. P.). --- \emph{On a connectedness theorem for a
    birational transformation at a simple point}, Amer. J. Math., t. 80, 1958,
    pp. 3--15.
  \item[{[5]}] Serre (Jean-Pierre). --- \emph{Géométrie algébrique et géométrie
    analytique}, Ann. Inst. Fourier Grenoble, t. 6, 1955--56, pp. 1--42.
  \item[{[6]}] Serre (Jean-Pierre). --- \emph{Espaces fibrés algébriques},
    Séminaire Chevalley, t. 2, 1958: Anneaux de Chow et applications, no 1.
\end{description}

\def\FgaUnitKey{e190}\clearpage
\chapter*{Fe de erratas del Exposé n\textsuperscript{o} 190}
\addcontentsline{toc}{chapter}{Fe de erratas del Exposé n\textsuperscript{o} 190}
\setcounter{footnote}{0}
\src{300}

\noindent\textbf{3. Técnica de descenso y teoremas de existencia en geometría
algebraica [TDTE I]. I: Generalidades. Descenso por morfismos fielmente planos.}

\begin{center}
[Séminaire Bourbaki, t. 12, 1959/60, n\textsuperscript{o} 190, 29 pp.]
\end{center}

\noindent\textbf{Página 1, línea 10.} --- Ahora parece excesivo decir que la
técnica de descenso está «en la base de la mayoría de los teoremas de
existencia en geometría algebraica». Esto es cierto en gran medida para las
técnicas no proyectivas que constituyen el objeto de los dos primeros exposés
de la serie TDTE I a VI, pero no para las técnicas proyectivas (TDTE IV, V,
VI).

\noindent\textbf{Página 5, línea 16.} --- No es necesario suponer que $\alpha$
sea un morfismo de $\mathcal F$-descenso.

\noindent\textbf{Página 20, observación.} --- Un morfismo
$S'\longrightarrow S$ cuasifinito, étale y sobreyectivo, o un morfismo
$\operatorname{Spec}(\widehat A)\longrightarrow\operatorname{Spec}(A)$, no es
siempre un morfismo de descenso estricto, aun si $A$ es un anillo local de una
curva algebraica sobre un cuerpo algebraicamente cerrado $k$ y
$S=\operatorname{Spec}(A)$. Así, se puede encontrar un morfismo propio y liso
$f:X\longrightarrow S$, que hace de $X$ un fibrado principal de base $S$ bajo
una curva elíptica $E$, tal que $f':X'\longrightarrow S'$ sea proyectivo, pero
$f$ no lo sea. Es, por tanto, al mismo tiempo un ejemplo de un fibrado
homogéneo principal no isotrivial bajo un esquema abeliano.

\noindent\textbf{Página 26, línea 9.} --- Léase «Chow--Lang» en lugar de
«Igusa--Lang».

Para diversos detalles acerca de la teoría del descenso, véanse SGA VI, VII y
VIII.

\def\FgaUnitKey{195}\unit{Exposición 195 -- Técnica de descenso II}
\sid{195}
\src{369}
\seminar{12.\textsuperscript{o} año, 1959/60, n.\textsuperscript{o} 195\\Febrero de 1960}
\paperhead{TÉCNICA DE DESCENSO Y TEOREMAS DE EXISTENCIA EN GEOMETRÍA ALGEBRAICA\\
II. EL TEOREMA DE EXISTENCIA EN LA TEORÍA FORMAL DE LOS MÓDULOS}%
  {por Alexander GROTHENDIECK}

\fsection{A}{Funtores representables y pro-representables}

\subsection*{1. Funtores representables}
\addcontentsline{toc}{subsection}{1. Funtores representables}

Sea $\mathcal C$ una categoría. Para todo $X\in\mathcal C$, sea $h_X$ el
funtor contravariante de $\mathcal C$ en la categoría $(\operatorname{Ens})$
de los conjuntos,
\[
  h_X:\mathcal C\longrightarrow(\operatorname{Ens}),
\]
definido por la fórmula
\[
  h_X(Y)=\operatorname{Hom}(Y,X).
\]
Un morfismo $X\longrightarrow X'$ en $\mathcal C$ induce de manera evidente un
homomorfismo de funtores $h_X\longrightarrow h_{X'}$: $h_X$ es un funtor
covariante en $X$; es decir, hemos definido un funtor covariante canónico
\[
  h:\mathcal C\longrightarrow
  \operatorname{Hom}\bigl(\mathcal C^o,(\operatorname{Ens})\bigr)
\]
de $\mathcal C$ en la categoría de los funtores covariantes de la dual
$\mathcal C^o$ de $\mathcal C$ en la categoría de los conjuntos. Recordemos
entonces:

\result{PROPOSICIÓN}{1.1} Este funtor $h$ es \emph{plenamente fiel}; en otros
términos, para todo par $X,X'$ de objetos de $\mathcal C$, la aplicación
natural
\[
  \operatorname{Hom}(X,X')\longrightarrow
  \operatorname{Hom}(h_X,h_{X'})
\]
es biyectiva.

En particular, si un funtor
$F\in\operatorname{Hom}(\mathcal C^o,(\operatorname{Ens}))$ es isomorfo a un
funtor de la forma $h_X$, entonces $X$ queda determinado salvo isomorfismo
único. Se dice en tal caso que el funtor $F$ es \emph{representable}. La
proposición anterior significa que el funtor canónico $h$ define una
equivalencia de la categoría $\mathcal C$ con la subcategoría plena de
$\operatorname{Hom}(\mathcal C^o,(\operatorname{Ens}))$ formada por los
funtores representables. Este hecho fundamenta la noción de «solución de un
problema universal»: tal problema consiste siempre en examinar si un funtor
dado (contravariante como aquí, o covariante en el caso dual) de $\mathcal C$
en $(\operatorname{Ens})$ es representable.

\src{370}
Observemos además que, por la propia definición de producto en una categoría
[1], el funtor $h:X\longmapsto h_X$ conmuta con los productos siempre que estos
existan (y, más generalmente, con los límites proyectivos finitos o infinitos;
en particular, con los productos fibrados, con la formación de «núcleos» [2],
etc., siempre que existan): se tiene un isomorfismo de funtores
\[
  h_{X\times X'}\xrightarrow{\sim}h_X\times h_{X'}
\]
siempre que exista $X\times X'$; es decir, se tienen biyecciones functoriales
en $Y$
\[
  h_{X\times X'}(Y)\xrightarrow{\sim}h_X(Y)\times h_{X'}(Y).
\]
En particular, dar un morfismo
\[
  X\times X'\longrightarrow X''
\]
en $\mathcal C$ (es decir, una «ley de composición» en $\mathcal C$ entre
$X,X',X''$) equivale a dar un morfismo
\[
  h_{X\times X'}=h_X\times h_{X'}\longrightarrow h_{X''},
\]
es decir, a dar, para todo $Y\in\mathcal C$, una ley de composición de
conjuntos
\[
  h_X(Y)\times h_{X'}(Y)\longrightarrow h_{X''}(Y)
\]
de tal modo que, para todo morfismo $Y\longrightarrow Y'$ en $\mathcal C$, el
sistema de aplicaciones entre conjuntos
\[
  h_{X^{(i)}}(Y)\longleftarrow h_{X^{(i)}}(Y')
  \qquad(\text{para }i=0,1,2)
\]
sea un morfismo respecto de las dos leyes de composición correspondientes a
$Y$ y $Y'$. Se ve así que la noción de estructura de «$\mathcal C$-grupo»,
«$\mathcal C$-anillo», etc. sobre un objeto $X$ de $\mathcal C$ se expresa del
modo más cómodo (tanto en la teoría como en la práctica) diciendo que, para
todo $Y\in\mathcal C$, el conjunto $h_X(Y)$ está provisto de una ley de grupo
(resp. de anillo, etc.) en el sentido usual, y que las aplicaciones
$h_X(Y)\longleftarrow h_X(Y')$ correspondientes a morfismos
$Y\longrightarrow Y'$ son homomorfismos de grupos (resp. de anillos, etc.).
Esta es, por ejemplo, la forma más intuitiva y cómoda de definir los diversos
grupos clásicos $G_a$, $G_m$, $\operatorname{Gl}(n)$, etc. sobre un preesquema
base arbitrario $S$ y de escribir las relaciones clásicas entre esos grupos;
de dotar al esquema afín $V(\mathcal F)$ sobre $S$, definido por un haz
cuasicoherente $\mathcal F$, de una estructura de «fibrado vectorial»; y de
definir y estudiar las diversas variedades de banderas asociadas
(grassmannianas, fibrados proyectivos), etc. El yoga general consiste
sencillamente en identificar, mediante el funtor canónico $h$, los objetos de
$\mathcal C$ con ciertos funtores contravariantes: los funtores representables,

\src{371}
de $\mathcal C$ en la categoría de los conjuntos.

El procedimiento habitual de invertir las flechas, necesario por ejemplo en el
caso de los esquemas afines para pasar del lenguaje geométrico al del álgebra
conmutativa, conduce a dualizar las consideraciones anteriores y, en
particular, a introducir también los funtores covariantes representables
$\mathcal C\longrightarrow(\operatorname{Ens})$; es decir, los de la forma
\[
  Y\longmapsto\operatorname{Hom}(X,Y)=h'_X(Y).
\]

\subsection*{2. Funtores pro-representables, proobjetos}
\addcontentsline{toc}{subsection}{2. Funtores pro-representables, proobjetos}

Sea $\widetilde X=(X_i)_{i\in I}$ un sistema proyectivo de objetos de
$\mathcal C$; le corresponde un funtor covariante
\[
  h'_{\widetilde X}=\varinjlim_i h'_{X_i},
\]
más explícitamente,
\[
  h'_{\widetilde X}(Y)=\varinjlim_i h'_{X_i}(Y)
  =\varinjlim_i\operatorname{Hom}(X_i,Y),
\]
de $\mathcal C$ en $(\operatorname{Ens})$. Un funtor de $\mathcal C$ en
$(\operatorname{Ens})$ que sea isomorfo a un funtor de este tipo con $I$
filtrante se llama \emph{pro-representable}. Por el número anterior, estos son
exactamente los funtores isomorfos a límites inductivos filtrantes de funtores
representables. Sea $\widetilde X'=(X'_j)_{j\in J}$ un segundo sistema
proyectivo filtrante en $\mathcal C$ (construido sobre un segundo conjunto de
índices preordenado filtrante $J$); se comprueba fácilmente que hay una
biyección canónica (que generaliza la proposición 1.1)
\[
  \operatorname{Hom}(h'_{\widetilde X'},h'_{\widetilde X})
  =\varprojlim_j\varinjlim_i\operatorname{Hom}(X_i,X'_j).
\]
Esto lleva a introducir la categoría $\operatorname{Pro}(\mathcal C)$ de los
proobjetos de $\mathcal C$; sus objetos son los sistemas proyectivos de objetos
de $\mathcal C$ (sobre conjuntos variables de índices preordenados filtrantes),
y si $\widetilde X=(X_i)_{i\in I}$ y
$\widetilde X'=(X'_j)_{j\in J}$ son dos objetos de esta clase, se pone
\[
  \operatorname{ProHom}(\widetilde X,\widetilde X')
  =\varprojlim_j\varinjlim_i\operatorname{Hom}(X_i,X'_j),
\]
siendo evidente la composición de los prohomomorfismos. Por su propia
construcción, $\widetilde X\longmapsto h'_{\widetilde X}$ puede considerarse
como un funtor contravariante en $\widetilde X$, que establece una equivalencia
entre la categoría dual de la categoría $\operatorname{Pro}(\mathcal C)$ de
los proobjetos de $\mathcal C$ y la categoría de los funtores covariantes
pro-representables de $\mathcal C$ en $(\operatorname{Ens})$. Naturalmente, un
objeto $X$ de $\mathcal C$ define canónicamente un proobjeto, denotado

\src{372}
también por $X$, de modo que $\mathcal C$ es equivalente a una subcategoría
plena de $\operatorname{Pro}(\mathcal C)$. Si entonces
$\widetilde X=(X_i)_{i\in I}$ es un proobjeto cualquiera de $\mathcal C$, se
tendrá (con la identificación anterior)
\[
  \widetilde X=\varprojlim_i X_i,
\]
tomándose el límite proyectivo en $\operatorname{Pro}(\mathcal C)$
(puesto que $h'_{\widetilde X}=\varinjlim_i h'_{X_i}$). Hay que observar que,
aun cuando el límite proyectivo de los $X_i$ exista en $\mathcal C$, en
general no será isomorfo al límite proyectivo $\widetilde X$ en
$\operatorname{Pro}(\mathcal C)$, como ya resulta evidente cuando $\mathcal C$
es la categoría de los conjuntos. Nótese que, por definición,
$\varprojlim_{\mathcal C}X_i=L$ queda definido por la condición de que el
funtor en $Y\in\mathcal C$
\[
  \varprojlim_i\operatorname{Hom}_{\mathcal C}(Y,X_i)
  =\operatorname{Hom}_{\operatorname{Pro}(\mathcal C)}(Y,\widetilde X)
\]
con valores en $(\operatorname{Ens})$ sea representable mediante $L$, es decir,
sea isomorfo a $\operatorname{Hom}_{\mathcal C}(Y,L)$; por consiguiente,
$\varprojlim_{\mathcal C}X_i$ ya está definido en términos del proobjeto
$\widetilde X$ y, más precisamente, depende functorialmente del proobjeto
$\widetilde X$ siempre que esté definido; no hay, pues, inconveniente en
denotarlo por $\varprojlim_{\mathcal C}(\widetilde X)$. Si los límites
proyectivos existen siempre en $\mathcal C$,
$\varprojlim_{\mathcal C}(\widetilde X)$ es un funtor de
$\operatorname{Pro}(\mathcal C)$ en $\mathcal C$, y hay un homomorfismo
canónico de funtores
$\varprojlim_{\mathcal C}(\widetilde X)\longrightarrow\widetilde X$.

Como todo funtor, que supondremos covariante para fijar las ideas,
\[
  F:\mathcal C\longrightarrow\mathcal C'
\]
se prolonga de manera evidente a un funtor
\[
  \operatorname{Pro}(F):\operatorname{Pro}(\mathcal C)
  \longrightarrow\operatorname{Pro}(\mathcal C'),
\]
se sigue que, si en $\mathcal C'$ existen siempre los límites proyectivos, $F$
define también canónicamente un funtor compuesto
\[
  \overline F:\operatorname{Pro}(\mathcal C)\longrightarrow\mathcal C',
  \qquad
  \overline F=\varprojlim_{\mathcal C'}\circ\operatorname{Pro}(F),
\]
que transforma $\widetilde X=(X_i)_{i\in I}$ en
$\varprojlim_{\mathcal C'}F(X_i)$.

Un proobjeto $\widetilde X$ se llama \emph{proobjeto estricto} si es isomorfo a
un proobjeto $(X_i)_{i\in I}$ cuyos morfismos de transición
$X_i\longrightarrow X_j$ son epimorfismos; un funtor definido por tal objeto se
llama \emph{estrictamente pro-representable}. Se puede exigir además que $I$
sea ordenado filtrante y que todo epimorfismo $X_i\longrightarrow X'$ sea
equivalente a un epimorfismo

\src{373}
$X_i\longrightarrow X_j$ para un $j\in I$ adecuado (determinado de manera
única por esta condición). En estas condiciones, el sistema proyectivo
$(X_i)_{i\in I}$ queda determinado salvo isomorfismo único (en el sentido usual
de los isomorfismos de sistemas proyectivos). Se sigue que en
$\operatorname{Pro}(\mathcal C)$ existe siempre el límite proyectivo de un
sistema proyectivo de proobjetos estrictos $\widetilde X^{(\alpha)}$ y que, con
las notaciones anteriores para $F$, $\overline F$, se tendrá
\[
  \overline F\left(\varprojlim_\alpha\widetilde X^{(\alpha)}\right)
  =\varprojlim_{\alpha,\mathcal C'}
  \overline F\left(\widetilde X^{(\alpha)}\right).
\]
En particular, si todo proobjeto de $\mathcal C$ es estricto (cf. el número
siguiente), el funtor prolongado $\overline F$ conmuta con los límites
proyectivos.

\subsection*{3. Caracterización de los funtores pro-representables}
\addcontentsline{toc}{subsection}{3. Caracterización de los funtores pro-representables}

Sean $\mathcal C$, $\mathcal C'$ dos categorías en las que existen los límites
proyectivos finitos (es decir, relativos a conjuntos preordenados finitos, no
necesariamente filtrantes), o, lo que es equivalente, en las que existen los
productos finitos y los productos fibrados finitos (lo cual implica, en
particular, la existencia de un «objeto final» $e$, tal que
$\operatorname{Hom}(X,e)$ se reduzca a un elemento para todo $X$). Sea $F$ un
funtor covariante de $\mathcal C$ en $\mathcal C'$. Las condiciones siguientes
son equivalentes:

\begin{enumerate}
\item[(i)] $F$ conmuta con los límites proyectivos finitos;
\item[(ii)] $F$ conmuta con los productos finitos y los productos fibrados finitos;
\item[(iii)] $F$ conmuta con los productos finitos y, para todo diagrama exacto
\[
  X\longrightarrow X'\rightrightarrows X''
\]
en $\mathcal C$ ([3], A, definición 2.1), el diagrama transformado por $F$
\[
  F(X)\longrightarrow F(X')\rightrightarrows F(X'')
\]
es exacto.
\end{enumerate}

Se dice entonces que $F$ es \emph{exacto por la izquierda}.

En lo sucesivo se supone que en $\mathcal C$ existen los límites proyectivos
finitos. De las definiciones resulta entonces inmediato que un funtor
representable es exacto por la izquierda y, pasando al límite, que un funtor
pro-representable es exacto por la izquierda.

Para obtener un recíproco, sea
\[
  F:\mathcal C\longrightarrow(\operatorname{Ens})
\]
un funtor covariante, y sean $X\in\mathcal C$ y $\xi\in F(X)$. Se dice que
$\xi$ (o el par \src{374}$(X,\xi)$) es \emph{minimal} si, para todo
$X'\in\mathcal C$ y $\xi'\in F(X')$, y todo monomorfismo estricto ([3], A, 2)
$u:X'\longrightarrow X$ tal que $\xi=F(u)(\xi')$, $u$ es un isomorfismo. Por
otra parte, se dice que un par $(X,\xi)$ \emph{domina} a $(X'',\xi'')$
($\xi\in F(X)$, $\xi''\in F(X'')$) si existe un morfismo
$v:X\longrightarrow X''$ tal que $\xi''=F(v)(\xi)$; si $\xi$ es minimal y $F$
es exacto por la izquierda, este morfismo $v$ es entonces único; si $\xi''$ es
minimal, $v$ es sobreyectivo. Se deduce fácilmente la proposición siguiente:

\result{PROPOSICIÓN}{3.1} Para que $F$ sea estrictamente pro-representable, es
necesario y suficiente que satisfaga las dos condiciones:
\begin{enumerate}
\item[(i)] $F$ es exacto por la izquierda;
\item[(ii)] todo par $(X,\xi)$, con $\xi\in F(X)$, está dominado por un par
minimal.
\end{enumerate}

Esta última condición es vacía si todo objeto de $\mathcal C$ es artiniano
(tomando un subobjeto $X'$ de $X$ minimal entre aquellos para los cuales existe
$\xi'\in F(X')$ tal que $\xi$ sea la imagen de $\xi'$). De aquí:

\result{COROLARIO}{} Sea $\mathcal C$ una categoría cuyos objetos son
artinianos y en la que existen los límites proyectivos finitos. Entonces los
funtores pro-representables de $\mathcal C$ en $(\operatorname{Ens})$ son
exactamente los funtores exactos por la izquierda, e incluso son estrictamente
pro-representables.

Este último hecho significa también que \emph{todo proobjeto de $\mathcal C$
es estricto}.

\subsection*{4. Ejemplo: grupos de tipo galoisiano, grupos proalgebraicos}
\addcontentsline{toc}{subsection}{4. Grupos de tipo galoisiano, grupos proalgebraicos}

Si $\mathcal C$ es la categoría de los grupos finitos ordinarios,
$\operatorname{Pro}(\mathcal C)$ es equivalente a la categoría de los grupos
topológicos compactos totalmente disconexos. Son grupos de este tipo y sus
generalizaciones, obtenidas sustituyendo los grupos finitos ordinarios por
esquemas en grupos finitos sobre un preesquema base dado (por ejemplo, los
grupos algebraicos finitos sobre un cuerpo $k$), los que harán las veces de
grupos fundamentales, de homotopía y de homología absolutos y relativos para
los preesquemas. En todos estos ejemplos se aplica el corolario de la
proposición 3.1, y es precisamente mediante el funtor asociado como debe
definirse $\pi_1$ [2]. Lo mismo sucede partiendo de la categoría de los grupos
algebraicos o cuasialgebraicos sobre un cuerpo (o, más generalmente, sobre un
preesquema noetheriano): se obtienen los «grupos proalgebraicos» de Serre [4].

\subsection*{5. Ejemplo: «variedades formales»}
\addcontentsline{toc}{subsection}{5. Ejemplo: variedades formales}

Sean $\Lambda$ un anillo noetheriano y $\mathcal C$ la categoría de las
$\Lambda$-álgebras $A$ que son módulos artinianos de tipo finito sobre
$\Lambda$ (o, más brevemente,

\src{375}
las $\Lambda$-álgebras artinianas). Se cumplen las condiciones del corolario de
la proposición 3.1. La categoría $\operatorname{Pro}(\mathcal C)$ es aquí
equivalente a la categoría de las álgebras topológicas $O$ sobre $\Lambda$
isomorfas a límites proyectivos topológicos
\[
  O=\varprojlim_i O_i
\]
de álgebras $O_i\in\mathcal C$, es decir, cuya topología es lineal, separada y
completa, y tal que para todo ideal abierto $J_i$ de $O$, $O/J_i$ sea un
álgebra artiniana sobre $\Lambda$. El funtor
$\mathcal C\longrightarrow(\operatorname{Ens})$ asociado a tal álgebra no es
otro que
\begin{align*}
  F(A)=h'_O(A)
  &=\{\text{homomorfismos continuos de $\Lambda$-álgebras topológicas }
       O\longrightarrow A\}\\
  &=\varinjlim_i
    \operatorname{Hom}_{\Lambda\text{-álgebras}}(O_i,A).
\end{align*}

Obsérvese, además, que la categoría $\mathcal C$ es esencialmente el producto
de las categorías análogas correspondientes a las completaciones de los
anillos locales $\Lambda_{\mathfrak m}$, para los ideales maximales
$\mathfrak m$ de $\Lambda$; por tanto, si se desea, se puede restringir al caso
en que $\Lambda$ sea un anillo local completo de esta clase. En cualquier caso,
$O$ se descompone canónicamente en el producto topológico de componentes
locales, correspondientes a los «puntos» del esquema formal [2] definido por
$O$. Un punto de esta clase está definido por un elemento $\xi$ de algún
$F(K)$, donde $K\in\mathcal C$ es un cuerpo (por ejemplo, el cuerpo residual
de la componente local considerada); dos pares $(\xi,K)$ y $(\xi',K')$
definen el mismo punto si y solo si están dominados por un mismo
$(\xi'',K'')$, o bien si dominan a un mismo $(\xi''',K''')$. (Si los
$\Lambda/\mathfrak m$ son algebraicamente cerrados, basta tomar la suma
disjunta de los $F(\Lambda/\mathfrak m)$.)

Es importante dar condiciones para que la componente local $O_\xi$ de $O$
correspondiente a un $\xi\in F(K)$ sea un anillo noetheriano. Cuando $\Lambda$
es un anillo local completo (noetheriano, recordémoslo), esto equivale a decir
que $O_\xi$ es isomorfo a un anillo cociente de un anillo de series formales
$\Lambda[[t_1,\ldots,t_n]]$ sobre $\Lambda$. Para dar un criterio de esta clase,
introduzcamos (para todo anillo $A$) la $A$-álgebra $I_A$ de los «números
duales» de $A$ mediante
\[
  I_A=A[t]/t^2A[t].
\]
Sea $\varepsilon:I_A\longrightarrow A$ el homomorfismo de aumento; este define
(si $A\in\mathcal C$) una aplicación
\[
  F(\varepsilon):F(I_A)\longrightarrow F(A),
\]

\src{376}
y, utilizando que $F$ es exacto por la izquierda, se define de manera
intrínseca una estructura de $A$-módulo en la parte
\[
  F(I_A,\xi)=F(\varepsilon)^{-1}(\xi)
\]
de $F(I_A)$ formada por los $\xi'\in F(I_A)$ que «se reducen a $\xi$»;
utilizando la forma explícita de $F$ en términos de $O$, se encuentra que este
$A$-módulo se identifica con
\[
  \operatorname{Hom}_{\Lambda}(\mathfrak m_\xi/\mathfrak m_\xi^2,A),
\]
donde $\mathfrak m_\xi$ es el núcleo del homomorfismo
$\xi:O\longrightarrow A$, es decir, si $A$ es un cuerpo, el ideal maximal de
la componente local $O_\xi$ de $O$. Se deduce inmediatamente la proposición
siguiente:

\result{PROPOSICIÓN}{5.1} Sea $\xi\in F(K)$, donde $K\in\mathcal C$ es un
cuerpo. Para que la componente local correspondiente $O_\xi$ de $O$ sea un
anillo noetheriano, es necesario y suficiente que el conjunto $F(I_K,\xi)$ de
los elementos de $F(I_K)$ que se reducen a $\xi$ sea un espacio vectorial
de dimensión finita sobre $K$. En estas condiciones, se tiene un isomorfismo
canónico
\[
  F(I_K,\xi)=
  \operatorname{Hom}\!\left(
    \frac{\mathfrak m_\xi/\mathfrak m_\xi^2}
         {\operatorname{im}(\mathfrak n_\xi/\mathfrak n_\xi^2)},K
  \right),
\]
donde $\mathfrak n_\xi$ es el ideal maximal de $\Lambda$, núcleo del
homomorfismo $\Lambda\longrightarrow K$.
\footnote{El texto impreso da aquí una fórmula que solo es correcta cuando
$\Lambda$ es un cuerpo. La fe de erratas de 1962 prescribe, en el caso general,
el cociente de $\mathfrak m_\xi/\mathfrak m_\xi^2$ por la imagen de
$\mathfrak n_\xi/\mathfrak n_\xi^2$.}
En particular, la dimensión del $K$-espacio vectorial $F(I_K,\xi)$ es igual a
la dimensión del espacio vectorial así corregido sobre el cuerpo
$O_\xi/\mathfrak m_\xi=K(\xi)$.

Supongamos que $O_\xi$ es noetheriano y, para simplificar la escritura, que
$\Lambda$ es local completo y $O=O_\xi$. Se dice que $O$ es \emph{liso sobre
$\Lambda$} si $O$ es un álgebra finita y étale sobre el álgebra completada del
localizado de $\Lambda[t_1,\ldots,t_n]$ en un ideal maximal de este anillo que
induce el ideal maximal de $\Lambda$; cuando la extensión residual de $O$ sobre
$\Lambda$ es trivial (por ejemplo, si el cuerpo residual de $\Lambda$ es
algebraicamente cerrado), esto equivale a decir que $O$ mismo es isomorfo a un
álgebra de series formales de esta clase.
\footnote{La fe de erratas de 1962 precisa que esta definición de «liso sobre
$\Lambda$» solo es correcta cuando la extensión residual $k'/k$ es separable;
en el caso general, remite a SGA III, 1.1.}
Finalmente, si ya no se supone necesariamente que $O$ sea noetheriano, se dirá
también que $O$ es liso sobre $\Lambda$ si $O$ es isomorfo a un límite
proyectivo topológico de $\Lambda$-álgebras cociente que son noetherianas y
lisas sobre $\Lambda$ en el sentido anterior. Se generaliza inmediatamente al caso
en que $\Lambda$ y $O$ ya no se suponen locales. Dicho esto, se tiene la
proposición siguiente:

\result{PROPOSICIÓN}{5.2} Para que $O$ sea liso sobre $\Lambda$, es necesario
y suficiente que el funtor asociado $F$ transforme epimorfismos en
epimorfismos.

Esto significa que, para todo homomorfismo sobreyectivo
$A\longrightarrow A'$ en $\mathcal C$, $F(A)\longrightarrow F(A')$ es
sobreyectivo. Naturalmente, basta comprobar esta condición cuando $A$ es local
y, procediendo paso a paso, cuando el ideal de $A$ que es núcleo de
$A\longrightarrow A'$ está anulado por el ideal maximal de $A$. En la práctica,
esto reduce la cuestión a comprobar que cierta obstrucción, ligada a invariantes
infinitesimales

\src{377}
de la situación que da origen al funtor $F$, es nula, problema de naturaleza
cohomológica.

Para terminar, digamos unas palabras, en el contexto anterior, sobre los
\emph{anillos de definición}. Sea todavía $F$ un funtor de $\mathcal C$ en
$(\operatorname{Ens})$, pro-representable mediante una $\Lambda$-álgebra
topológica $O$. Entonces, para todo $A\in\mathcal C$ y todo $\xi\in F(A)$,
existe un menor subanillo $A'$ de $A$ tal que $\xi$ sea imagen de un elemento
$\xi'$ de $F(A')$ (que queda entonces determinado de manera única): en efecto,
basta interpretar $\xi$ mediante un homomorfismo de $O$ en $A$ y tomar como
$A'$ la imagen de $O$ por este último. Se dice entonces que $A'$ es el anillo
de definición del objeto $\xi\in F(A)$. Si $u:A\longrightarrow B$ es un
homomorfismo de álgebras y $\eta=F(u)(\xi)$, el anillo de definición de $\eta$
es la imagen por $u$ del anillo de definición de $\xi$. Cuando se parte de un
funtor $F$ de $\mathcal C$ en $(\operatorname{Ens})$, la existencia de los
anillos de definición y las propiedades que acabamos de dar son, con pequeñas
salvedades, equivalentes a que $F$ sea pro-representable; es decir, por lo
general distan mucho de ser triviales.

\fsection{B}{Los dos teoremas de existencia}

Conservemos las notaciones de A, párrafo 5, y partamos de un funtor covariante
\[
  F:\mathcal C\longrightarrow(\operatorname{Ens});
\]
se buscan criterios manejables para que $F$ sea pro-representable, es decir,
expresable mediante una $\Lambda$-álgebra $O$ como arriba. En virtud del
corolario de A, proposición 3.1, para ello es necesario y suficiente que $F$ sea
exacto por la izquierda. En el estado actual de la técnica de descenso (cf. las
cuestiones planteadas en [3], p. 9), este criterio no puede comprobarse
directamente bajo esta forma en los casos más importantes, y se necesitan
criterios aparentemente menos exigentes.

\result{TEOREMA}{1} Para que el funtor $F$ sea pro-representable, es necesario
y suficiente que satisfaga las dos condiciones siguientes:
\begin{enumerate}
\item[(i)] $F$ conmuta con los productos finitos;
\item[(ii)] para toda álgebra $A\in\mathcal C$ y todo homomorfismo
$A\longrightarrow A'$ en $\mathcal C$ tal que el diagrama
\[
  A\longrightarrow A'\rightrightarrows A'\otimes_A A'
\]
sea exacto ([3], A, definición 1.2), el diagrama transformado
\[
  F(A)\longrightarrow F(A')\rightrightarrows F(A'\otimes_A A')
\]

\src{378}
es exacto.
\end{enumerate}
Además, basta comprobar (ii) cuando $A$ es local y se está, además, en uno de
los dos casos siguientes:
\begin{enumerate}
\item[(a)] $A'$ es un módulo libre sobre $A$;
\item[(b)] el módulo cociente $A'/A$ es un $A$-módulo de longitud 1.
\end{enumerate}

La demostración de este teorema es bastante delicada y no puede esbozarse aquí.
Limitémonos a señalar que descansa esencialmente en un estudio de las relaciones
de equivalencia (en el sentido de las categorías) en el espectro de un álgebra
artiniana (estudio que plantea todavía varios problemas cuya solución parece
indispensable para el desarrollo ulterior de la teoría).

En las aplicaciones, la comprobación de la condición (i) es siempre trivial.
La de (ii) se descompone en dos: el caso en que $A'$ es un $A$-módulo libre
corresponde a la técnica de descenso por morfismos planos ([3], teoremas 1, 2 y
3) y no ofrece dificultad. Para tratar el caso (b), se utilizará el resultado
siguiente:

\result{TEOREMA}{2} Sean $A$ un anillo local artiniano de ideal maximal
$\mathfrak m$ y $A'$ una $A$-álgebra que contiene a $A$, tal que
$\mathfrak mA'\subset A$ y que
\[
  A\longrightarrow A'\rightrightarrows A'\otimes_A A'
\]
sea exacto (lo cual ocurre, en particular, si $A'/A$ es un $A$-módulo de
longitud 1). Sea $\mathcal F$ la categoría fibrada ([3], A, definición 1.1) de
los haces cuasicoherentes y planos sobre preesquemas variables. Entonces el
morfismo $\operatorname{Spec}(A')\longrightarrow\operatorname{Spec}(A)$ es un
morfismo de $\mathcal F$-descenso estricto ([3], A, definición 1.7).

En otros términos, el dato de un $A$-módulo plano $M$ es completamente
equivalente al dato de un $A'$-módulo plano $M'$, provisto de un
$A'\otimes_AA'$-isomorfismo de $M'\otimes_AA'$ sobre $A'\otimes_AM'$ que
satisface la condición de transitividad usual para un dato de descenso (loco
citato).

El teorema 2 se demuestra probando primero que
\[
  H^i(A'/A,G_a)=0\qquad\text{para }i\geq 1
\]
([3], A, 4(e)), pues la hipótesis $\mathfrak mA'\subset A$ permite reducirse
fácilmente al caso en que $A$ es un cuerpo $k$ (a saber, $A/\mathfrak m$). Se
aplican entonces las equivalencias señaladas en [3], página 16.

\fsection{C}{Aplicaciones a algunos casos particulares}

\subsection*{1. Observaciones generales sobre los funtores representados por preesquemas}
\addcontentsline{toc}{subsection}{1. Observaciones generales sobre los funtores representados por preesquemas}

Sea $S$ un preesquema localmente noetheriano. Se dice que un preesquema $X$
sobre $S$ es \src{379}localmente de tipo finito sobre $S$ si, para todo
$x\in X$ que se proyecta en $y\in S$,\footnote{El texto impreso dice aquí
« $y\in Y$ » y vuelve a emplear $Y$ al final del párrafo. Como el morfismo
considerado es $X\longrightarrow S$, se lee $S$ en ambas apariciones.} existen
un entorno afín de $y$, de anillo $A$, y un entorno afín de $x$ situado sobre
el anterior, de anillo $B$, tales que $B$ sea un álgebra de tipo finito sobre
$A$. Se encuentran muchos ejemplos importantes de preesquemas localmente de
tipo finito sobre $S$ que no son de tipo finito sobre $S$, como soluciones de
problemas universales clásicos; así, importa poder considerar el esquema de
Picard de una curva como la reunión de una infinidad de componentes conexas
(que no deben confundirse con la componente conexa del elemento neutro, es
decir, la «variedad de Picard»). En ocasiones resulta entonces conveniente
trabajar en la categoría $\mathcal C$ de los preesquemas localmente de tipo
finito sobre $S$, para examinar en ella la cuestión de la representabilidad de
un funtor contravariante $F$. El objetivo principal de estas exposiciones es
desarrollar una técnica general que permita reconocer si tal funtor $F$ es
representable y estudiar las propiedades del $S$-preesquema $X$ correspondiente
a partir de las de $F$. Señalemos de paso que, en este estudio, aparecen
ejemplos no patológicos de preesquemas sobre $S$ no separados sobre $S$, en
particular como «preesquemas de Picard» de excelentes $S$-esquemas; no deben,
por tanto, desterrarse de la geometría algebraica los preesquemas que no son
esquemas.

Sea $X$ un preesquema localmente de tipo finito sobre $S$, y sea
\[
  F(Y)=\operatorname{Hom}_S(Y,X)
\]
el funtor contravariante asociado. Se puede considerar la restricción $F_0$ de
$F$ a la subcategoría $\mathcal C_0$ de $\mathcal C$ formada por los
preesquemas $Y$ sobre $S$ que son artinianos y finitos sobre $S$: cuando
$S=\operatorname{Spec}(\Lambda)$, $\mathcal C_0$ es, pues, la categoría dual
de la categoría de las $\Lambda$-álgebras artinianas considerada en B. Si
$Y=\operatorname{Spec}(A)$, donde $A$ es un anillo local artiniano, $Y$ se
reduce a un solo punto $y$ situado sobre un punto cerrado $s$ de $S$, y un
$S$-homomorfismo de $Y$ en $X$ (es decir, un elemento de $F(Y)$) queda definido
por el dato de un punto $x\in X$ sobre $s$ y de un
$\mathcal O_{s,S}$-homomorfismo de $\mathcal O_{x,X}$ en $A$. Si existe tal
homomorfismo, entonces $x$ es necesariamente un punto cerrado de $X$ (pues su
cuerpo residual es algebraico sobre el de $s$). Esto muestra que la restricción
$F_0$ de $F$ a las «$\Lambda$-álgebras artinianas» es pro-representable, y está
representada por el álgebra topológica sobre $\Lambda$ cuyas componentes
locales son los completados $\widehat{\mathcal O}_{x,X}$ de los anillos locales
de $X$ en los puntos $x$ de $X$ que son cerrados y están situados sobre puntos
cerrados de $S$. Esto muestra que el solo conocimiento de $F_0$ proporciona
información precisa sobre la estructura de $X$ (a saber, la estructura de los
completados de sus anillos locales en los puntos indicados). Obsérvese que,
incluso cuando $S$

\src{380}
es el espectro de un cuerpo algebraicamente cerrado, solo mediante la
consideración sistemática de «variedades» $Y$ tales que $\mathcal O_Y$ pueda
tener elementos nilpotentes, en particular trabajando con los espectros de
anillos artinianos locales, se puede llegar a la «formulación correcta» de los
problemas universales clásicos y captar su aspecto «infinitesimal».

Cuando se parte de un funtor $F$ dado de antemano y se quiere decidir si es
representable, el estudio del funtor $F_0$ (con ayuda de los teoremas 1 y 2)
proporcionará indicaciones casi completas: o bien, como ocurre con frecuencia
(comprobando simplemente, por ejemplo, la naturaleza de los conjuntos
$F(I_K,\xi)$ y su comportamiento funtorial, cf. A), se constata que $F_0$ ya no
es pro-representable (lo que explica el fracaso de los intentos realizados
hasta ahora para definir de manera más o menos natural variedades de módulos
para la clasificación de los fibrados vectoriales de rango $>1$); o bien se
logra comprobar que $F_0$ es representable, pero los espacios vectoriales
$F(I_K,\xi)$ no son de dimensión finita, en cuyo caso hay que conformarse con
la solución «formal»; o bien, por último, $F_0$ está representado por un
producto de anillos locales completos noetherianos, lo que aporta indicios muy
sólidos de que el propio $F$ es representable y, unido a propiedades análogas
pero de naturaleza más global que pensamos desarrollar más adelante, bastará
sin duda para implicar que así sea efectivamente. Finalmente, se encuentran
problemas geométricos interesantes (véanse los párrafos 4 y 5 más abajo) en los
que solo se dispone del funtor $F_0$ (que no procede de ningún funtor «global»
$F$), y en los que habrá motivos para darse por satisfecho si se logra asociarle
un «esquema formal de módulos».

Para terminar estas generalidades, indiquemos cómo la teoría de esquemas
explica anomalías aparentes, como la superficie de Igusa $V$, cuya «variedad
de Picard» $P$ se reduce a un punto y para la cual, sin embargo, se tiene
\[
  H^1(V,\mathcal O_V)\ne 0 ;
\]
en este caso, $P$ es un grupo «puramente infinitesimal» que no se reduce al
grupo unidad, es decir, está definido por un álgebra local $O$ de rango finito
sobre el cuerpo base $k$ y dotado de una aplicación diagonal correspondiente a
la estructura multiplicativa de $P$; si $\mathfrak m$ es el ideal maximal de
$O$, el dual de $\mathfrak m/\mathfrak m^2$ es canónicamente isomorfo a
$H^1(V,\mathcal O_V)$ (cf. el párrafo 3 más abajo). Solo cuando el grupo de
Picard es un grupo algebraico en el sentido clásico, es decir, liso sobre el
cuerpo base $k$, la dimensión de $H^1(V,\mathcal O_V)$ (siempre igual a la de
$\mathfrak m/\mathfrak m^2$) es igual a la del grupo de Picard.

\subsection*{2. Esquemas $\underline{\operatorname{Hom}}_S(X,Y)$,
$\Pi_{X/S}Z$, $\underline{\operatorname{Aut}}_S(X)$, etc.}
\addcontentsline{toc}{subsection}{2. Esquemas Hom, Pi y Aut}

Sean $X,Y$ dos preesquemas sobre $S$; para todo preesquema $T$ sobre $S$,

\src{381}
sean $X_T=X\times_ST$, $Y_T=Y\times_ST$, y consideremos el conjunto
\[
  F(T)=\operatorname{Hom}_T(X_T,Y_T)
      =\operatorname{Hom}_S(X_T,Y)
      =\operatorname{Hom}_S(X\times_ST,Y)
\]
como funtor contravariante en $T$. Si es representable, se denota por
$\underline{\operatorname{Hom}}_S(X,Y)$ el preesquema sobre $S$ que lo
representa; se tiene, por tanto, un isomorfismo funtorial
\[
  \operatorname{Hom}_S
  \bigl(T,\underline{\operatorname{Hom}}_S(X,Y)\bigr)
  \xrightarrow{\sim}
  \operatorname{Hom}_S(T\times_SX,Y).
\]
Hay muchas variantes de este problema universal cuya solución se reduce a él:
el preesquema de los $S$-automorfismos de un $S$-preesquema $X$ (que será un
preesquema en grupos), el preesquema de los $S$-homomorfismos de un
$S$-preesquema en grupos en otro (que será un preesquema en grupos conmutativos
si el segundo esquema en grupos es conmutativo), etc. También se puede
generalizar la definición de $\underline{\operatorname{Hom}}_S(X,Y)$
considerando un preesquema $Z$ sobre el preesquema $X$ sobre $S$, y el funtor
\[
  F(T)=\operatorname{Hom}_{X_T}(X_T,Z_T)
\]
(conjunto de las «secciones» de $Z_T$ sobre $X_T$); si este funtor es
representable, el $S$-preesquema que lo representa se denotará por
$\Pi_{X/S}Z$; se tendrá, pues, por definición, un isomorfismo funtorial
\[
  \operatorname{Hom}_S(T,\Pi_{X/S}Z)
  =\operatorname{Hom}_{X_T}(X_T,Z_T).
\]
Tomando $Z=Y\times_SX$, se recupera
$\underline{\operatorname{Hom}}_S(X,Y)$. De estas definiciones resulta para los
nuevos preesquemas así introducidos un formulario tan trivial como útil, que no
daremos aquí (puesto que es válido en toda categoría en la que existen
productos y productos fibrados). Más seria es la cuestión de la existencia de
esquemas del tipo $\underline{\operatorname{Hom}}_S(X,Y)$. Se constata primero
que, para $X$ fijo, $\underline{\operatorname{Hom}}_S(X,Y)$ solo puede existir
para todo $Y$ sobre $S$ (o $\Pi_{X/S}Z$ para todo $Z$ sobre $X$) si $X$ es plano
sobre $S$. Además, se advierte que solo es razonable esperar la existencia de
una solución, para $Y$ suficientemente general, si $X$ es además propio sobre
$S$. Por otra parte, parece que estas condiciones son suficientes para la
existencia de $\underline{\operatorname{Hom}}_S(X,Y)$ y $\Pi_{X/S}Z$, a
condición de imponer, cuando proceda, una hipótesis del tipo «cuasiproyectivo»
sobre $Y/S$, respectivamente sobre $Z/X$; esto puede comprobarse, en cualquier
caso, en bastantes situaciones (por ejemplo, cuando $Y$ es afín sobre $S$; o,
mediante construcciones directas elementales, cuando $X$ es finito sobre $S$).
He aquí lo que dan los teoremas 1 y 2:

\src{382}
\result{PROPOSICIÓN}{2.1} Sean $\Lambda$ un anillo noetheriano, $X$ e $Y$ dos
preesquemas cualesquiera sobre $\Lambda$, y consideremos el funtor
\[
  F(A)=\operatorname{Hom}_A(X_A,Y_A)
\]
en la categoría $\mathcal C_0$ de las $\Lambda$-álgebras artinianas. Si $X$ es
plano sobre $\Lambda$, este funtor es pro-representable.

Además, se comprueba que para todo $A\in\mathcal C_0$ y $\xi\in F(A)$ se tiene
un isomorfismo canónico
\[
  F(I_A,\xi)=H^0\!\left(
    X_A,
    \underline{\operatorname{Hom}}_{\mathcal O_{X_A}}
    \bigl(\xi^*(\Omega^1_{Y_A/A}),\mathcal O_{X_A}\bigr)
  \right),
\]
donde $\Omega^1_{Y_A/A}$ es el haz de $1$-diferenciales de Kähler de $Y_A$
respecto de $A$. Tomando como $A$ un cuerpo, se obtiene, utilizando A.5.1 y el
teorema de finitud de [2], el corolario siguiente:

\result{COROLARIO}{} Supongamos que $X$ es plano y propio sobre $S$, y que $Y$
es localmente de tipo finito sobre $S$. Entonces $F$ es pro-representable y las
componentes locales del álgebra topológica correspondiente sobre $\Lambda$ son
anillos noetherianos.

\result{OBSERVACIONES}{} Los problemas considerados en este apartado, y muchos
otros, se estudiaban habitualmente, en el marco de la geometría algebraica
clásica, mediante las «coordenadas de Chow» de los ciclos en el espacio
proyectivo, que permitían considerar dichos ciclos como puntos de variedades
proyectivas adecuadas. Este procedimiento, y en general el uso de las
coordenadas de Chow, parece irremediablemente insuficiente desde el punto de
vista de los esquemas, pues destruye los elementos nilpotentes en las
variedades de parámetros y, en particular, no se presta a un estudio
satisfactorio de las variaciones infinitesimales de ciclos (por no hablar de su
naturaleza no intrínseca, ligada al espacio proyectivo). El lenguaje de las
coordenadas de Chow era, por desgracia, el único utilizado por muchos
geómetras algebraicos para estudiar familias de variedades o familias de
ciclos, lo que parece haber constituido un obstáculo serio para aclarar estas
nociones, pese a su indudable interés técnico (probablemente provisional). Si
se quiere obtener el análogo de las variedades de Chow en la teoría de
esquemas, se llega al problema universal siguiente: sea $X$ un preesquema sobre
$S$; para todo preesquema $T$ sobre $S$, se considera el conjunto $F(T)$ de los
subpreesquemas cerrados de $X_T=X\times_ST$ que son planos sobre $T$. Se intenta
representar este funtor en $T$ mediante un preesquema sobre $S$. Más
generalmente, se puede dar un haz cuasicoherente $G$ sobre $X$ y tomar como
$F(T)$ el conjunto de los haces

\src{383}
cuociente de $G_T$ que son planos sobre $T$. Parece que el problema tiene una
solución, mediante un esquema $C$ localmente de tipo finito sobre $S$, cuando
$X$ es propio sobre el preesquema localmente noetheriano $S$ y $G$ es además
coherente. En cualquier caso, suponiendo únicamente que $S$ es localmente
noetheriano, la restricción de $F$ a las «$S$-álgebras artinianas» es
pro-representable; y si, además, $X$ es propio sobre $S$ y $G$ es coherente, las
componentes locales del anillo topológico correspondiente $O$ son
noetherianas. Por supuesto, aun después de demostrar la existencia del
«esquema de Chow» de $X$ sobre $S$, quedará por
encontrar una descomposición en abiertos disjuntos $C_i$ (correspondientes a
la fijación de invariantes continuos, como el grado y la dimensión de los
ciclos que se hacen variar) que sean de tipo finito (y no solo localmente de
tipo finito) sobre $S$, precisar las relaciones de este esquema con las
variedades de Chow clásicas y precisar cuándo un $C_i$ es incluso proyectivo,
o al menos cuasiproyectivo, sobre $S$.\footnote{La fe de erratas de 1962
precisa que el caso proyectivo queda completamente resuelto por los esquemas de
Hilbert (TDTE IV), mientras que ejemplos de Nagata y Hironaka muestran que, sin
una hipótesis de proyectividad, estos funtores pueden no ser representables,
incluso para las subvariedades de dimensión cero de una variedad completa no
singular de dimensión tres; este fenómeno está ligado a la posible
inexistencia del cuadrado simétrico.}

\subsection*{3. Esquemas de Picard}
\addcontentsline{toc}{subsection}{3. Esquemas de Picard}

Sea $f:X\longrightarrow S$ un $S$-preesquema,\footnote{Para un estudio más
completo, la fe de erratas de 1962 remite a TDTE V.} consideremos el haz
multiplicativo $\mathcal O_X^*$ de las unidades del haz estructural de $X$, y
el grupo
\[
  P(X/S)=H^0\bigl(S,R^1f_*(\mathcal O_X^*)\bigr),
\]
llamado grupo de Picard relativo de $X/S$. Un elemento de este grupo queda,
pues, definido dando una cubierta abierta $(U_i)$ de $S$ y un haz invertible
$L_i$ sobre cada $f^{-1}(U_i)$, de modo que, para todo par $(i,j)$,
$L_i|f^{-1}(U_i\cap U_j)$ sea isomorfo a $L_j|f^{-1}(U_i\cap U_j)$, al menos
localmente sobre $U_i\cap U_j$ (es decir, estos dos haces sean «equivalentes»
en el sentido de [3], B, 4). Si $X/S$ admite una sección, entonces $P(X/S)$ no
es sino el conjunto de las clases de haces invertibles sobre $X/S$ módulo
«equivalencia» (loco citato). Pongamos ahora, para todo $T$ sobre $S$,
\[
  F(T)=P(X_T/T) ;
\]
se obtiene un funtor contravariante en $T$, que podrá llamarse funtor de Picard de
$X/S$; si este funtor es representable, el preesquema sobre $S$ que lo
representa se llamará preesquema de Picard de $X/S$ y se denotará por
$\underline P(X/S)$. Se tendrá, pues, un isomorfismo de funtores
\[
  \operatorname{Hom}_S\bigl(T,\underline P(X/S)\bigr)
  \xrightarrow{\sim}P(X_T/T).
\]
La formación de preesquemas de Picard es compatible con la extensión de la
base; en particular, los preesquemas de Picard de las fibras de $X$ sobre $S$
(que son preesquemas sobre los cuerpos residuales $K(s)$ de los $s\in S$) son
las fibras de $\underline P(X/S)$. Por supuesto, como $P(X_T/T)=F(T)$ es un
grupo conmutativo, los preesquemas de Picard son preesquemas en grupos.
Obsérvese también que las jacobianas generalizadas de Rosenlicht no son sino
las componentes conexas de la unidad en los esquemas de Picard de curvas
completas que pueden tener singularidades, lo que debería hacer evidentes la
mayoría de sus propiedades (una vez demostrada la existencia).

\result{OBSERVACIÓN}{} La definición adoptada aquí solo es razonable cuando
todo punto de $S$\footnote{El texto impreso dice «$Y$»; aquí no se ha definido
ninguna base $Y$, y la condición se refiere a la base $S$.} admite un entorno
abierto $U$ sobre el cual $X$ admite una sección. En el caso general, hay que
modificar ligeramente la definición del funtor de Picard para poder obtener
también un teorema de existencia.

Aquí, las condiciones plausibles para la existencia de un preesquema de Picard
son las siguientes: $X$ es propio y plano sobre $S$,
$f_*(\mathcal O_X)=\mathcal O_S$, y $X$ admite localmente una sección sobre
$S$. Esta condición aparece de manera natural al aplicar la técnica de
descenso, para eliminar los automorfismos de un haz invertible $L$ sobre $X$
dotándolo de una sección marcada sobre la sección $s$ ([3], B, 4). Se obtiene,
en particular:

\src{384}
\result{PROPOSICIÓN}{3.1} Supongamos que $X$ es plano sobre
$S=\operatorname{Spec}(\Lambda)$, con $\Lambda$ noetheriano, y que, para todo
$T$ de tipo finito sobre $S$, se tiene
\[
  f_{T*}(\mathcal O_{X_T})=\mathcal O_T
\]
(si $f$ es propio y separable,\footnote{La fe de erratas de 1962 prescribe
explícitamente la hipótesis «$f$ propio y separable».} o si $S$ es el espectro
de un cuerpo, resulta de Künneth que la última condición equivale a
$f_*(\mathcal O_X)=\mathcal O_S$). Entonces el funtor de Picard de $X/S$ en la
categoría de las $\Lambda$-álgebras artinianas es pro-representable.

Además, aquí se tendrá
\[
  F(I_A,\xi)=H^1(X_A,\mathcal O_{X_A}),
\]
en particular:

\result{COROLARIO}{} Si $X$ es propio sobre $S$, entonces las componentes
locales de la $\Lambda$-álgebra topológica $O$ correspondiente al funtor de
Picard son noetherianas.

\result{OBSERVACIONES}{} Se pueden generalizar las definiciones y los resultados
de este apartado a la clasificación de los fibrados principales sobre $X$, con
grupo estructural un esquema en grupos $G$ sobre $S$ que sea afín, plano sobre
$S$ y conmutativo. Cuando $G$ no es conmutativo, y por tanto el fibrado en
grupos adjunto de un fibrado principal (cuyas secciones son los automorfismos
del fibrado principal) deja de ser trivial, la proposición 3.1 ya no es válida
tal cual. No obstante, se puede modificar el problema universal de manera que
siga teniendo solución (al menos, por ahora, en geometría formal). La regla de
oro que debe recordarse, en el contexto de este \src{385}apartado y de los
siguientes, y siempre que se busquen «esquemas de módulos» para clases de
objetos definidos solo salvo isomorfismo, es siempre la siguiente: eliminar
los posibles automorfismos de los objetos que se quieren clasificar mediante
la introducción, si es necesario, de estructuras auxiliares (puntos o
elementos de secciones marcadas, fijación de formas diferenciales, etc.), que
se elegirán lo bastante inocuas como para no modificar sustancialmente el
problema inicial.

\subsection*{4. Módulos formales de una variedad}
\addcontentsline{toc}{subsection}{4. Módulos formales de una variedad}

Sea $\Lambda$ un anillo local noetheriano de cuerpo residual $k$ (con mayor
frecuencia, $\Lambda$ será igual a $k$, o a un $p$-anillo de Cohen), y sea
$X_0$ un preesquema sobre $k$. Para toda $\Lambda$-álgebra artiniana local $A$,
se considera el conjunto $F(A)$ de las clases, salvo isomorfismo, de
$A$-preesquemas $X$ planos sobre $A$, dotados de un isomorfismo
\[
  (*)\qquad
  X\otimes_A\kappa(A)\xleftarrow{\sim}
  X_0\otimes_k\kappa(A),
\]
donde $\kappa(A)$ es el cuerpo residual de $A$; por supuesto, los isomorfismos
entre tales $A$-preesquemas planos deben respetar el isomorfismo precedente
incluido en la estructura. Si $A$ es una $\Lambda$-álgebra artiniana no
necesariamente local, de componentes locales $A_i$, se toma como $F(A)$ el
producto de los $F(A_i)$. De este modo, $F$ se convierte en un funtor
multiplicativo en $A$, que podrá llamarse funtor de módulos para $X_0$ (y
$\Lambda$). Si este funtor es representable, le corresponde una
$\Lambda$-álgebra topológica $O$ local, de cuerpo residual $k$, y el espectro
formal de $O$ se llamará esquema formal de módulos para $X_0$ (y $\Lambda$);
cf. [2] para algunos detalles sobre esta noción.

Aquí, para aplicar la técnica de descenso, los automorfismos «finitos» de $X_0$
son inofensivos, pues no influyen en la existencia de automorfismos (en el
sentido precisado más arriba) de los $A$-preesquemas $X$; cuando $A$ no se
reduce a un cuerpo, la condición necesaria y suficiente para que $X$ no tenga
ningún $A$-automorfismo no trivial es que se tenga
\[
  H^0(X_0,\Theta_{X_0/k})=0,
\]
donde $\Theta_{X_0/k}$ denota el haz de las $k$-derivaciones (= haz tangente)
de $X_0$. Por otra parte, se obtiene fácilmente (al menos si $X_0$ es liso
sobre $k$)
\[
  F(I_A,\xi)=H^1(X_A,\Theta_{X_A/A}).
\]
Se concluye entonces, como de costumbre:

\src{386}
\result{PROPOSICIÓN}{4.1} Supongamos que
$H^0(X_0,\Theta_{X_0/k})=0$. Entonces existe el esquema formal de módulos para
$X_0$. Si, además, $X_0$ es propio sobre $k$, el esquema formal de módulos es
noetheriano.

\result{OBSERVACIONES}{} 1\textsuperscript{o} Cuando no se supone que $X_0$ sea
liso sobre $k$, $F(I_A,\xi)$ se identifica con un sub-$A$-módulo de
\[
  \operatorname{Ext}^1_{\mathcal O_{P_A}}
  \bigl(P_A;\mathcal I_{X_A},\mathcal O_{X_A}\bigr),
\]
donde se pone $P_A=X_A\times_\Lambda X_A$, donde $\mathcal O_{X_A}$ se
considera como un haz coherente sobre $P_A$ mediante el morfismo diagonal
$X_A\longrightarrow P_A$, y donde $\mathcal I_{X_A}$ denota el haz coherente
de ideales sobre $P_A$ definido por el morfismo diagonal. Más precisamente,
una globalización sencilla de la teoría de Hochschild muestra que el
$\operatorname{Ext}^1$ anterior se identifica con el conjunto de las clases,
salvo isomorfismo, de haces de $I_A$-álgebras planas $\mathcal O$ sobre $X_A$,
dotados de un isomorfismo de aumentación
\[
  \mathcal O\otimes_{I_A}A\longrightarrow\mathcal O_{X_A}
\]
(se recuerda que se pone $I_A=A[t]/(t^2)$). El submódulo $F(I_A,\xi)$ es el
que corresponde a los haces de álgebras conmutativas. Las hipótesis de lisura
no son, pues, esenciales en la teoría de módulos, contrariamente a lo que [2]
daba a entender.

2\textsuperscript{o} Recordemos (loco citato) que, en particular, toda curva
algebraica $X_0$ lisa y propia sobre $k$ admite un esquema formal de módulos,
liso sobre $\Lambda$, de dimensión relativa $3g-3$ si el género $g$ es
$\geq2$, y $g$ si $g=0$ o $1$. Estos dos últimos casos ya no entran directamente
en la proposición 4.1. Sin embargo, en el caso de las curvas elípticas
($g=1$), podrían reducirse a ella mediante las consideraciones que siguen.

Por supuesto, se puede variar ad libitum la proposición 4.1 considerando
sistemas de esquemas sobre $k$ dotados de diversas estructuras. Supongamos,
por ejemplo, que $X_0$ sea un esquema abeliano sobre $k$, con punto origen
marcado (es decir, que $X_0$ se considere como esquema en grupos sobre $k$), y
sea $F(A)$ el conjunto de las clases, salvo isomorfismo, de esquemas abelianos
sobre $A$ (es decir, de esquemas en grupos propios y lisos sobre $A$) dotados
de un isomorfismo de esquemas abelianos $(*)$. Se comprueba que la imposición
de una estructura multiplicativa (e incluso de una mera «sección unidad»)
elimina los automorfismos infinitesimales y que, por consiguiente, existe un
esquema formal de módulos, correspondiente a un anillo local completo
noetheriano $O$. Por otra parte, puede demostrarse que, si $X$ es un esquema
propio y liso con fibras «absolutamente conexas» sobre un preesquema localmente
noetheriano $S$, toda estructura multiplicativa sobre $X$ que admita una
sección unidad es necesariamente asociativa y conmutativa (siempre que lo sea,
al menos, sobre una fibra y que $S$ sea conexo), y además queda determinada de
manera única por el conocimiento

\src{387}
de la sección unidad. Además, supongamos que $S$ es el espectro de un anillo
local artiniano $A$ de cuerpo residual $k$, que $X$ es propio sobre $A$ y está
dotado de una sección $s$, y que, por último, $X\otimes_Ak$ está dotado de una
estructura de esquema abeliano sobre $k$ cuyo elemento neutro es el punto de
$X\otimes_Ak$ correspondiente a $s$; un cálculo sencillo de obstrucciones,
unido a un razonamiento debido a Serre, permite demostrar que existe
efectivamente sobre $X$ una estructura multiplicativa que admite la sección
$s$ como sección unidad. (A partir de ahí, utilizando el «teorema de
existencia» de [2] para pasar al caso en que $A$ es noetheriano local completo,
y luego la técnica de descenso de [3] para el caso general, se puede demostrar
el enunciado análogo para todo $S$ localmente noetheriano conexo.) Esto prueba
que el funtor $F(A)$ considerado aquí es isomorfo al funtor análogo definido al
principio de este apartado prescindiendo de la estructura multiplicativa sobre
$X_0$. Se sigue, en particular, que, si $\mathfrak m$ es el ideal maximal de
$O$, $\mathfrak m/\mathfrak m^2$ es canónicamente isomorfo al dual de
$H^1(X_0,\Theta_{X_0/k})$ y tiene, por tanto, dimensión $n^2$, donde
$n=\dim X_0$. Sería muy interesante determinar si $O$ es liso sobre $\Lambda$,
es decir, isomorfo a un álgebra de series formales en $n^2$ indeterminadas
sobre $\Lambda$. El párrafo 1, proposición 5.2, permite dar una formulación
equivalente de este problema como problema de existencia de esquemas abelianos
que se reduzcan según un esquema abeliano dado. En cualquier caso, por métodos
trascendentes se ve que la respuesta es afirmativa si $k$ es de característica
cero. En característica $p\ne0$, basta evidentemente limitarse al caso en que
$\Lambda$ es el anillo de vectores de Witt construido sobre un cuerpo
algebraicamente cerrado $k$. Tal vez haya llegado el momento de que el «funtor
de Greenberg» demuestre su valía\ldots\footnote{La fe de erratas de 1962
anuncia que el esquema formal de módulos de una variedad abeliana sobre un
cuerpo es liso sobre el anillo de Witt: todo esquema abeliano sobre un cociente
artiniano local se levanta al anillo del cual es cociente. La demostración
utiliza la varianza de la clase de obstrucción al levantamiento introducida en
la exposición 182, p. 12. También recuerda las construcciones de Mumford para
las curvas de género $g$ y los esquemas abelianos polarizados (cf. SHMT).}

\subsection*{5. Extensión de los revestimientos}
\addcontentsline{toc}{subsection}{5. Extensión de los revestimientos}

Sean $\mathfrak X$ un preesquema formal noetheriano [2] y $U$ una parte
abierta de $\mathfrak X$, definida localmente por la «no anulación» de una
sección de $\mathcal O_{\mathfrak X}$ que no sea divisor de cero no solo sobre
$\mathfrak X$, sino sobre todo $X_n$,\footnote{Precisión impuesta por la fe de
erratas de 1962.} y que sea, por tanto, lo bastante grande para que toda sección
de $\mathcal O_{\mathfrak X}$ sobre un abierto $V$ que sea nula sobre $U\cap V$
sea nula. Sea $\mathcal J$ un «ideal de definición» para $\mathfrak X$, y sea
\[
  X_n=(\mathfrak X,\mathcal O_{\mathfrak X}/\mathcal J^{n+1}),
\]
que es, por tanto, un preesquema noetheriano ordinario. Si entonces
$\mathfrak X',\mathfrak X''$ son dos revestimientos planos de $\mathfrak X$ (es
decir, dos preesquemas sobre $\mathfrak X$ definidos por haces de álgebras que
son coherentes y localmente libres como haces de módulos), no ramificados sobre
$U$, la aplicación evidente
\[
  \operatorname{Hom}_{\mathfrak X}(\mathfrak X',\mathfrak X'')
  \longrightarrow
  \operatorname{Hom}_{X_0}(X'_0,X''_0)
\]
es inyectiva; en particular, un automorfismo de $\mathfrak X'$ que induce la
identidad sobre $X'_0$ es la identidad. Esto permite aplicar a la situación la
técnica de descenso. Partamos, en particular, de un revestimiento plano $X'_0$
de $X_0$, no ramificado

\src{388}
sobre $U_0$; sea $G(\mathfrak X)$ el conjunto de las clases, salvo isomorfismo
(que induzca la identidad sobre $X'_0$), de revestimientos planos $\mathfrak X'$
de $\mathfrak X$ que inducen $X'_0$ sobre $X_0$ (y son, por tanto,
necesariamente no ramificados sobre $U$). Se define de igual modo $G(V)$ para
toda parte abierta $V$ de $\mathfrak X$, y más generalmente $G(\mathfrak Y)$
para todo preesquema formal $\mathfrak Y$ sobre $\mathfrak X$. Dicho esto, los
resultados de [2] y [3] implican en primer lugar los resultados siguientes:

\noindent a. Cuando $V$ es un abierto variable de $\mathfrak X$, los $G(V)$
forman un haz sobre $\mathfrak X$, sea $G_{\mathfrak X}=G$. La restricción de
este haz a $U$ es el haz constante cuyas fibras se reducen a un elemento.

Más generalmente, la determinación de las fibras de $G_{\mathfrak X}$ es una
cuestión de anillos locales completos; más precisamente:

\noindent b. Para todo $x\in\mathfrak X$, se tiene
\[
  G_x=G\bigl(\operatorname{Spec}(\mathcal O_{\mathfrak X,x})\bigr)
  \subset
  G\bigl(\operatorname{Spec}(\widehat{\mathcal O}_{\mathfrak X,x})\bigr),
\]
(= clases salvo isomorfismo de álgebras finitas y libres $B$ sobre
$\widehat{\mathcal O}_{\mathfrak X,x}$ dotadas de un isomorfismo de
$B\otimes_{\widehat{\mathcal O}_{\mathfrak X,x}}\mathcal O_{X_0,x}$ con
$\widehat{\mathcal O}'_{0x}$, donde $\mathcal O'_0$ es el haz de álgebras sobre
$X_0$ que define $X'_0$).\footnote{La fórmula impresa identifica directamente
$G_x$ con el término completado y omite el sombrero sobre
$\mathcal O'_{0x}$. La inclusión y el completado son los prescritos por la fe
de erratas de 1962.}

\noindent c. Se tiene un isomorfismo canónico
\[
  G_{\mathfrak X}=\varprojlim_nG_{X_n} ;
\]
en otros términos, para todo abierto $V$ de $\mathfrak X$, se tiene
\[
  G(V)=\varprojlim_nG(V_n).
\]

\noindent d. Supongamos que $\mathfrak X$ se obtiene de un esquema ordinario
propio $X$ sobre un anillo local noetheriano completo $\Lambda$ que posee un
ideal de definición $\mathfrak m$, tomando el completado $\mathcal J$-ádico de
$\mathcal O_X$, donde $\mathcal J=\mathfrak m\mathcal O_X$. Entonces
$G(\mathfrak X)$ es canónicamente isomorfo al conjunto de las clases de
revestimientos planos del esquema ordinario $X$ que «se reducen a $X'_0$».

De manera gráfica, puede decirse que (a) y (b) establecen las relaciones
fundamentales entre el aspecto local y el aspecto global del problema, (c) da
las relaciones entre el aspecto «finito» y el aspecto «infinitesimal» y, por
último, (d) recuerda (bajo las condiciones precisadas) la identidad entre el
aspecto «formal» y el aspecto «algebraico».

Supongamos ahora que $\mathfrak X$ está definido sobre un anillo local
noetheriano completo $\Lambda$, con
$\mathcal J=\mathfrak m\mathcal O_{\mathfrak X}$ (por tanto, $X_0$ es un
preesquema sobre $\Lambda/\mathfrak m$). Para toda álgebra $A$ finita sobre
$\Lambda$, pongamos
\[
  F(A)=G(\mathfrak X\times_\Lambda A),
\]
donde el marcado se toma después del cambio de base al álgebra reducida $A_{\mathrm{red}}$, y no literalmente a $A/\mathfrak mA$. Es un funtor covariante en $A$, con valores en la categoría de los conjuntos,
y por (c) este funtor queda completamente determinado si se lo conoce para $A$
artiniana; y decir que este funtor es pro-representable, es decir, de la forma
\[
  F(A)=\operatorname{Hom}_{\Lambda\text{-algèbres top.}}(O,A),
\]

\src{389}
donde $O$ es un álgebra topológica sobre $\Lambda$ del tipo considerado en A,
párrafo 5, equivale a decir que lo mismo ocurre al restringirse a los $A$
artinianos. La conjunción de los teoremas 1 y 2 implica entonces efectivamente:

\result{PROPOSICIÓN}{5.1} Siempre que toda familia artiniana clasificada por el funtor sea plana sobre su álgebra de coeficientes y que el abierto escogido permanezca esquemáticamente denso después de los cambios de base pertinentes, el funtor precedente es pro-representable.

Por supuesto, según (a), si $U=\mathfrak X$, $G(\mathfrak Y)$ se reduce a un
solo elemento para todo $\mathfrak Y$ sobre $\mathfrak X$, y el funtor $F$
tiene entonces escaso interés (se tendrá $O=\Lambda$). Parece que, en casi
cualquier otro caso, el anillo local topológico $O$ no es noetheriano. No
obstante, su existencia pone de manifiesto, de manera sorprendente, la
naturaleza «continua» del conjunto $G(\mathfrak X)$ de las soluciones (que
corresponde intuitivamente al hecho de que hay una elección «continua» en el
modo en que la ramificación se propaga al extender $X'_0$). Compárese este
resultado con el punto de vista de J.-P. Serre [4] en la teoría del cuerpo de
clases local, que pone asimismo de relieve el carácter continuo del grupo de
Galois topológico de la extensión abeliana maximal de un cuerpo local
«geométrico»; allí, el grupo dual (en el sentido de Pontryagin) aparece como un
límite directo de grupos algebraicos (o al menos cuasialgebraicos); también
aquí, la clasificación de las extensiones viene dada por «variedades» de
dimensión infinita. Por otra parte, en lo anterior se puede tomar como
$\mathfrak X$ el espectro formal de un anillo local completo (del cual
$\Lambda$ será, por ejemplo, un subanillo de Cohen), y cabe esperar que los
resultados de este apartado puedan utilizarse en el estudio de las extensiones
de un anillo local completo de dimensión $>1$. Tanto en el caso local como en
el global, quizá permitan formular relaciones precisas entre los fenómenos de
ramificación superior y fenómenos en característica cero (abordables por
métodos trascendentes). En cualquier caso, el análisis preliminar de la
proposición 5.1 es el que permite extender al caso «moderadamente ramificado»
los métodos expuestos en [2] para estudiar el grupo fundamental y resolver por
métodos trascendentes el «problema de los tres puntos».

Para terminar, señalemos que la situación se simplifica cuando $X_0$ tiene dimensión 1: entonces, en virtud de (a), complementado por el pegado formal en los esquemas locales completados perforados,
\[
  G(\mathfrak X)\simeq
  \prod_iG\bigl(\operatorname{Spec}(\widehat{\mathcal O}_{\mathfrak X,x_i})\bigr),
\]
donde los $x_i$ son los puntos de $X_0-U$: pueden darse arbitrariamente las
extensiones «locales» en los puntos de ramificación. Además, la normalidad de $X_0$ por sí sola no implica que el esquema formal de módulos de 5.1 sea liso sobre $\operatorname{Spec}(\Lambda)$; la última afirmación impresa es falsa en general.

\src{390}
\begin{center}
  \textbf{BIBLIOGRAFÍA}
\end{center}

\begin{description}
  \item[{[1]}] Grothendieck (Alexander). --- \emph{Sur quelques points
    d'algèbre homologique}, Tohoku Math. J., t. 9, 1957, p. 119--221.
  \item[{[2]}] Grothendieck (Alexander). --- \emph{Géométrie formelle et
    géométrie algébrique}, Séminaire Bourbaki, t. 11, 1958/59, fasc. 3,
    n\textsuperscript{o} 182.
  \item[{[3]}] Grothendieck (Alexander). --- \emph{Technique de descente et
    théorèmes d'existence en géométrie algébrique, I : Généralités, Descente
    par morphismes fidèlement plats}, Séminaire Bourbaki, t. 12, 1959/60,
    fasc. 1, n\textsuperscript{o} 190.
  \item[{[4]}] Serre (Jean-Pierre). --- \emph{Corps locaux et isogénies},
    Séminaire Bourbaki, t. 11, 1958/59, fasc. 3, n\textsuperscript{o} 185.
\end{description}

\def\FgaUnitKey{e195}\clearpage
\chapter*{Fe de erratas de la exposición n.\textsuperscript{o} 195}
\addcontentsline{toc}{chapter}{Fe de erratas de la exposición n.\textsuperscript{o} 195}
\setcounter{footnote}{0}
\src{300}

\noindent\textbf{4. TDTE II: El teorema de existencia en la teoría formal de
los módulos.}

\begin{center}
[Seminario Bourbaki, t. 12, 1959/60, n.\textsuperscript{o} 195, 22 p.]
\end{center}

\noindent\textbf{Página 8.} --- La fórmula escrita en la proposición 5.1 solo es
correcta cuando $\Lambda$ es un cuerpo; en el caso general, hay que sustituir
$\mathfrak m_\xi/\mathfrak m_\xi^2$ por el cociente de este espacio por la
imagen de $\mathfrak n_\xi/\mathfrak n_\xi^2$, donde $\mathfrak n$ es el ideal
maximal de $\Lambda$. Además, la definición dada para que $O$ sea liso sobre
$\Lambda$ solo es correcta cuando la extensión residual $k'/k$ es separable.
En el caso general, cf. SGA III, 1.1.

\noindent\textbf{Página 14, observación.} --- Los problemas planteados aquí
quedan completamente resueltos en el caso proyectivo mediante los «esquemas de
Hilbert» (cf. TDTE IV). Ejemplos de Nagata y Hironaka muestran, en cambio, que
los funtores considerados no son necesariamente representables si no se impone
una hipótesis de proyectividad, incluso limitándose a la clasificación de las
subvariedades de dimensión cero de una variedad completa no singular de
dimensión 3; esto está ligado al hecho de que el cuadrado simétrico de una
variedad tal puede no existir.

\src{301}
\noindent\textbf{Página 15, n.\textsuperscript{o} 3.} --- Para un estudio más
completo, véase TDTE V.

\noindent\textbf{Página 16, proposición 3.1.} --- En lugar de «si $f$ es
propio», léase «si $f$ es propio y separable».

\noindent\textbf{Páginas 18 y 19, observaciones, 2.\textsuperscript{o}} --- He
demostrado recientemente que el esquema formal de módulos de una variedad
abeliana sobre un cuerpo es efectivamente liso sobre el anillo de Witt; en
otros términos, todo esquema abeliano sobre un anillo artiniano local que sea
cociente de otro proviene, por reducción, de un esquema abeliano sobre este
último. La demostración utiliza únicamente las propiedades de varianza de la
clase de obstrucción al levantamiento introducida en la exposición
n.\textsuperscript{o} 182, página 12. Recordemos asimismo que los esquemas de
módulos de las curvas de género $g$ y de los esquemas abelianos polarizados han
sido construidos por Mumford (cf. SHMT).

\noindent\textbf{Página 19, párrafo 5.} --- Hay que suponer que la sección
considerada para definir $U$ sea un no divisor de cero no solo sobre
$\mathfrak X$, sino sobre todo $X_n$.

\noindent\textbf{Página 20, líneas 12 y 13.} --- En lugar de
\[
  G_x=G\bigl(\operatorname{Spec}(\widehat{\mathcal O}_{\mathfrak X,x})\bigr),
\]
léase
\[
  G_x=G\bigl(\operatorname{Spec}(\mathcal O_{\mathfrak X,x})\bigr)
  \subset
  G\bigl(\operatorname{Spec}(\widehat{\mathcal O}_{\mathfrak X,x})\bigr),
\]
y, en lugar de «$\mathcal O'_{0x}$», léase
«$\widehat{\mathcal O}'_{0x}$».

\end{document}
